\documentclass[10pt,a4paper]{article}

\usepackage[T1]{fontenc}
\usepackage[utf8]{inputenc}
\usepackage[french]{babel}
\usepackage{lmodern}
\usepackage{microtype}
\usepackage[a4paper,margin=2.2cm,headheight=20pt]{geometry}
\usepackage{enumitem}
\usepackage{needspace}
\usepackage{xcolor}
\usepackage{fancyhdr}
\usepackage[autostyle]{csquotes}
\usepackage[hidelinks]{hyperref}
\usepackage[backend=biber,style=authoryear,maxcitenames=2,maxbibnames=12]{biblatex}
\addbibresource{bibliographie/references.bib}
\hypersetup{
  pdftitle={L'émergence de l'intéressant - Catalogue provisoire des idées},
  pdfauthor={François Pachet}
}

\definecolor{metadata}{HTML}{59636B}
\definecolor{rulecolor}{HTML}{C8CDD1}
\setlength{\parindent}{0pt}
\setlength{\parskip}{0.55em}
\setlength{\emergencystretch}{2em}
\setcounter{tocdepth}{1}
\setlist[itemize]{leftmargin=1.5em,itemsep=0.25em}
\renewcommand{\familydefault}{\sfdefault}
\pagestyle{fancy}
\fancyhf{}
\fancyhead[L]{\footnotesize Catalogue des idées}
\fancyhead[R]{\footnotesize\color{metadata} État de travail}
\fancyfoot[C]{\thepage}
\renewcommand{\headrulewidth}{0.3pt}
\renewcommand{\headrule}{\hbox to\headwidth{\color{rulecolor}\leaders\hrule height \headrulewidth\hfill}}

\title{L'émergence de l'intéressant\\\large Catalogue provisoire des idées}
\author{François Pachet}
\date{État de travail du 22 juillet 2026}

\begin{document}
\maketitle

\begin{abstract}
Ce document rassemble les 123 propositions actuellement extraites du corpus du projet de thèse. Il suit l'ordre des sept familles argumentatives de l'index de travail. Cette organisation est provisoire : elle sert à partager, comparer et discuter les propositions avant de fixer un plan de thèse. Les formulations et les attributions documentaires pourront encore être révisées.
\end{abstract}

\tableofcontents
\clearpage
\section{Définition et émergence de l'intéressant}

\Needspace{8\baselineskip}
\subsection{Naissance d'une idée, compréhension et fascination sont trois faces de l'émergence d'une forme}
{\small\color{metadata} \texttt{idea\_0083}\quad\textbullet\quad hypothèse\quad\textbullet\quad conceptuelle\par}

L'idée qui se cristallise, le problème qui devient soudain intelligible et l'objet qui appelle la reprise semblent relever de domaines différents. L'hypothèse du projet est qu'ils manifestent un même événement : quelque chose se détache pour un esprit comme une organisation saillante et structurée.

La naissance, la compréhension et la fascination ne seraient donc pas trois problèmes juxtaposés. Elles décrivent respectivement l'apparition, la saisie et la persistance attentionnelle d'une forme devenue intelligible.

\begin{quote}
\textbf{Intérêt pour la thèse.} Cette proposition donne son unité au projet doctoral. Elle permet de faire de l'intéressant une catégorie philosophique plutôt qu'un qualificatif psychologique ou mondain.
\end{quote}
{\footnotesize\color{metadata}\textbf{Provenance :} projet thèse philo.pdf. PDF p. 1\par}

\Needspace{8\baselineskip}
\subsection{L'intéressant est la rencontre entre une forme organisée et une mémoire capable d'en éprouver la singularité}
{\small\color{metadata} \texttt{idea\_0084}\quad\textbullet\quad définition\quad\textbullet\quad conceptuelle\par}

L'intéressant ne réside ni entièrement dans l'objet ni entièrement dans une préférence privée. Il émerge lorsqu'une organisation de tensions et de résolutions rencontre un esprit historiquement, affectivement et culturellement constitué. Sans structure dans l'objet, rien ne se détache ; sans mémoire, attentes et aptitudes perceptives, cette structure ne peut être éprouvée comme singulière.

La relativité à l'auditeur n'implique pas que tout jugement se vaille. Une propriété relationnelle peut être robuste lorsqu'on peut décrire le problème que la forme résout, les attentes qu'elle déjoue et les compétences nécessaires pour entendre cette résolution. Une analyse réussie explicite alors la relation au lieu de remplacer l'expérience par une liste de traits.

Cette distribution des causes entre objet et sujet possède un précédent spinoziste. Dans la démonstration de la proposition XXVII de la partie III de l'\emph{Éthique}, les idées des affections enveloppent à la fois la nature du corps propre et celle du corps extérieur. Sans assimiler l'intéressant à une passion spinoziste déterminée, ce schème interdit d'en chercher la cause complète dans un seul terme : la forme contribue par son organisation et ses prises ; le sujet par sa constitution, son histoire et ses capacités.

Cette rencontre conjugue trois conditions : une intelligibilité suffisante pour ne pas rester confuse, une singularité qui résiste au déjà-vu, et une fécondité qui ouvre des reprises, variations ou pensées nouvelles.

La relation minimale est un binôme entre une forme et un sujet, mais ce binôme est indexé par un horizon collectif et par un moment : \texttt{I(F, S | H, t)}. Le milieu social n'est pas nécessairement un troisième sommet indépendant, car il façonne déjà les attentes du sujet, les catégories disponibles et les formes qui peuvent lui parvenir. Il devient un troisième terme explicite lorsqu'on étudie non plus l'expérience d'un sujet, mais la circulation, la stabilisation ou la disparition collective d'un intérêt. La thèse doit donc distinguer la structure relationnelle minimale de son explication sociale complète.

La mémoire pertinente n'est pas seulement le souvenir explicite d'une œuvre. Des écoutes passives, des erreurs de disque, des gestes instrumentaux et des contextes oubliés préparent des attentes qui se réveillent beaucoup plus tard. L'histoire de l'oreille est en partie irreconstructible : apprendre à entendre ouvre de nouveaux détails tout en rendant l'écoute antérieure inaccessible.

\begin{quote}
\textbf{Intérêt pour la thèse.} Cette définition relationnelle évite à la fois l'objectivisme des traits universels et le subjectivisme qui rendrait toute explication impossible.
\end{quote}
{\footnotesize\color{metadata}\textbf{Provenance :} projet thèse philo.pdf; The Mystery of Jotney Songs -full.pdf; PACHET HISTOIRE OREILLE BAT.pdf. PDF p. 1-2 Dossier Jotney, propriété relationnelle robuste et effet de l'analyse, PDF p. 3-4 Histoire d'une oreille, apprentissages passifs qui préparent des effets futurs, PDF p. 27-34 ; changement de goût et dépendance au contexte, p. 281-295 ; impossibilité de retrouver l'oreille antérieure, p. 306-310. Spinoza, Éthique, partie III, proposition XXVII, démonstration : les idées des affections enveloppent la nature de notre corps et celle du corps extérieur.\par}
{\footnotesize\color{metadata}\textbf{Références :} \parencite{pachet2018oreille,spinoza1861oeuvres}\par}

\Needspace{8\baselineskip}
\subsection{L'intéressant peut être défini comme la zone de flow entre ennui et anxiété}
{\small\color{metadata} \texttt{idea\_0121}\quad\textbullet\quad hypothèse\quad\textbullet\quad conceptuelle\par}

Pour un sujet donné à un moment donné, l'intéressant peut être défini comme la zone où la difficulté d'une tâche sollicite ses compétences sans être ni trop faible ni trop forte. À compétences fixes, une difficulté croissante produit la succession \textbf{ennui → intéressant/flow → anxiété}. L'intéressant n'est donc pas seulement ce qui s'oppose à l'ennuyeux : il occupe une région intermédiaire entre sous-sollicitation et perte de prise.

Cette définition est relationnelle. Une même tâche peut être ennuyeuse pour un expert, intéressante pour un sujet suffisamment préparé et anxiogène ou opaque pour un novice. Elle est aussi dynamique : pour une tâche fixe, l'apprentissage fait parcourir le diagramme en sens inverse, \textbf{anxiété → intéressant/flow → ennui}. L'intérêt apparaît, se maintient et s'épuise avec le déplacement du rapport entre compétences et exigences.

Il faut distinguer la \textbf{zone de flow}, qui fournit cette structure, de l'\textbf{état de flow}, expérience d'absorption intense décrite par Csíkszentmihályi. L'hypothèse identifie l'intéressant à la première ; elle ne suppose pas que tout objet intéressant produit continuellement le second. Surprise, hésitation et rupture peuvent appartenir à la zone intéressante précisément parce que le sujet conserve une prise sur une difficulté qui le sollicite.

La zone n'est cependant pas homogène. Les travaux sur la curiosité artificielle suggèrent que l'attention se porte plus précisément vers les situations où le progrès d'apprentissage estimé est maximal. L'ajustement compétences-difficulté donne la géographie générale ; le gradient de progrès de \texttt{idea\_0071} indique où, dans cette zone, une construction est effectivement en train d'aboutir.

\begin{quote}
\textbf{Intérêt pour la thèse.} Cette hypothèse donne à l'intéressant une définition positive, relationnelle et opérationnalisable. Elle unifie ennui, apprentissage, difficulté et attention dans un même espace de variation, tout en fournissant des conditions d'échec observables.
\end{quote}
{\footnotesize\color{metadata}\textbf{Provenance :} De l'impossibilité de créer.pdf. PDF p. 59-62 pour le modèle du flow ; l'identification à l'intéressant est une hypothèse propre à la thèse\par}

\Needspace{8\baselineskip}
\subsection{Pour les objets temporels, l'intérêt appartient à l'agencement plus qu'à l'item isolé}
{\small\color{metadata} \texttt{idea\_0001}\quad\textbullet\quad définition\quad\textbullet\quad conceptuelle\par}

Un objet n'est pas intéressant par la seule possession de certaines propriétés. Son intérêt dépend de la séquence dans laquelle il apparaît : ce qui le précède fait naître des attentes, son apparition les confirme ou les déplace, puis la répétition transforme l'effet en reconnaissance ou en fatigue. Le même objet peut donc captiver ou ennuyer selon son ordre d'apparition.

Cette proposition déplace l'unité d'analyse : il faut modéliser des trajectoires d'attention et des agencements, non classer une collection d'objets une fois pour toutes.

Sa portée doit cependant être délimitée. Certains dispositifs semblent surtout attirants par les items qu'ils montrent et supportent une inversion de l'ordre sans perte majeure; d'autres utilisent des matériaux pauvres dont toute la force vient de l'organisation. Inverser ou permuter la séquence fournit ainsi un test empirique du poids propre de la structure temporelle.

\begin{quote}
\textbf{Intérêt pour la thèse.} Cette carte peut servir de définition de départ : l'intéressant n'est pas seulement ce qui attire, mais ce qui organise une durée d'attention.
\end{quote}
{\footnotesize\color{metadata}\textbf{Provenance :} Synopsis MIT Press.doc; projet thèse philo.pdf; TBKLullyNOTES.doc. Projet thèse philo, PDF p. 1-3 Synopsis MIT Press, rendu PDF p. 7-10 et 18-20 TBKLullyNOTES, distinction temporel/atemporal, rendu PDF p. 5\par}

\Needspace{8\baselineskip}
\subsection{Le goût n'est pas un dictionnaire statique}
{\small\color{metadata} \texttt{idea\_0002}\quad\textbullet\quad argument\quad\textbullet\quad conceptuelle\par}

Un dictionnaire de préférences suppose que le sujet possède déjà des jugements stables qu'il suffirait d'appliquer aux objets. Or l'interaction avec les objets modifie ce que le sujet sait reconnaître, attendre et désirer. Une technologie de création ou de recommandation ne doit donc pas seulement retrouver un goût préexistant : elle participe à sa formation.

\begin{quote}
\textbf{Intérêt pour la thèse.} Cette idée attaque directement les systèmes de recommandation et les théories plates de la préférence. Elle ouvre la voie à une théorie de l'intéressant comme relation évolutive.
\end{quote}
{\footnotesize\color{metadata}\textbf{Provenance :} ERCGrantPachetInterestingness.pdf; ERCInteractiveReflexions2.docx; TBKLullyNOTES.doc. Grant ERC PDF p. 1 Projet REFLEX, rendu PDF p. 2-4 TBKLullyNOTES, paradoxe aimer/aimer faire, rendu PDF p. 2\par}

\Needspace{8\baselineskip}
\subsection{L'intéressant mesure un progrès de compression, pas une surprise}
{\small\color{metadata} \texttt{idea\_0027}\quad\textbullet\quad argument\quad\textbullet\quad articulation\par}

L'intérêt ne vient ni de la complexité d'un objet, ni de sa nouveauté, ni de sa surprise prises séparément. Il vient du progrès que fait l'observateur lorsqu'il découvre une régularité qui rend les données plus simples à représenter ou à prédire. Autrement dit, ce qui est récompensé n'est pas la compression déjà acquise, mais l'amélioration du compresseur ou du modèle.

Cette proposition explique dans un même mouvement pourquoi le parfaitement prévisible et le bruit blanc sont ennuyeux. Le premier est déjà compressé et ne produit plus de progrès ; le second reste incompressible et ne permet aucun progrès durable. Une surprise n'est intéressante que si elle annonce une régularité apprenable.

Dans le papier de 1997, Schmidhuber formule ce mécanisme en termes de régularités inattendues faciles à apprendre et de récompense pour croissance des connaissances. Il ne parle pas encore explicitement de \texttt{compression progress}. Le projet ERC reformule ensuite l'intérêt comme dérivée d'une beauté comprise en termes de compressibilité. La carte distingue donc la proposition générale du papier source et cette reformulation, au lieu d'attribuer anachroniquement l'expression au rapport de 1997.

Le mécanisme de confiance est essentiel. Une erreur de prédiction ne suffit pas : le système ne récompense une surprise que lorsqu'une partie s'était engagée sur un résultat et que l'autre pouvait raisonnablement anticiper son erreur. Cette condition empêche le bruit blanc de devenir une réserve illimitée de récompenses.

\begin{quote}
\textbf{Intérêt pour la thèse.} Cette proposition donne un critère opérationnel plus fort que la simple "zone entre trivialité et bruit" : on peut rechercher une variation de performance prédictive ou descriptive chez un observateur donné. Elle permet aussi de distinguer l'intéressant de l'étonnant.
\end{quote}
{\footnotesize\color{metadata}\textbf{Provenance :} interestingness.pdf; ERCGrantPachetInterestingness.pdf. Schmidhuber, What's Interesting?, PDF p. 1-4 et 18-21 ERCGrantPachetInterestingness, beauté compressible et intérêt comme dérivée, PDF p. 2\par}
{\footnotesize\color{metadata}\textbf{Références :} \parencite{schmidhuber1997interesting}\par}

\Needspace{8\baselineskip}
\subsection{L'intérêt suit le gradient de progrès d'un observateur, non son erreur actuelle}
{\small\color{metadata} \texttt{idea\_0071}\quad\textbullet\quad hypothèse\quad\textbullet\quad articulation\par}

Si l'intérêt correspond à un progrès de compression, il ne peut pas être une propriété stable de l'objet. Il dépend de ce que l'observateur sait déjà, de ce qu'il est capable de calculer et du moment de son apprentissage. La régularité qui provoque aujourd'hui un gain de compréhension devient demain une connaissance acquise et donc une source d'ennui.

Cette proposition est plus précise que l'affirmation générale selon laquelle les goûts sont subjectifs : elle indique quelle transformation interne fait varier l'intérêt. Apprendre détruit progressivement la récompense qui a motivé l'apprentissage et force la curiosité à changer d'objet.

Il n'existe donc pas de proportion universelle entre information attendue et inattendue. La frontière entre trivial et aléatoire se déplace avec les savoirs et les capacités de calcul de l'observateur.

Les travaux de Kaplan, Oudeyer et Hafner permettent de préciser cette transformation comme un gradient. Si \texttt{E\_t(o, S)} désigne l'erreur de prédiction du sujet \texttt{S} confronté à l'objet ou à la situation \texttt{o}, un signal simple de progrès est \texttt{P\_t(o, S) = E\_\{t-\ensuremath{\Delta}\}(o, S) - E\_t(o, S)}. L'intérêt ne suit donc pas le niveau actuel de l'erreur, mais sa diminution. Une situation maîtrisée produit peu d'erreur mais aucun progrès ; une situation aléatoire ou trop difficile produit beaucoup d'erreur mais aucun progrès ; une situation intéressante est celle où la compétence est en train d'augmenter.

Cette mesure définit des \textbf{niches de progrès}, locales et temporaires dans l'espace des situations accessibles. Un agent curieux choisit les régions où son modèle s'améliore le plus vite, puis les quitte lorsque le progrès s'aplatit. Il construit ainsi sans curriculum prédéfini une trajectoire allant vers des activités de difficulté croissante. Le gradient n'est pas une propriété de l'objet seul : il naît de la rencontre entre ce que l'objet permet de percevoir ou d'essayer et les compétences, biais, moyens d'action et expériences déjà acquis par le sujet.

Schmidhuber formule un principe voisin en termes de progrès de compression : la récompense de curiosité vient de l'amélioration du compresseur ou du prédicteur, non de la surprise brute. Les deux familles de travaux convergent donc sur une dérivée temporelle de l'apprentissage, même si leurs architectures et leurs mesures ne sont pas identiques.

La neuveclamine en donne une micro-histoire concrète : découverte fascinante, reproduction intensive, saturation, puis réapprentissage d'un usage rare et presque indétectable. L'effet ne disparaît pas simplement ; sa valeur dépend de la manière dont l'attente est de nouveau préparée après son assimilation.

\begin{quote}
\textbf{Intérêt pour la thèse.} Elle permet de définir l'intéressant comme une relation dynamique entre une forme et l'état d'un observateur. Elle donne aussi une mesure candidate du travail constructif : une construction est en cours lorsque les erreurs diminuent ou que les prises s'améliorent. Elle prédit enfin qu'une œuvre peut rester intéressante si elle offre plusieurs régularités découvrables à des échelles ou des moments différents, plutôt qu'une seule surprise rapidement épuisée.
\end{quote}
{\footnotesize\color{metadata}\textbf{Provenance :} interestingness.pdf; The Mystery of Jotney Songs.pdf; PACHET HISTOIRE OREILLE BAT.pdf. Schmidhuber, What's Interesting?, PDF p. 1-2 et 18-20 Kaplan, Oudeyer et Hafner, progrès d'apprentissage, niches de progrès et développement autonome, IEEE 2007, p. 265-286 Jotney, variation de la surprise avec la formation de l'oreille, PDF p. 3 Histoire d'une oreille, fascination, saturation puis usage parcimonieux de la neuveclamine, PDF p. 23-24 ; solutions devenues évidentes après réécoute, p. 285.\par}
{\footnotesize\color{metadata}\textbf{Références :} \parencite{schmidhuber1997interesting,schmidhuber2009compression,oudeyer2007intrinsic,oudeyer2007typology,pachet2018oreille}\par}

\Needspace{8\baselineskip}
\subsection{L'intéressant se situe à la frontière mouvante de ce qui est presque apprenable}
{\small\color{metadata} \texttt{idea\_0107}\quad\textbullet\quad hypothèse\quad\textbullet\quad conceptuelle\par}

Un objet n'est pas intéressant parce qu'il contient une proportion fixe de connu et d'inconnu. Il l'est lorsqu'il présente, pour un observateur donné, une régularité encore inattendue mais accessible avec ses connaissances et ses moyens de calcul actuels. Le déjà appris est trivial, le bruit est inépuisablement imprévisible, et une structure trop éloignée des capacités présentes reste opaque : ces trois cas échouent pour des raisons différentes.

Cette frontière n'est donc ni une propriété de l'objet ni un ratio informationnel universel. Elle dépend du temps et de l'effort requis pour apprendre. Schmidhuber se distingue ici d'une simplicité fondée seulement sur la longueur du plus court programme : une régularité courte en description peut rester pratiquement inaccessible si son calcul est trop coûteux.

\begin{quote}
\textbf{Intérêt pour la thèse.} La proposition transforme la vague idée d'un juste milieu en hypothèse testable : l'intérêt devrait culminer pour les formes qui provoquent un gain atteignable au prix d'un effort non nul. Elle rapproche l'intéressant de la zone de difficulté crédible de la virtuosité sans les confondre. Elle précise aussi le modèle du flow : dans la succession \textbf{ennui → intéressant/flow → anxiété}, le presque-apprenable situe une part de l'intérêt vers la frontière supérieure de la zone intermédiaire, au moment où une difficulté devient accessible et promet un gain de compétence ou de compréhension.
\end{quote}
{\footnotesize\color{metadata}\textbf{Provenance :} interestingness.pdf; The Mystery of Jotney Songs.pdf. Schmidhuber, What's Interesting?, PDF p. 1-3 et 20 Jotney, bande étroite entre mobilité, densité, monotonie et surprise, PDF p. 3\par}
{\footnotesize\color{metadata}\textbf{Références :} \parencite{schmidhuber1997interesting}\par}

\Needspace{8\baselineskip}
\subsection{Syntaxe et sens ne suffisent pas à expliquer ce qui retient l'attention}
{\small\color{metadata} \texttt{idea\_0003}\quad\textbullet\quad distinction\quad\textbullet\quad conceptuelle\par}

Deux messages syntaxiquement bien formés et sémantiquement intelligibles peuvent avoir des pouvoirs d'attention radicalement différents. Une théorie de la communication qui s'arrête à la syntaxe et au sens laisse donc inexpliquée une troisième dimension : la capacité d'une forme à attirer puis soutenir l'activité d'un destinataire. Cette dimension exige ses propres observations et modèles ; elle ne se déduit pas automatiquement des deux premières.

\begin{quote}
\textbf{Intérêt pour la thèse.} Cette distinction donne une architecture conceptuelle forte : l'intéressant n'est ni réductible à la forme, ni au sens, ni à l'utilité.
\end{quote}
{\footnotesize\color{metadata}\textbf{Provenance :} ERCGrantPachetInterestingness.pdf; ERCInteractiveReflexions2.docx. Grant ERC PDF p. 2 Projet REFLEX, notes d'hypothèses, rendu PDF p. 5\par}

\Needspace{8\baselineskip}
\subsection{L'inattendu ne suffit pas}
{\small\color{metadata} \texttt{idea\_0005}\quad\textbullet\quad objection\quad\textbullet\quad conceptuelle\par}

Une chose inattendue n'est pas automatiquement intéressante. L'inattendu peut produire une surprise locale, mais aussi une fatigue ou un ennui infini s'il n'est pas inscrit dans une structure interprétable. Les interfaces qui multiplient les anomalies et les scénarios qui ne font que violer les attentes illustrent cette limite : l'événement doit être non seulement inattendu, mais pertinent pour ce qui précède et susceptible de recevoir une interprétation.

Le bruit fournit ainsi beaucoup d'information au sens de Shannon et peut faire échouer indéfiniment un prédicteur sans devenir intéressant. Chez Schmidhuber, une surprise digne de récompense suppose au contraire une attente confiante et la possibilité qu'une régularité soit apprise.

\begin{quote}
\textbf{Intérêt pour la thèse.} Cette carte évite une définition trop courte de l'intéressant par la simple surprise.
\end{quote}
{\footnotesize\color{metadata}\textbf{Provenance :} Synopsis MIT Press.doc; interestingness.pdf. Synopsis MIT Press, rendu PDF p. 4 et 28 Schmidhuber, What's Interesting?, PDF p. 2 et 4\par}
{\footnotesize\color{metadata}\textbf{Références :} \parencite{schmidhuber1997interesting}\par}

\Needspace{8\baselineskip}
\subsection{Un concept hypertrophié peut cacher les mécanismes qu'il prétend expliquer}
{\small\color{metadata} \texttt{idea\_0079}\quad\textbullet\quad objection\quad\textbullet\quad conceptuelle\par}

Lorsqu'un mot comme \texttt{créativité}, \texttt{innovation}, \texttt{talent} ou \texttt{génie} est appliqué à des phénomènes trop différents, il ne les rend plus lisibles. Il déclenche des schémas de reconnaissance globaux qui recouvrent les contraintes locales, les processus lents, l'appropriation et les transformations subjectives.

Cette ombre ontologique vient d'un excès de concept, non de son absence. Comme le dessinateur voit mieux les lignes d'un visage lorsqu'il cesse momentanément de le reconnaître comme visage, l'enquête sur la création doit parfois suspendre son vocabulaire central pour observer ce qu'il empêche de distinguer.

\begin{quote}
\textbf{Intérêt pour la thèse.} Cette proposition offre une discipline méthodologique : ne pas laisser le mot \texttt{intéressant} devenir à son tour une explication globale qui dispense d'observer.
\end{quote}
{\footnotesize\color{metadata}\textbf{Provenance :} De l'impossibilité de créer.pdf. PDF p. 71-73\par}

\Needspace{8\baselineskip}
\subsection{Un modèle trop général de l'intéressant devient difficile à falsifier et inutile pour produire}
{\small\color{metadata} \texttt{idea\_0095}\quad\textbullet\quad objection\quad\textbullet\quad articulation\par}

Les courbes en U, la surprise bayésienne ou le progrès de compression expliquent pourquoi le trivial et le bruit sont peu attirants dans de nombreux domaines. Mais leur indépendance au contenu leur permet aussi de s'accorder avec presque tout résultat après coup. Elles indiquent une forme générale sans préciser quelles régularités un sujet peut percevoir ni comment produire un objet particulier.

Une théorie computationnelle utile doit donc descendre vers des hypothèses risquées et situées - rareté dans un espace défini, difficulté de production, réflexivité, transformation du goût - puis construire des expériences capables de les séparer.

Le papier de Schmidhuber rend cette limite visible : son langage universel autorise en principe presque toute régularité calculable, mais cette généralité rend les dynamiques effectives difficiles à comprendre. Il propose lui-même de réduire la puissance des systèmes futurs pour rendre leurs comportements plus analysables.

\begin{quote}
\textbf{Intérêt pour la thèse.} Cette objection permet d'utiliser Schmidhuber et la théorie de l'information comme fondations sans les prendre pour une explication complète de l'intéressant.
\end{quote}
{\footnotesize\color{metadata}\textbf{Provenance :} ERCGrantPachetInterestingness.pdf; interestingness.pdf. PDF p. 2-4 Schmidhuber, What's Interesting?, généralité et difficultés d'analyse, PDF p. 12 et 20-21\par}
{\footnotesize\color{metadata}\textbf{Références :} \parencite{schmidhuber1997interesting}\par}

\Needspace{8\baselineskip}
\subsection{Une forme peut être belle parce qu'elle est simple et ennuyeuse parce qu'elle est déjà connue}
{\small\color{metadata} \texttt{idea\_0108}\quad\textbullet\quad distinction\quad\textbullet\quad conceptuelle\par}

La beauté et l'intérêt ne sélectionnent pas nécessairement les mêmes objets. Une forme peut être belle parce qu'elle possède une description simple relativement au codage de l'observateur, tout en étant ennuyeuse si cette simplicité est déjà connue. L'intérêt exige en plus que la régularité simple soit encore inattendue ou nouvellement saisissable.

Inversement, un objet peut être intéressant sans constituer le meilleur exemple de beauté : ce qui compte alors est la transformation du modèle de l'observateur, pas seulement l'économie de la forme. Identifier beau, bon et intéressant effacerait donc le rôle propre de la découverte.

\begin{quote}
\textbf{Intérêt pour la thèse.} Cette distinction interdit de traiter l'intéressant comme un simple synonyme de valeur esthétique. Elle permet de demander séparément ce qui rend une forme économique, ce qui la rend nouvelle pour quelqu'un et ce qui maintient son pouvoir d'attention.
\end{quote}
{\footnotesize\color{metadata}\textbf{Provenance :} interestingness.pdf. Schmidhuber, What's Interesting?, distinction entre beauty et interestingness, PDF p. 20-21\par}
{\footnotesize\color{metadata}\textbf{Références :} \parencite{schmidhuber1997interesting}\par}

\Needspace{8\baselineskip}
\subsection{L'intéressant occupe un angle mort entre intelligibilité et valeur}
{\small\color{metadata} \texttt{idea\_0116}\quad\textbullet\quad hypothèse\quad\textbullet\quad conceptuelle\par}

Les disciplines de la connaissance disposent de concepts puissants pour distinguer le bien formé du mal formé, le sens du non-sens, le vrai du faux, le pertinent de l'inopportun et le beau du laid. Elles décrivent beaucoup moins directement ce qui, parmi les formes admissibles, intelligibles ou même vraies, mérite et soutient l'attention d'un sujet. L'intéressant pourrait ainsi désigner un angle mort : une valeur relationnelle assez forte pour orienter l'enquête, la lecture ou la création, mais moins terminale et moins substantielle que le vrai, le beau ou le bien.

Cette proposition ne suppose pas une progression historique simple allant de la syntaxe à la sémantique, puis à l'intérêt et enfin à la beauté. Elle avance une hypothèse de travail plus limitée : la correction et l'intelligibilité définissent des conditions d'admissibilité, tandis que l'intéressant sélectionne, à l'intérieur de cet espace, les formes capables de transformer l'attention ou la compréhension. En ce sens, il constitue peut-être un objet plus opératoire que la beauté pour étudier la dynamique de la découverte et de la création.

\begin{quote}
\textbf{Intérêt pour la thèse.} La thèse ne viserait donc pas seulement à définir un mot négligé. Elle devrait montrer pourquoi plusieurs traditions philosophiques et scientifiques rencontrent les effets de l'intéressant sous d'autres noms sans en faire un objet autonome, puis dégager les propriétés que ce déplacement rend visibles : temporalité, relation à une mémoire, sélection parmi des formes correctes, transformation de l'attention et épuisement possible par familiarisation.

Le rôle de Chomsky et du structuralisme serait celui d'un jalon dans cette histoire. Leur limitation pertinente n'est pas qu'ils confondraient correction et vérité : la grammaticalité chomskyenne ne dit précisément rien de la vérité d'une phrase. C'est plutôt que leur découpage méthodologique privilégie les structures, les règles de formation et les conditions d'interprétation, en laissant hors modèle la différence de fécondité ou de puissance attentionnelle entre deux formes également admissibles.

Les travaux neuroesthétiques d'Ishizu et Zeki cherchent, quant à eux, un corrélat neuronal commun à des expériences déclarées de beauté. Cette recherche d'un invariant peut masquer la relation historique qui intéresse la thèse, mais Zeki reconnaît aussi que la beauté des artefacts varie selon les individus, les cultures et l'exposition. La critique ne doit donc pas lui attribuer une négation absolue du sujet ou du social : elle doit montrer que l'invariant neuronal ne décrit pas la transformation temporelle par laquelle une forme devient intéressante pour un sujet situé. La beauté n'est pas à abolir ; elle cesse d'être la catégorie directrice de cette enquête.
\end{quote}
{\footnotesize\color{metadata}\textbf{Provenance :} projet-these-v1-fr.tex. Section finale sur l'originalité : l'intéressant comme catégorie intermédiaire entre vrai, beau, bien, sens et jugement\par}
{\footnotesize\color{metadata}\textbf{Références :} \parencite{chomsky1957syntactic,ishizu2011brainbeauty,zeki2018mathematicalbeauty}\par}

\Needspace{8\baselineskip}
\subsection{Préférer un objet et le trouver intéressant sont deux relations distinctes}
{\small\color{metadata} \texttt{idea\_0118}\quad\textbullet\quad distinction\quad\textbullet\quad conceptuelle\par}

« J'aime ce morceau » n'implique pas « je trouve ce morceau intéressant ». Le premier énoncé exprime une relation affective ou préférentielle ; le second attribue au morceau une capacité à provoquer ou soutenir une activité d'attention, de compréhension, d'exploration ou de reprise. La réciproque est également fausse : je peux trouver un morceau intéressant sans l'aimer.

Préférer un objet peut consister à l'aimer davantage, à le choisir ou à anticiper qu'il procurera plus de satisfaction qu'une alternative. Ces relations peuvent coïncider avec l'intérêt, mais aucune ne l'implique.

On peut préférer une forme familière, confortable ou identitaire tout en la trouvant prévisible et sans intérêt. Inversement, une œuvre déplaisante, une idée inquiétante, un échec ou un comportement que l'on désapprouve peuvent être intéressants par l'enquête, la résistance ou la transformation qu'ils provoquent. Les quatre combinaisons sont possibles : préféré et intéressant, préféré mais inintéressant, non préféré mais intéressant, ni préféré ni intéressant.

L'absence d'implication ne signifie pas indépendance causale. Une exposition rendue possible par la préférence peut faire naître un intérêt tardif ; un intérêt répété peut former un attachement ou une préférence ; la familiarité peut renforcer le goût tout en épuisant l'intérêt. La recherche doit étudier ces passages au lieu d'utiliser la préférence comme mesure indirecte automatique de l'intéressant.

\begin{quote}
\textbf{Intérêt pour la thèse.} Cette distinction empêche de réduire l'intéressant à une psychologie du goût ou à un score de recommandation. Elle justifie de mesurer séparément choix, plaisir déclaré, durée d'attention, reprise, transformation de l'objet et évolution ultérieure du sujet.

Elle renforce également la neutralité axiologique de l'enquête : un objet peut devenir cognitivement ou attentionnellement puissant sans être aimé, approuvé ou recherché comme une fin.
\end{quote}
{\footnotesize\color{metadata}\textbf{Provenance :} TBKLullyNOTES.doc; Synopsis MIT Press.doc; projet-these-fr.tex. TBKLullyNOTES, distinction entre aimer un objet et aimer l'activité qu'il provoque, rendu PDF p. 2 Synopsis MIT Press, préférence déclarée et comportement effectif, rendu PDF p. 10 Question directrice : l'intéressant ne se réduit pas à une préférence subjective donnée sans histoire\par}

\Needspace{8\baselineskip}
\subsection{L'intérêt du tout ne se déduit pas de l'intérêt de ses parties}
{\small\color{metadata} \texttt{idea\_0119}\quad\textbullet\quad hypothèse\quad\textbullet\quad conceptuelle\par}

Une axiomatique de l'intéressant doit commencer par préciser sa signature. Soient un sujet \texttt{p}, des objets ou séquences \texttt{o}, un horizon \texttt{H}, un moment \texttt{t}, une évaluation \texttt{I(p, o | H, t)} et une opération \texttt{o1 + o2} qui combine deux objets. \texttt{I} peut d'abord être booléen, puis devenir une grandeur ou une trajectoire si les cas l'exigent.

La première propriété candidate est une non-compositionnalité :

\texttt{I(p, o1) et I(p, o2)} n'impliquent pas \texttt{I(p, o1 + o2)}.

Deux morceaux, idées ou mécanismes intéressants peuvent former un ensemble inintéressant, confus ou saturé. La réciproque échoue aussi : \texttt{I(p, o1 + o2)} n'implique pas que \texttt{I(p, o1)} ou \texttt{I(p, o2)}. Une relation, un contraste ou un ordre peut rendre le tout intéressant alors qu'aucune partie ne l'est isolément. Il n'existe donc pas nécessairement de fonction indépendante du contexte qui calcule l'intérêt du tout à partir du seul intérêt de ses parties.

Cette proposition ouvre un programme de propriétés à tester plutôt qu'une liste d'axiomes déjà admis :

- non-commutativité : changer l'ordre des parties peut changer l'intérêt ; - non-idempotence : répéter un objet peut l'épuiser, le renforcer ou le révéler ; - non-monotonie : ajouter une partie peut augmenter ou détruire l'intérêt ; - dépendance historique : le même couple sujet-objet change avec l'exposition ; - non-universalité : l'intérêt pour un sujet ne se transmet pas automatiquement à un autre.

Ces « non-lois » ne sont pas des défauts de la théorie. Elles peuvent caractériser un objet dont la valeur appartient à des agencements et à des transformations plutôt qu'à une somme de propriétés locales.

\begin{quote}
\textbf{Intérêt pour la thèse.} Cette carte donne une forme rigoureuse à plusieurs intuitions du corpus : primauté de l'agencement temporel, émergence de relations globales, variation avec la mémoire du sujet et insuffisance des traits locaux. Elle permettrait de comparer musique, jeux, textes et systèmes génératifs dans un même langage formel.

Le terme « axiomatique » doit rester provisoire. Chaque propriété doit être accompagnée de cas positifs, de contre-exemples et de conditions de validité. Si toutes les exceptions sont absorbées dans le contexte \texttt{H}, le formalisme deviendrait irréfutable ; il faut donc spécifier quelles variables contextuelles sont observables et quelles prédictions la propriété exclut.
\end{quote}
{\footnotesize\color{metadata}\textbf{Provenance :} BUT DE LA THESE.md; projet thèse philo.pdf. Programme de formalisation : relation entre forme, sujet, horizon collectif et temps Projet de thèse V1, temporalité de l'agencement et émergence relationnelle, PDF p. 1-2\par}

\Needspace{8\baselineskip}
\subsection{L'intéressant est une modalité dynamique qui reconfigure le possible}
{\small\color{metadata} \texttt{idea\_0120}\quad\textbullet\quad hypothèse\quad\textbullet\quad conceptuelle\par}

On peut noter \texttt{I[p,H,t](phi)} le fait qu'une proposition, une forme ou une possibilité \texttt{phi} soit intéressante pour un sujet \texttt{p}, dans un horizon \texttt{H} et au moment \texttt{t}. Cet opérateur ne se comporte probablement pas comme les modalités normales \texttt{nécessaire}, \texttt{possible}, \texttt{obligatoire} ou \texttt{permis}.

Plusieurs règles classiques semblent échouer. Une vérité logique n'est pas nécessairement intéressante, donc la nécessitation ne vaut pas. Si \texttt{phi} implique \texttt{psi}, l'intérêt de \texttt{phi} ne se transmet pas à \texttt{psi}. L'intérêt de deux propositions ne se transmet pas à leur conjonction. Enfin, deux propositions logiquement équivalentes peuvent différer par leur formulation, leur ordre d'apparition et leur effet sur le sujet. Une logique de l'intéressant devrait donc être non normale ou porter sur des formes présentées, pas seulement sur leur contenu de vérité abstrait.

La piste la plus forte est dynamique : une forme intéressante ne qualifie pas seulement une possibilité dans un espace déjà donné ; elle peut modifier la relation d'accessibilité elle-même. Elle rend visible un problème, ouvre une possibilité jusque-là impensable ou fait apparaître rétrospectivement une solution comme nécessaire. L'état modal du sujet avant la rencontre n'est alors pas le même qu'après elle.

L'ennui ne fournit pas un opérateur dual simple. Si \texttt{B(phi)} signifie que \texttt{phi} ennuie, alors \texttt{B(phi)} n'est pas équivalent à \texttt{non I(phi)} : l'absence d'intérêt peut être indifférence, incompréhension ou distraction sans ennui. Réciproquement, \texttt{non B(phi)} n'implique pas \texttt{I(phi)}. L'ennui désigne plutôt une trajectoire dans laquelle une demande d'attention n'est plus satisfaite ; il peut suivre l'épuisement de l'intérêt ou motiver la recherche d'une nouvelle prise.

\begin{quote}
\textbf{Intérêt pour la thèse.} Cette formalisation relie trois noyaux du projet : l'intéressant comme transformation du sujet, l'Adjacent Possible et la nécessité inventée. Elle permet de demander si l'intéressant sélectionne des mondes accessibles, modifie leur accessibilité ou change les catégories mêmes sous lesquelles le sujet les décrit.

Elle ouvre un choix technique entre logique modale non normale, sémantique de voisinages, logique épistémique dynamique et théorie de la révision des croyances. Ce choix ne doit pas être fait par analogie : il dépendra des lois effectivement soutenues par les cas du corpus.
\end{quote}
{\footnotesize\color{metadata}\textbf{Provenance :} BUT DE LA THESE.md; projet thèse philo.pdf; De l'impossibilité de créer.pdf. Programme formel : opérateur relationnel d'intérêt indexé par sujet, horizon et temps Projet de thèse V1, nécessité inventée et problème rendu visible par sa solution, PDF p. 2 Essai sur la création, impossible productif et Adjacent Possible, PDF p. 14-27\par}

\section{Temporalité, attention et apprentissage}

\Needspace{8\baselineskip}
\subsection{L'ennui est le prix de l'excitation parce que le plaisir vient d'une variation}
{\small\color{metadata} \texttt{idea\_0004}\quad\textbullet\quad argument\quad\textbullet\quad conceptuelle\par}

Une excitation continue et sans résistance s'annule elle-même. Comme le plaisir de la vitesse vient de l'accélération et du ralentissement avant le virage, le plaisir attentionnel vient d'une variation entre tension, attente et relance. L'ennui ou la contrainte ne sont donc pas seulement des échecs du dispositif : ils constituent le contraste et le prix sans lesquels l'excitation ne peut pas être perçue.

L'ennui n'est pourtant pas la négation logique de l'intéressant. Ne pas être intéressé peut signifier rester indifférent, ne pas comprendre ou ne pas remarquer, sans éprouver l'ennui. Ne pas s'ennuyer ne signifie pas davantage être intéressé. L'ennui suppose une durée et une demande attentionnelle insatisfaite : quelque chose continue, mais ne produit plus assez de variation, de progression ou de prise. Il peut donc être la suite de l'épuisement d'un intérêt, coexister avec une attente encore ouverte ou provoquer l'exploration qui fera naître un nouvel intérêt.

Le modèle du flow fournit une première géométrie de cette relation. À compétences fixes, lorsque la difficulté augmente, on traverse successivement trois zones : l'ennui si la tâche est trop facile, l'intéressant et le flow lorsqu'elle sollicite les compétences disponibles, puis l'anxiété, l'opacité ou le décrochage lorsqu'elle les dépasse trop fortement. L'intéressant occupe ainsi une zone intermédiaire plutôt qu'un pôle simplement opposé à l'ennui.

Le parcours s'inverse si l'on fixe la tâche et que le sujet apprend : une même difficulté peut être d'abord anxiogène ou opaque, devenir intéressante puis fluide, et enfin ennuyeuse lorsqu'elle ne sollicite plus assez les capacités acquises. La zone d'intérêt se déplace donc avec les compétences.

On peut alors poser comme hypothèse que l'intéressant est cette zone de flow intermédiaire. Il s'agit de la zone structurelle d'ajustement entre compétences et difficulté, non du seul état subjectif d'absorption continue. Une hésitation ou une rupture peut rester intéressante tant que le sujet conserve une prise sur la difficulté et peut réorganiser ses attentes.

\begin{quote}
\textbf{Intérêt pour la thèse.} Cette carte permet de penser l'intéressant comme région instable entre saturation, attente et relance, et comme trajectoire relative aux compétences du sujet.
\end{quote}
{\footnotesize\color{metadata}\textbf{Provenance :} Synopsis MIT Press.doc. Synopsis MIT Press, rendu PDF p. 6-7, 10 et 27-30\par}

\Needspace{8\baselineskip}
\subsection{L'intéressant peut être opérationnalisé comme choix du meilleur prochain événement}
{\small\color{metadata} \texttt{idea\_0032}\quad\textbullet\quad méthode\quad\textbullet\quad scientifique\par}

On peut remplacer la question trop globale « cet objet est-il intéressant ? » par une décision séquentielle : quel prochain événement prolongera l'attention sans être une simple répétition ni une rupture arbitraire ? Dans le jeu proposé, le producteur choisit successivement des éléments et l'auditeur peut attendre, s'exciter, s'ennuyer ou quitter. L'intérêt devient ainsi une politique de relance évaluée par la trajectoire de l'auditeur, non une propriété attribuée après coup à la séquence complète.

Le signal le plus simple n'est pas l'excitation déclarée à chaque instant, mais la durée pendant laquelle l'auditeur accepte de rester dans le jeu. Attendre peut faire partie de l'expérience intéressante si une résolution reste possible. Le problème expérimental devient alors : quelle politique de prochains événements retarde l'abandon sans exploiter une répétition compulsive ?

Le dossier Jotney propose une variante analytique de cette opération : demander, à chaque bifurcation, ce qu'aurait fait une chanson simplement compétente, puis mesurer l'écart entre cette continuation acceptable et celle qui a effectivement été choisie. Cet écart ne suffit pas à prouver l'intérêt, mais il localise les décisions où une forme refuse l'automatisme sans tomber dans l'arbitraire.

L'écoute superposée et l'interdifference donnent deux techniques pour produire ce contrefactuel. La première transforme mentalement rythme, mélodie ou harmonie en conservant certains points d'appui ; la seconde soustrait à l'original une reprise ou une imitation ratée afin de faire apparaître la petite différence qui soutient sa qualité. Comprendre devient ainsi une activité de reconstruction, pas seulement la réception d'un effet.

\begin{quote}
\textbf{Intérêt pour la thèse.} Cette carte donne un dispositif expérimental très simple pour penser l'intéressant comme gestion temporelle de l'attention.
\end{quote}
{\footnotesize\color{metadata}\textbf{Provenance :} Excited Bored.doc; ExcitingSequences.html; Synopsis MIT Press.doc; The Mystery of Jotney Songs -full.pdf; PACHET HISTOIRE OREILLE BAT.pdf. Excited Bored, jeu producteur-auditeur et plan du livre, rendu PDF p. 1-4 ExcitingSequences.html : version courte du même jeu ; mêmes actions de l'auditeur et même objectif de retard maximal de l'abandon, sans règle ni proposition supplémentaire. Synopsis MIT Press, rendu PDF p. 27-30 Dossier Jotney, méthode contrefactuelle du prochain choix, PDF p. 2-4 Histoire d'une oreille, écoute superposée et variation intérieure, PDF p. 17-21 ; interdifference par comparaison avec reprises et imitations ratées, p. 133-135 et 259-260.\par}
{\footnotesize\color{metadata}\textbf{Références :} \parencite{pachet2018oreille}\par}

\Needspace{8\baselineskip}
\subsection{Mesurer une préférence la modifie lorsque l'exposition fait évoluer le goût}
{\small\color{metadata} \texttt{idea\_0097}\quad\textbullet\quad hypothèse\quad\textbullet\quad scientifique\par}

Si le goût résulte d'une interaction historique avec les artefacts, présenter un objet pour recueillir un jugement n'est pas une observation neutre. L'exposition ajoute une expérience, modifie les attentes et peut changer les évaluations suivantes, y compris pour des objets déjà rencontres.

La préférence doit alors être modélisée comme un processus potentiellement chaotique et dépend de l'ordre des essais. Une expérience ne cherche plus une valeur stable cachée dans le sujet ; elle doit enregistrer la trajectoire par laquelle ses jugements se transforment.

Le livre décrit plusieurs trajectoires complètes : rejet, détail qui accroche, digestion, indifférence ; ou saturation des effets locaux suivie d'une attention nouvelle à la structure globale. L'exposition ne fait donc pas seulement monter ou baisser une note de préférence : elle peut changer la dimension même sur laquelle l'auditeur écoute et juge.

\begin{quote}
\textbf{Intérêt pour la thèse.} Cette hypothèse donne une conséquence expérimentale forte au refus du goût comme dictionnaire statique.
\end{quote}
{\footnotesize\color{metadata}\textbf{Provenance :} ERCGrantPachetInterestingness.pdf; PACHET HISTOIRE OREILLE BAT.pdf. PDF p. 4-5 Histoire d'une oreille, trajectoire de goût de La Montagne, PDF p. 50-53 ; attention élargie de l'effet à la structure, p. 184-185 et 276-278 ; morceau ignoré puis devenu attirant, p. 294-295.\par}
{\footnotesize\color{metadata}\textbf{Références :} \parencite{pachet2018oreille}\par}

\Needspace{8\baselineskip}
\subsection{L'intéressant est un déclencheur de constructions qui donnent prise sur l'objet}
{\small\color{metadata} \texttt{idea\_0123}\quad\textbullet\quad définition\quad\textbullet\quad conceptuelle\par}

Un objet ne peut rester intéressant qu'en déclenchant un travail chez le sujet. Ce travail peut consister à comparer, anticiper, discriminer, former une catégorie, réviser une croyance ou construire un modèle partiel de ce qui est perçu. L'intéressant n'est donc pas seulement une stimulation reçue : il inaugure une relation dans laquelle le sujet acquiert de nouvelles prises et se transforme.

Cette proposition fournit un candidat pour le cœur de la thèse : \textbf{l'intéressant est ce qui déclenche et soutient une construction par laquelle un sujet acquiert une prise nouvelle sur un objet}. La construction est le terme médiateur entre l'attention initiale et la compréhension. Elle explique ce que l'intéressant fait, plutôt que de le définir seulement par sa position entre ennui et anxiété.

Ce travail n'est ni nécessairement conscient ni extérieurement visible. Un bébé qui regarde longuement un triangle jaune peut sembler passif tout en apprenant à distinguer des invariants de couleur, de contour ou de forme. Cet exemple constitue une hypothèse sur la construction perceptive, non un résultat expérimental établi ici. Les \emph{language games} de Luc Steels fournissent un modèle plus précis du mécanisme général : sous la pression répétée d'une tâche de discrimination et de communication, des agents construisent et coordonnent des catégories perceptives et les mots qui les désignent.

Dans l'optique du flow, ce travail explique comment l'intérêt se maintient provisoirement dans la zone intermédiaire. Chaque réussite augmente les compétences et ouvre une compréhension nouvelle ; elle rend aussi la tâche plus facile et prépare ainsi une dérive vers l'ennui. À l'inverse, si le sujet ne parvient plus à construire de modèle, les échecs s'accumulent et la relation dérive vers l'anxiété ou l'abandon. Soutenir l'intérêt exige donc que l'objet rende possible une construction dont les résultats renouvellent, plutôt qu'ils n'épuisent immédiatement, ce qui reste à comprendre.

Une forme privilégiée de cette construction est la résolution de problème à rebours. Le sujet rencontre d'abord une solution, puis infère progressivement les contraintes auxquelles elle répond, les alternatives qu'elle exclut et le problème qu'elle rend visible. Demander « de quel problème cet objet est-il la solution ? » devient ainsi une opération générale pour prolonger l'intérêt et transformer la perception en compréhension.

\paragraph{Distinction}
La passivité corporelle ou l'absence d'action visible ne sont pas la passivité cognitive. Écouter, regarder ou lire peuvent mobiliser un travail intense de modélisation. La proposition vise la passivité forte : une réception qui ne modifie ni les attentes, ni les catégories, ni les capacités d'anticipation du sujet. Une telle relation peut capturer momentanément l'attention par la seule variation des stimuli, mais elle ne peut soutenir l'intéressant sans renouvellement externe continu.

La construction ne doit pas non plus devenir le nom flatteur de toute interprétation. Elle doit produire au moins une prise contrôlable : une distinction reproductible, une prédiction, une comparaison, une possibilité d'intervention, un transfert ou un objet construit dont les réussites et les échecs peuvent être examinés. Ce critère distingue la construction d'une simple impression de profondeur.

\begin{quote}
\textbf{Intérêt pour la thèse.} Cette définition relie phénoménologie de l'attention, apprentissage, formation des concepts et pratique de l'ingénieur. Elle suggère également une méthode : mesurer l'intéressant non seulement par la durée d'exposition, mais par les distinctions, prédictions, stratégies, modèles ou artefacts que le sujet devient capable de produire au cours de la relation.
\end{quote}
{\footnotesize\color{metadata}\textbf{Provenance :} PACHET HISTOIRE OREILLE BAT.pdf; interestingness.pdf. Histoire d'une oreille, apprentissages perceptifs et construction progressive d'une écoute, PDF p. 27-34 et 281-295 Schmidhuber, apprentissage de régularités et progrès de compression, PDF p. 1-3 et 20 Le rapprochement avec les language games de Steels est une articulation proposée par la thèse\par}
{\footnotesize\color{metadata}\textbf{Références :} \parencite{pachet2018oreille,schmidhuber1997interesting,steels2005coordinating,steels2015talkingheads}\par}

\Needspace{8\baselineskip}
\subsection{Se comprendre devient intéressant lorsque l'affect est reconstruit comme un système de causes}
{\small\color{metadata} \texttt{idea\_0126}\quad\textbullet\quad hypothèse\quad\textbullet\quad conceptuelle\par}

Comprendre les raisons de ses émotions peut être intéressant au sens fort retenu par la thèse : une expérience d'abord subie et opaque déclenche la construction d'un modèle qui donne au sujet une prise nouvelle sur lui-même. Le sujet devient ici simultanément l'enquêteur, le matériau de l'enquête et l'un des effets transformés par celle-ci.

Spinoza fournit un point d'appui précis. Dans les propositions III et IV de la partie V de l'\emph{Éthique}, former une idée claire et distincte d'une affection en réduit le caractère passif. La proposition XV ajoute que « celui qui comprend ses passions et soi-même [...] ressent de la joie ». La compréhension n'est donc pas une représentation extérieure et inerte de l'émotion : elle modifie la manière dont le sujet en pâtit et peut accroître sa puissance d'agir.

Sous l'angle de la thèse, le terme décisif n'est pourtant pas d'abord le plaisir, mais la \textbf{construction}. Se comprendre demande d'identifier des éléments, de reconstruire des enchaînements causaux, de distinguer des déclencheurs et des conditions, de comparer des hypothèses et parfois d'imaginer des interventions possibles. Cette activité ressemble au travail d'un ingénieur lorsqu'il construit un agencement intelligible pour expliquer le comportement d'un système. Elle ne réduit pas pour autant la personne à une machine : la métaphore porte sur le mode d'enquête, pas sur une ontologie du sujet.

Deux titres donnent un vocabulaire suggestif à cette articulation. \emph{The Pleasure of Finding Things Out} de Richard Feynman nomme le plaisir épistémique attaché au fait de découvrir comment les choses fonctionnent. \emph{The Existential Pleasures of Engineering} de Samuel Florman nomme le plaisir de concevoir, d'agencer et de faire tenir une construction dans le monde. La compréhension de soi peut conjuguer ces deux dimensions : découvrir les causes d'un affect et construire un modèle qui permette d'agir sur la relation dont il procède.

Ces références ne démontrent pas que toute découverte ou toute construction est intéressante. Elles permettent de distinguer deux plaisirs susceptibles d'accompagner le même processus : le plaisir de trouver une intelligibilité et celui d'éprouver la prise produite par un agencement qui fonctionne.

L'objet intéressant est ici réflexif. L'émotion appartient au sujet, mais ses causes débordent son intériorité : elles peuvent engager une forme rencontrée, l'état du corps, une mémoire, une attente, une relation sociale ou une circonstance historique. La construction de soi prolonge ainsi la causalité relationnelle de \texttt{idea\_0084} et de \texttt{idea\_0125} : comprendre « pourquoi je ressens cela » revient à reconstruire le système momentané dont l'affect est un effet.

\paragraph{Distinction}
La joie de comprendre ne doit pas être confondue avec l'intéressant. Elle peut en être un effet, un signal ou une conséquence tardive. Le processus de compréhension peut rester intéressant tout en étant pénible, inquiétant ou frustrant ; inversement, une explication séduisante peut procurer du plaisir sans produire de prise réelle. L'intéressant désigne la dynamique constructive, tandis que le plaisir qualifie l'un de ses devenirs affectifs possibles.

\paragraph{Critère}
Une compréhension de soi ne vaut pas seulement parce qu'elle donne une impression de clarté. Elle doit permettre au moins une opération nouvelle : distinguer deux affects jusque-là confondus, reconstruire un enchaînement, anticiper une réaction, formuler un contre-factuel, déplacer une conduite ou reconnaître les limites du modèle construit. Ces effets fournissent des critères pour séparer la prise réflexive de la simple promesse de profondeur décrite dans \texttt{idea\_0124}.

\begin{quote}
\textbf{Intérêt pour la thèse.} Cette proposition fait de la compréhension de soi un cas privilégié de la définition centrale de l'intéressant. Elle relie directement la naturalisation des affects à l'activité de construction : rechercher les causes n'est pas seulement produire un savoir sur un phénomène psychique, c'est transformer la relation que le sujet entretient avec ce phénomène. Elle ouvre ainsi un terrain où l'intéressant, l'attention, la joie et la puissance d'agir peuvent être articulés sans être identifiés.
\end{quote}
{\footnotesize\color{metadata}\textbf{Provenance :} BUT DE LA THESE.md. Spinoza, Éthique, partie V, propositions III-IV : former une idée plus adéquate d'une affection diminue la passivité. Spinoza, Éthique, partie V, proposition XV : comprendre clairement ses passions et soi-même fait ressentir de la joie. Rapprochement proposé avec l'intéressant comme déclencheur de construction dans idea\_0123. Feynman, The Pleasure of Finding Things Out, recueil édité par Jeffrey Robbins, 1999. Florman, The Existential Pleasures of Engineering, deuxième édition, 1996 ; édition originale, 1976.\par}
{\footnotesize\color{metadata}\textbf{Références :} \parencite{spinoza1861oeuvres,macherey1995vieaffective,feynman1999pleasure,florman1996existential}\par}

\Needspace{8\baselineskip}
\subsection{L'attention observée renseigne mieux sur l'intérêt que la préférence déclarée}
{\small\color{metadata} \texttt{idea\_0011}\quad\textbullet\quad méthode\quad\textbullet\quad scientifique\par}

Les jugements déclarés réduisent une expérience temporelle à une réponse finale et sont exposés aux rationalisations. Les durées d'attention, reprises, abandons et transformations effectuées sur l'objet donnent accès à ce que celui-ci fait effectivement au sujet. Une science computationnelle de l'intéressant doit donc traiter le comportement observable comme une donnée centrale, et non comme un simple supplément au questionnaire.

\begin{quote}
\textbf{Intérêt pour la thèse.} Cette carte propose une méthode : regarder l'intéressant comme comportement situé plutôt que jugement abstrait.
\end{quote}
{\footnotesize\color{metadata}\textbf{Provenance :} ERCGrantPachetInterestingness.pdf; Paper citations updated.docx. Grant ERC PDF p. 3-4 Projet SenseOfDirection, comportement de skip, rendu PDF p. 6\par}

\Needspace{8\baselineskip}
\subsection{Une mesure d'intérêt peut diriger la découverte avant de pouvoir juger sa valeur}
{\small\color{metadata} \texttt{idea\_0028}\quad\textbullet\quad argument\quad\textbullet\quad scientifique\par}

Dans un espace de concepts et de conjectures trop vaste pour être exploré exhaustivement, l'intéressant n'est pas seulement une note attribuée au résultat final : il devient un opérateur de recherche qui décide quoi construire et examiner ensuite. Cette évaluation doit pourtant être faite immédiatement, alors que la valeur mathématique véritable d'un résultat peut n'apparaître que beaucoup plus tard.

Il faut donc combiner plusieurs indices imparfaits - plausibilité empirique, nouveauté, surprise, applicabilité, compréhensibilité et pertinence pour une tâche - plutôt que chercher une mesure unique de valeur. La sélection heuristique participe ainsi au processus de découverte qu'elle prétend évaluer.

Une même mesure peut intervenir à trois niveaux croissants. Elle peut filtrer les résultats après leur production, ordonner la recherche heuristique pour choisir les pistes à prolonger, ou définir l'espace de recherche de sorte que certaines formes seulement puissent être engendrées. Ce passage de la justification à la découverte transforme un critère d'évaluation en principe génératif. Les mesures doivent aussi pouvoir évoluer : la nouveauté ou l'applicabilité d'un concept change quand de nouveaux exemples, concepts ou conjectures entrent dans la théorie.

\begin{quote}
\textbf{Intérêt pour la thèse.} Cette carte rejoint le prologue sur les mathématiques : ce qui compte dans la recherche n'est pas seulement le résultat prouve, mais la sélection des pistes fertiles.
\end{quote}
{\footnotesize\color{metadata}\textbf{Provenance :} interestingness-ijhcs.pdf. Colton, Bundy et Walsh, usages et classification des mesures, PDF p. 15-22\par}
{\footnotesize\color{metadata}\textbf{Références :} \parencite{colton2000interestingness}\par}

\Needspace{8\baselineskip}
\subsection{Un motif fréquent n'est intéressant que par l'effet qu'il produit sur des croyances}
{\small\color{metadata} \texttt{idea\_0029}\quad\textbullet\quad argument\quad\textbullet\quad scientifique\par}

La fréquence ou la régularité statistique d'un motif ne suffit pas à établir son intérêt. Le même motif peut être trivial pour un utilisateur qui l'attend et décisif pour un autre dont il contredit les croyances. Une mesure subjective de l'intéressant doit donc représenter l'état de croyance avant la découverte et évaluer l'ampleur de sa révision, plutôt que noter le motif seul.

Une contradiction n'appelle toutefois pas toujours la même révision. Si elle heurte une croyance tenue pour une contrainte dure, elle signale d'abord une erreur possible dans les données. Si elle heurte une croyance souple et que les données sont confirmées, c'est le degré de croyance qui doit changer. L'intérêt d'un motif ne mesure donc pas seulement un écart : il distribue le doute entre observation et modèle du monde.

\begin{quote}
\textbf{Intérêt pour la thèse.} Cette carte donne une formulation opérationnelle du contexte subjectif : l'intéressant dépend de ce qu'un sujet croyait déjà.
\end{quote}
{\footnotesize\color{metadata}\textbf{Provenance :} kdd95.pdf; tkde.pdf. Silberschatz et Tuzhilin, croyances dures et souples, révision et mesure bayésienne, PDF p. 3-7 Version étendue, formalismes de degrés de croyance et processus de découverte, PDF p. 5-12\par}
{\footnotesize\color{metadata}\textbf{Références :} \parencite{silberschatz1995subjective,silberschatz1996patterns}\par}

\Needspace{8\baselineskip}
\subsection{Les séquences intéressantes mettent à jour un système de croyances}
{\small\color{metadata} \texttt{idea\_0035}\quad\textbullet\quad argument\quad\textbullet\quad scientifique\par}

Dans le séquence mining, la fréquence ne suffit pas. Les règles deviennent intéressantes lorsqu'elles confirment, contredisent ou transforment un système de croyances, puis orientent la prochaine session de fouille.

L'inattendu doit lui-même être qualifié : une règle peut dévier statistiquement, produire une prédiction différente, reposer sur des conditions inattendues, ou n'avoir aucun précédent comparable. Ces cas n'appellent pas la même révision. Le post-mining peut ainsi transformer certaines règles attendues en croyances, remplacer des croyances devenues caduques et utiliser chaque type d'écart pour construire la requête suivante. Le contexte d'évaluation se forme donc en même temps que les découvertes qu'il évalue.

\begin{quote}
\textbf{Intérêt pour la thèse.} Cette carte renforce l'idée que l'intéressant n'est pas seulement une propriété statistique : il dépend d'une histoire interpretative et d'une mémoire du chercheur.
\end{quote}
{\footnotesize\color{metadata}\textbf{Provenance :} C PKDD99.pdf. Spiliopoulou, PostMine et catégories d'unexpectedness, PDF p. 3-7\par}
{\footnotesize\color{metadata}\textbf{Références :} \parencite{spiliopoulou1999rules}\par}

\Needspace{8\baselineskip}
\subsection{Actionabilité et unexpectedness}
{\small\color{metadata} \texttt{idea\_0036}\quad\textbullet\quad distinction\quad\textbullet\quad scientifique\par}

Un motif peut être intéressant parce qu'il est inattendu, parce qu'il permet d'agir, ou les deux. Ces dimensions sont indépendantes : une information attendue peut autoriser une décision utile, tandis qu'une anomalie spectaculaire peut ne donner aucune prise à son destinataire. L'actionabilité est souvent difficile à formaliser directement ; l'inattendu ne doit donc pas être traité comme son proxy universel.

\begin{quote}
\textbf{Intérêt pour la thèse.} Cette distinction évite de réduire l'intéressant à une surprise contemplative : parfois est intéressant ce qui ouvre une action.
\end{quote}
{\footnotesize\color{metadata}\textbf{Provenance :} kdd95.pdf; tkde.pdf. Silberschatz et Tuzhilin, classification et indépendance des deux mesures, PDF p. 2-4 et 7 Version étendue, actionabilité comme notion centrale mais difficile à formaliser, PDF p. 3-5 et 12\par}
{\footnotesize\color{metadata}\textbf{Références :} \parencite{silberschatz1995subjective,silberschatz1996patterns}\par}

\Needspace{8\baselineskip}
\subsection{L'éthologie active provoque des situations au lieu de seulement observer}
{\small\color{metadata} \texttt{idea\_0070}\quad\textbullet\quad méthode\quad\textbullet\quad scientifique\par}

L'éthologie active ne se contente pas d'attendre qu'un comportement révèle une préférence. Un dispositif en boucle fermée produit une variation, observe la réaction de l'animal et adapte l'interaction suivante afin de provoquer des situations discriminantes.

Cette intervention permet de tester si certains principes de l'intéressant sont spécifiquement humains. Les chants artificiels de canaris, qui exagèrent des traits de virtuosité au-delà des performances naturelles, montrent notamment que la préférence peut porter sur une possibilité produite par l'expérience elle-même.

\begin{quote}
\textbf{Intérêt pour la thèse.} Cette carte renforce la dimension expérimentale de l'intéressant : il faut fabriquer des scènes ou l'intérêt devient observable.
\end{quote}
{\footnotesize\color{metadata}\textbf{Provenance :} ERCGrantPachetInterestingness.pdf. Grant ERC PDF p. 6-7\par}

\Needspace{8\baselineskip}
\subsection{La curiosité aide surtout lorsque la récompense externe est rare}
{\small\color{metadata} \texttt{idea\_0072}\quad\textbullet\quad hypothèse\quad\textbullet\quad scientifique\par}

Une récompense interne pour la découverte de régularités peut accélérer l'apprentissage orienté vers un but lorsque les récompenses externes sont trop rares pour guider efficacement l'exploration. Mais cet avantage n'est pas universel : lorsque le but externe fournit déjà beaucoup de signal, la curiosité peut détourner du perfectionnement de la politique utile et devenir coûteuse.

Les simulations de Schmidhuber soutiennent cette hypothèse de façon nuancée : le système curieux trouve d'abord plus vite la récompense externe, puis son avantage disparaît avec l'accumulation des exemples orientés vers le but. La curiosité n'est donc pas posée comme une vertu absolue, mais comme une stratégie dont la valeur dépend du régime de rareté du signal externe.

Dans l'expérience 2a, les résultats moyens du système curieux dépassent fortement ceux du système sans récompense interne, mais un essai curieux ne trouve jamais le but. Dans l'expérience 2b, l'avantage initial disparaît après un apprentissage plus long et le système non curieux finit même légèrement devant. Le papier présente ces résultats comme exploratoires et ne prétend pas établir une loi générale.

\begin{quote}
\textbf{Intérêt pour la thèse.} Cette proposition relie une théorie de l'intéressant à une fonction pratique : l'intéressant peut servir de guide provisoire lorsqu'aucun objectif externe ne donne encore de direction suffisante. Elle fournit aussi une limite testable aux discours qui identifient curiosité et performance.
\end{quote}
{\footnotesize\color{metadata}\textbf{Provenance :} interestingness.pdf. Schmidhuber, What's Interesting?, PDF p. 3 et 13-17 ; limites p. 20-21\par}
{\footnotesize\color{metadata}\textbf{Références :} \parencite{schmidhuber1997interesting}\par}

\Needspace{8\baselineskip}
\subsection{Une similarité culturelle dépend du corpus interrogé, pas d'une vérité terrain unique}
{\small\color{metadata} \texttt{idea\_0106}\quad\textbullet\quad distinction\quad\textbullet\quad scientifique\par}

Une similarité culturelle calculée par cooccurrence ne mesure pas une distance qui existerait indépendamment des données. Elle dépend de la portion de culture choisie, des noms reconnus, des contextes conservés et des pratiques qui ont laissé des traces dans le corpus. La comparer à une unique vérité terrain revient donc à masquer la question culturelle sous une procédure d'évaluation technique.

Cette dépendance peut devenir un paramètre plutôt qu'un défaut. Choisir un corpus, c'est choisir le monde dans lequel deux objets seront dits proches : discours de fans, presse, catalogue spécialisé ou usages d'une communauté. Plusieurs mesures peuvent être valides si elles rendent explicites les populations et relations qu'elles représentent.

\begin{quote}
\textbf{Intérêt pour la thèse.} Les parcours intéressants héritent de cette construction. Leur continuité et leurs écarts ne sont pas neutres : ils dépendent du corpus qui définit les voisinages possibles.
\end{quote}
{\footnotesize\color{metadata}\textbf{Provenance :} TBKLullyNOTES.doc. TBKLullyNOTES, similarités calculables et choix du corpus, rendu PDF p. 6\par}

\Needspace{8\baselineskip}
\subsection{Une réponse adaptée au contexte peut rester inintéressante si elle ignore l'histoire}
{\small\color{metadata} \texttt{idea\_0015}\quad\textbullet\quad distinction\quad\textbullet\quad conceptuelle\par}

Satisfaire exactement une demande est un premier niveau ; adapter la réponse au contexte présent en est un second. Mais ces deux niveaux peuvent répéter ce qui a déjà fonctionné jusqu'à l'ennui. Une réponse vraiment intéressante doit aussi tenir compte de la chronologie des interactions : ce qui a déjà été montré, appris, épuisé ou rendu possible. L'historique n'est donc pas un contexte supplémentaire, mais une condition de la relance.

\begin{quote}
\textbf{Intérêt pour la thèse.} Cette distinction donne une progression très utile pour classer les systèmes : service, recommandation, expérience temporelle.
\end{quote}
{\footnotesize\color{metadata}\textbf{Provenance :} TBKLullyNOTES.doc. TBKLullyNOTES, trois ordres d'interestingness, rendu PDF p. 3\par}

\Needspace{8\baselineskip}
\subsection{Un objet raté peut soutenir l'attention en transformant le spectateur en diagnosticien}
{\small\color{metadata} \texttt{idea\_0022}\quad\textbullet\quad hypothèse\quad\textbullet\quad conceptuelle\par}

Un objet réussi peut être consommé sans reste, tandis qu'un objet raté ouvre une activité de diagnostic : chercher pourquoi il échoue, anticiper la prochaine faute, comparer avec ce qui aurait fonctionné. L'intérêt peut alors porter moins sur la valeur de l'objet que sur l'enquête cognitive qu'il déclenche. Cela explique pourquoi vérifier qu'une chose est mauvaise peut parfois retenir davantage que recevoir passivement une chose bonne.

Il faut ainsi distinguer aimer l'objet et aimer l'activité qu'il provoque. On peut désapprouver une émission, un livre ou un style tout en recherchant l'expérience de le regarder, de le lire, de le critiquer ou de se mesurer à lui. Cette divergence n'est pas une incohérence du goût : elle indique que l'unité pertinente est parfois l'activité temporelle plutôt que l'objet évalué statiquement.

\begin{quote}
\textbf{Intérêt pour la thèse.} Cette carte contredit l'identification rapide entre intéressant, bon, beau ou réussi.
\end{quote}
{\footnotesize\color{metadata}\textbf{Provenance :} TBKLullyNOTES.doc. TBKLullyNOTES, aimer faire et mauvais candidats, rendu PDF p. 2\par}

\Needspace{8\baselineskip}
\subsection{Une relation temporelle minimale peut rendre saillant un contenu sans valeur}
{\small\color{metadata} \texttt{idea\_0023}\quad\textbullet\quad hypothèse\quad\textbullet\quad conceptuelle\par}

Un message sans valeur propre peut devenir momentanément saillant s'il installe une attente : objet énigmatique, faux contexte, date future, puis relance qui semble répondre au premier envoi. Le spam montre à échelle minimale que l'intérêt peut être produit par une relation entre occurrences plutôt que par la qualité de leur contenu. La seconde apparition transforme rétrospectivement la première en début d'une séquence.

\begin{quote}
\textbf{Intérêt pour la thèse.} Le spam est un laboratoire minuscule de capture attentionnelle : il montre comment une promesse narrative ou temporelle peut suffire à ouvrir une boucle.
\end{quote}
{\footnotesize\color{metadata}\textbf{Provenance :} TBKLullyNOTES.doc. TBKLullyNOTES, titres, fausse date et relance, rendu PDF p. 1 et 3-5\par}

\Needspace{8\baselineskip}
\subsection{Apprendre est intéressant parce que l'erreur ne se répète pas}
{\small\color{metadata} \texttt{idea\_0024}\quad\textbullet\quad argument\quad\textbullet\quad articulation\par}

Un système qui apprend devient intéressant lorsqu'il rejoue une situation similaire sans refaire la même erreur. Le plaisir vient de la perception d'une différence temporelle : quelque chose a été retenu.

Le même mécanisme vaut pour l'utilisateur. Un jeu peut maintenir l'attention en donnant un retour clair sans expliciter sa règle : le sujet essaie, distingue les réponses du système et découvre progressivement quelle transformation est attendue. L'intérêt porte alors sur le passage d'une hypothèse fausse à une stratégie meilleure, non sur la simple récompense d'une réponse correcte.

\begin{quote}
\textbf{Intérêt pour la thèse.} Cette carte donne une micro-structure de l'intéressant : répétition, écart, mémoire, amélioration.
\end{quote}
{\footnotesize\color{metadata}\textbf{Provenance :} TBKLullyNOTES.doc. TBKLullyNOTES, jeu de règle et erreur non répétée, rendu PDF p. 2 et 5\par}

\Needspace{8\baselineskip}
\subsection{Une coïncidence devient intéressante quand ses exemples forment eux-mêmes un motif}
{\small\color{metadata} \texttt{idea\_0030}\quad\textbullet\quad hypothèse\quad\textbullet\quad conceptuelle\par}

Une ressemblance isolée peut être accidentelle et ne rien autoriser. Lorsque des exemples indépendants s'accumulent, leur multiplicité devient elle-même une régularité à expliquer : la coïncidence change de statut et acquiert un pouvoir heuristique. Le seuil pertinent n'est donc pas seulement statistique ; il dépend du moment où la série d'exemples suggère un mécanisme ou une classe encore absente.

La source conservée est une note très courte attribuant cette intuition a Lenat ; elle établit la provenance de la proposition mais n'en fournit pas la démonstration complète.

\begin{quote}
\textbf{Intérêt pour la thèse.} Cette carte donne une petite hypothèse heuristique : l'intéressant apparaît quand plusieurs occurrences commencent à faire signe ensemble.
\end{quote}
{\footnotesize\color{metadata}\textbf{Provenance :} Lenat and interestingness.eml. Courriel du 30 septembre 2004 : rapprochement en deux lignes entre le concept d'intéressant de Lenat et une théorie de la coïncidence exigeant de nombreux exemples ; lien externe ancien, sans développement conserve.\par}

\Needspace{8\baselineskip}
\subsection{La relation musicale optimale n'est ni trop proche ni trop lointaine}
{\small\color{metadata} \texttt{idea\_0031}\quad\textbullet\quad hypothèse\quad\textbullet\quad conceptuelle\par}

Une relation musicale intéressante suppose une distance optimale : assez proche pour que la recommandation soit comprise et digne de confiance, assez lointaine pour qu'elle apporte une vraie nouveauté.

Cette distance ne peut pas rester fixe. Les recommandations réussies augmentent la confiance et rapprochent progressivement les profils ; mais cette convergence finit par rendre les recommandations prévisibles. Une relation durable doit donc produire de la confiance sans abolir l'altérité qui lui permet encore de surprendre.

\begin{quote}
\textbf{Intérêt pour la thèse.} Cette carte précise le problème de la recommandation comme compromis dynamique entre confiance, similarité et surprise.
\end{quote}
{\footnotesize\color{metadata}\textbf{Provenance :} An interesting musical relationship.doc. Document intégral, rendu PDF p. 1\par}

\Needspace{8\baselineskip}
\subsection{Thue-Morse : non-répétition structurée}
{\small\color{metadata} \texttt{idea\_0037}\quad\textbullet\quad exemple\quad\textbullet\quad articulation\par}

La suite de Prouhet-Thue-Morse est engendrée par une règle morphique très simple, tout en étant sans chevauchement et non ultimement périodique. Elle constitue donc un contre-exemple formel à l'assimilation du non-répétitif au hasard : une séquence peut être facile à produire, fortement structurée et pourtant éviter des répétitions triviales.

Son ubiquité dans plusieurs domaines tient notamment à l'équivalence de nombreuses définitions qui font apparaître la même structure sous des problèmes différents. L'article établit ces propriétés mathématiques, mais aucun effet esthétique ou attentionnel ; son usage pour penser l'intéressant est un transfert théorique de la thèse, pas un résultat du papier.

\begin{quote}
\textbf{Intérêt pour la thèse.} Cette carte donne un modèle formel du compromis entre répétition et variation, utile pour penser musique, scénarios, jeux et génération.
\end{quote}
{\footnotesize\color{metadata}\textbf{Provenance :} ubiq.pdf. Allouche et Shallit, The ubiquitous Prouhet-Thue-Morse séquence, pp. 1-4 et 12-13 : définition morphique, absence de chevauchements, régularité non ultimement périodique et équivalence de nombreuses caractérisations.\par}
{\footnotesize\color{metadata}\textbf{Références :} \parencite{allouche1999thueMorse}\par}

\Needspace{8\baselineskip}
\subsection{La popularité mesure une dynamique sociale, l'importance une transformation historique}
{\small\color{metadata} \texttt{idea\_0053}\quad\textbullet\quad argument\quad\textbullet\quad conceptuelle\par}

La popularité amplifie de petites différences initiales jusqu'à produire des succès inégaux et imprévisibles. Elle mesure ce que les autres font d'une œuvre, non ce que cette œuvre transforme. L'importance, au contraire, se reconnait aux bifurcations, filiations et changements de langage qu'une œuvre produit dans le temps.

Une IA entraînée sur les clics et les usages apprend donc principalement la popularité. Elle risque de stabiliser les dominances accidentelles du présent au lieu de reconnaître ce qui pourrait réorganiser durablement le paysage.

\begin{quote}
\textbf{Intérêt pour la thèse.} Cette carte relie interestingness, valeur et temporalité longue de la reconnaissance.
\end{quote}
{\footnotesize\color{metadata}\textbf{Provenance :} De l'impossibilité de créer.pdf. PDF p. 79-83\par}

\Needspace{8\baselineskip}
\subsection{L'exposition passive aux médias fabrique des modes d'attention que l'industrie ne devrait pas seule définir}
{\small\color{metadata} \texttt{idea\_0093}\quad\textbullet\quad hypothèse\quad\textbullet\quad conceptuelle\par}

L'attention n'est pas une capacité fixe que les contenus se disputent. L'exposition répétée à des médias et à leurs rythmes entraîne des manières de regarder, de changer de cible et d'attendre une stimulation. Les dispositifs industriels ne captent donc pas seulement une attention préexistante : ils contribuent a en fabriquer les modes.

L'enjeu politique et technique est d'éviter que cette évolution soit déterminée uniquement par l'optimisation de l'engagement. Des systèmes qui favorisent l'exploration et la création active pourraient exercer d'autres formes d'attention que la consommation passive de flux.

\begin{quote}
\textbf{Intérêt pour la thèse.} Cette hypothèse relie l'intéressant à la formation historique de l'attention et donne une portée normative au design des systèmes culturels.
\end{quote}
{\footnotesize\color{metadata}\textbf{Provenance :} ERCGrantPachetInterestingness.pdf; ERCInteractiveReflexions2.docx. PDF p. 1 Projet REFLEX, background, rendu PDF p. 2\par}

\Needspace{8\baselineskip}
\subsection{Une préférence déclarée peut être une image de soi plutôt qu'une anticipation du plaisir}
{\small\color{metadata} \texttt{idea\_0101}\quad\textbullet\quad hypothèse\quad\textbullet\quad conceptuelle\par}

Dire que l'on préfère les Beatles a Metallica ne décrit pas nécessairement ce qui retiendra effectivement l'attention dans une situation donnée. Une préférence déclarée peut être une croyance relativement stable sur soi, entretenue parce qu'elle soutient une identité, alors que le plaisir est produit par une séquence, un contexte et un moment.

Il faut donc distinguer au moins deux objets de mesure : l'image de ses goûts que le sujet reconnait comme sienne, et la dynamique de son attention pendant une expérience. Leur divergence n'est pas une erreur de mesure mais un phénomène à expliquer.

\begin{quote}
\textbf{Intérêt pour la thèse.} Cette proposition interdit de traiter les profils déclarés comme vérité directe du plaisir et suggère de comparer discours identitaire et trajectoire d'attention.
\end{quote}
{\footnotesize\color{metadata}\textbf{Provenance :} Synopsis MIT Press.doc. Synopsis MIT Press, rendu PDF p. 10\par}

\Needspace{8\baselineskip}
\subsection{L'abondance transforme la sélection d'objets en construction de trajectoires}
{\small\color{metadata} \texttt{idea\_0105}\quad\textbullet\quad argument\quad\textbullet\quad conceptuelle\par}

Dans un catalogue assez petit, choisir peut consister à parcourir les objets puis à nommer celui que l'on veut. Dans un catalogue immense, l'utilisateur ne connaît ni l'ensemble disponible ni, souvent, l'objet précis qu'il cherche. Une liste de résultats ou de recommandations lui restitue alors le problème sous une autre forme.

L'unité pertinente devient un programme ordonné : une trajectoire qui négocie le désir de retrouver du familier et celui de découvrir du nouveau. Un titre médiocre isolé peut être juste au bon moment parce que les éléments précédents lui donnent une fonction. L'accès doit donc optimiser des propriétés de séquence, pas seulement classer des objets indépendants.

Le seuil de l'abondance n'est pas un nombre absolu : « beaucoup » signifie déjà trop pour être appréhende par les moyens disponibles. Une technologie qui augmente l'accès sans augmenter les capacités de parcours peut donc transformer une grande collection en infini pratique, plutôt qu'en richesse effectivement explorable.

\begin{quote}
\textbf{Intérêt pour la thèse.} Cette proposition relie la surabondance culturelle à la temporalité de l'intéressant et reformule la recommandation comme composition.
\end{quote}
{\footnotesize\color{metadata}\textbf{Provenance :} Synopsis MIT Press.doc; TBKLullyNOTES.doc. Synopsis MIT Press, chapitre Active Listening, rendu PDF p. 14-20 TBKLullyNOTES, le beaucoup relatif à ce qui peut être appréhende, rendu PDF p. 8\par}

\Needspace{8\baselineskip}
\subsection{L'intéressant permet d'étudier l'attachement sans le confondre avec un jugement moral}
{\small\color{metadata} \texttt{idea\_0117}\quad\textbullet\quad distinction\quad\textbullet\quad conceptuelle\par}

Qualifier un objet ou une activité d'intéressante ne revient ni à les recommander ni à les condamner. L'intéressant peut servir de catégorie descriptive pour étudier comment une prise attentionnelle devient attente, retour volontaire, attachement, habitude ou conduite compulsive. Cette suspension provisoire du jugement moral permet d'interroger les mécanismes avant de décider si leurs effets sont désirables, nocifs ou ambivalents.

Les usages intensifs des modèles de langage et les jeux comme Tetris constituent des cas programmatiques. Au lieu de demander d'abord si leur usage améliore ou détériore la cognition, l'enquête demande ce qui fait revenir : ajustement de la difficulté, résolution incomplète, nouveauté locale, sentiment de progression, disponibilité immédiate, renforcement intermittent, personnalisation ou production d'une attente. Ces exemples ne prouvent aucun mécanisme particulier ; ils indiquent des variables à isoler et à comparer.

Il faut cependant distinguer intérêt, attachement et addiction. L'intérêt peut amorcer et entretenir une pratique sans suffire à produire une addiction clinique. Inversement, une conduite automatisée ou compulsive peut persister lorsque l'objet n'est plus vécu comme intéressant. La thèse doit donc étudier les passages et les dissociations entre ces régimes plutôt que les disposer sur une échelle unique.

\begin{quote}
\textbf{Intérêt pour la thèse.} Cette distinction fait de l'intéressant une manière d'enquêter sur notre propre comportement et notre cognition. Elle autorise l'étude commune d'objets admirés, ordinaires, triviaux ou jugés dangereux, en prenant pour données les reprises, les abandons, les durées, les anticipations et les transformations du sujet.

La neutralité n'est que méthodologique. Expliquer un mécanisme de capture ne le rend pas acceptable, et le critiquer moralement ne dispense pas de l'expliquer. L'évaluation éthique et politique intervient dans un second temps, lorsqu'on examine qui conçoit le dispositif, quelles conduites il favorise et quelles possibilités il ferme.
\end{quote}
{\footnotesize\color{metadata}\textbf{Provenance :} BUT DE LA THESE.md; TBKLullyNOTES.doc. Cadrage méthodologique : neutralité axiologique provisoire dans l'étude de l'attention et de l'attachement TBKLullyNOTES, divergence entre aimer un objet et aimer l'activité qu'il provoque, rendu PDF p. 2\par}

\Needspace{8\baselineskip}
\subsection{Un objet intéressant est soumis à deux dérives constantes vers l'ennui et l'anxiété}
{\small\color{metadata} \texttt{idea\_0122}\quad\textbullet\quad hypothèse\quad\textbullet\quad conceptuelle\par}

Un objet intéressant à l'instant \texttt{t} n'occupe pas durablement une position fixe dans la zone de flow. Il est soumis à deux dérives opposées. La première conduit vers l'ennui : la répétition, la familiarisation et l'acquisition de compétences réduisent la difficulté effective jusqu'à ce que l'objet ne sollicite plus assez le sujet. La seconde conduit vers l'anxiété : l'accumulation d'incompréhensions, d'échecs ou de frustrations réduit la prise du sujet jusqu'à faire apparaître un mur de difficulté.

On peut noter \texttt{D\_t} la difficulté effective de la tâche et \texttt{C\_t} les compétences disponibles du sujet. L'intéressant correspondrait à un intervalle de l'écart \texttt{D\_t - C\_t}, plutôt qu'à une égalité exacte. L'apprentissage tend à diminuer cet écart ; la complexification de la tâche, la fatigue ou la perte de repères tendent à l'augmenter. L'intérêt est donc un équilibre métastable entre deux sorties toujours possibles de la zone de flow.

Ce déplacement peut être suivi par le gradient de progrès plutôt que par le seul écart \texttt{D\_t - C\_t}. Tant que l'erreur diminue ou que les prises s'améliorent, la construction reste intéressante. Lorsque le gradient tend vers zéro parce que tout est maîtrisé, la relation dérive vers l'ennui ; lorsqu'il tend vers zéro malgré une erreur élevée, elle dérive vers l'anxiété. Le même gradient nul possède donc deux interprétations opposées selon le niveau de difficulté résiduelle.

L'écoute musicale rend ces deux forces particulièrement visibles. Une chanson peut devenir ennuyeuse lorsque l'auditeur en a épuisé les ressorts, par exemple après avoir appris l'algèbre de ses cadences. Inversement, une écoute imposée dont les régularités restent inaccessibles, comme peut l'être une musique dodécaphonique pour un auditeur sans prises adéquates, peut accumuler incompréhension et frustration jusqu'au rejet. Ces exemples ne qualifient pas les œuvres en elles-mêmes : ils décrivent leur relation temporelle à un auditeur donné.

\begin{quote}
\textbf{Intérêt pour la thèse.} Cette hypothèse transforme la zone statique du flow en modèle dynamique de la naissance, du maintien et de l'épuisement de l'intérêt. Elle permet de comparer une chanson, un jeu, un apprentissage ou une interaction avec un système adaptatif selon la manière dont ils compensent, exploitent ou accélèrent ces deux dérives.
\end{quote}
{\footnotesize\color{metadata}\textbf{Provenance :} De l'impossibilité de créer.pdf. PDF p. 59-62 pour le modèle du flow ; les deux dérives et leur application à l'écoute sont une hypothèse propre à la thèse\par}

\section{Forme, problème, contrainte et nécessité}

\Needspace{8\baselineskip}
\subsection{L'objet intéressant peut inventer le problème dont il apparaît ensuite comme la solution nécessaire}
{\small\color{metadata} \texttt{idea\_0085}\quad\textbullet\quad hypothèse\quad\textbullet\quad conceptuelle\par}

Certaines formes ne répondent pas à une question déjà formulée. Leur apparition fait simultanément comprendre la contrainte qu'elles résolvent. Elles semblent alors nécessaires après coup, bien qu'aucune méthode n'ait permis de les déduire avant leur existence.

On peut en tirer une définition candidate : \textbf{un objet est intéressant lorsqu'il engage le sujet à reconstruire le problème dont cet objet apparaît comme une solution}. Le travail attentionnel ne consiste alors pas seulement à reconnaître des traits, mais à inférer des contraintes, imaginer des alternatives et comprendre pourquoi la forme rencontrée réussit là où les solutions voisines échoueraient.

Cette définition s'inscrit dans la zone de flow. Si le problème et sa solution sont immédiatement évidents, la reconstruction n'exige aucun travail et dérive vers l'ennui. Si les contraintes restent impossibles à inférer, l'objet demeure opaque et la relation dérive vers l'anxiété ou l'abandon. L'intéressant occupe la région où le sujet peut progressivement formuler le problème, éprouver la difficulté et reconnaître la prise offerte par la solution.

Cette \texttt{nécessité inventée} distingue la singularité intéressante de la solution simplement correcte. L'objet reconfigure l'espace des problèmes : il ne remplit pas une place vide, il rend cette place visible en l'occupant.

Le sentiment d'optale fournit une opération perceptive pour cette reconfiguration : les alternatives voisines échouent après que la forme a été entendue, faisant apparaître simultanément la solution et les contraintes qu'elle satisfait. Cette épreuve doit toutefois être distinguée de la pseudoptale, où la familiarité seule fait passer une trajectoire répétée pour nécessaire.

\begin{quote}
\textbf{Intérêt pour la thèse.} Cette proposition relie l'intéressant à la naissance des problèmes et explique pourquoi certaines trouvailles paraissent à la fois surprenantes et évidentes. Elle fournit une forme déterminée au travail constructif du sujet : comprendre ce que la forme résout.
\end{quote}
{\footnotesize\color{metadata}\textbf{Provenance :} projet thèse philo.pdf; The Mystery of Jotney Songs.pdf; PACHET HISTOIRE OREILLE BAT.pdf. PDF p. 2 Jotney, chanson comme problème melodico-harmonique et preuve de sa solution, PDF p. 1-2 Histoire d'une oreille, sentiment d'optale obtenu par variations voisines, PDF p. 17-21 et 64-66 ; solution d'ouverture trouvée après recherche, p. 282.\par}
{\footnotesize\color{metadata}\textbf{Références :} \parencite{pachet2018oreille}\par}

\Needspace{8\baselineskip}
\subsection{Un objet peut être intéressant parce qu'il apparaît comme une solution rare d'un problème implicite}
{\small\color{metadata} \texttt{idea\_0096}\quad\textbullet\quad hypothèse\quad\textbullet\quad articulation\par}

La rareté pertinente n'est pas la faible fréquence brute d'un objet. En percevant une forme, le sujet peut reconstruire implicitement un espace combinatoire de contraintes dans lequel cette forme constitue une solution peu nombreuse, stable ou difficile à atteindre. L'objet devient intéressant parce qu'il rend sensible la petitesse de la classe de solutions dont il fait partie.

Tester cette hypothèse suppose d'inférer depuis l'objet le problème qu'il semble résoudre, puis de comparer ses propriétés aux autres solutions légales. La rareté est donc relative à une formulation du problème, non absolue dans le monde.

Le compromis est délicat : une contrainte trop générale laisse tant de solutions qu'elle ne produit aucune rareté perceptible; une contrainte trop particulière peut fabriquer artificiellement une solution unique sans lui donner de portée. La rareté devient signifiante lorsque des contraintes assez génériques relient globalement les parties de l'objet tout en laissant une petite classe de solutions.

\begin{quote}
\textbf{Intérêt pour la thèse.} Cette proposition donne une interprétation computationnelle à la nécessité inventée et à l'autostabilité des formes.
\end{quote}
{\footnotesize\color{metadata}\textbf{Provenance :} ERCGrantPachetInterestingness.pdf; Synopsis MIT Press.doc; TBKLullyNOTES.doc. PDF p. 4 Synopsis MIT Press, rendu PDF p. 22-23 TBKLullyNOTES, palindromes comme solutions rares, rendu PDF p. 6\par}

\Needspace{8\baselineskip}
\subsection{Produire de l'intéressant exige souvent d'articuler sampling et résolution de contraintes}
{\small\color{metadata} \texttt{idea\_0109}\quad\textbullet\quad hypothèse\quad\textbullet\quad articulation\par}

Produire un objet intéressant demande souvent deux opérations hétérogènes. Le sampling en marche avant ouvre des possibles en prolongeant localement un matériau, sans exiger qu'un but complet soit connu. La résolution de contraintes traite au contraire chaque choix comme une étape d'un problème global : elle vérifie les continuations, revient en arrière, répare ou sélectionne les chemins encore compatibles avec une forme à atteindre.

Chacune échoue seule d'une manière caractéristique. Le sampling pur peut produire une suite d'événements localement plausibles mais incapable d'aboutir à une fin, un mètre ou une organisation globale. Le solveur pur peut énumérer toutes les solutions légales sans savoir lesquelles sont singulières ou intéressantes. La génération de l'intéressant exige alors leur articulation : proposer du matériau, faire apparaître ou imposer un problème, tester les contraintes, puis relancer la génération depuis ce que la résolution a rendu possible.

Cette articulation n'impose pas toujours un plan antérieur à la production. Dans le doodling, le sampling peut faire apparaître après coup la forme ou le problème qui orientera les étapes suivantes. Ailleurs, une contrainte explicite précède le premier geste. L'hypothèse porte sur la coopération des deux régimes, pas sur leur ordre fixe ni sur tous les cas possibles de l'intéressant.

NetNeg fournit un précédent scientifique concret : chaque réseau neuronal émet une distribution de prochains sons et les agents cherchent ensuite une paire qui maximise les préférences apprises tout en respectant les contraintes du contrepoint. Le résultat retenu est réinjecté comme nouveau contexte. Le système ne démontre pas une théorie générale de l'intéressant, mais il réalise déjà la boucle proposition, négociation, contrainte et relance décrite ici.

\paragraph{Statut de la source}
Cette carte est une synthèse pour la thèse, non une conclusion explicite de \texttt{Hidden Biases}. Ce papier établit la tension computationnelle entre sampling autoregressif local et contraintes globales. Les notes sur le doodling et le grant ERC fournissent le lien proprement dit avec la production de formes intéressantes. NetNeg constitue une construction scientifique antécédente que la thèse peut réinterpréter, sans lui attribuer cette généralisation philosophique.

\begin{quote}
\textbf{Intérêt pour la thèse.} La proposition donne un mécanisme commun à l'exploration sans but, au sens de la direction et à l'objet perçu comme solution rare. Elle suggère aussi une architecture expérimentale : comparer sampling seul, résolution seule et boucle hybride, puis mesurer la diversité, la légalité et l'intérêt des objets obtenus.
\end{quote}
{\footnotesize\color{metadata}\textbf{Provenance :} Hidden Biases in Conditioning Autoregressive Models.pdf; Notes thèse.pdf; ERCGrantPachetInterestingness.pdf; aaai99A.pdf; netneg.pdf. Hidden Biases, insuffisance du sampling local pour les propriétés globales, PDF p. 1-3 et 10-14 Notes thèse, doodling comme échantillonnage peu contraint, PDF p. 1 ERC Grant, objet comme solution rare, harmonisations légales et exploration combinatoire, PDF p. 4-6 NetNeg, prédictions du réseau négociées avec les contraintes de contrepoint, PDF p. 3-6 NetNeg développe, architecture et comparaison des modules, PDF p. 9-18\par}
{\footnotesize\color{metadata}\textbf{Références :} \parencite{pachet2026biases,goldman1999netneg}\par}

\Needspace{8\baselineskip}
\subsection{Créer commence par une impossibilité}
{\small\color{metadata} \texttt{idea\_0008}\quad\textbullet\quad argument\quad\textbullet\quad conceptuelle\par}

La création ne peut pas être simplement commandée. Dire "crée" expose une contradiction : si le geste est programmable, il n'est déjà plus création au sens fort ; s'il est vraiment créateur, il excède l'ordre donné.

\begin{quote}
\textbf{Intérêt pour la thèse.} Cette carte peut devenir un pivot philosophique : l'intéressant naît peut-être dans la zone ou une tâche ne peut pas être exécutée directement.
\end{quote}
{\footnotesize\color{metadata}\textbf{Provenance :} De l'impossibilité de créer.pdf. PDF p. 3, 9, 19-20 et 94-96\par}

\Needspace{8\baselineskip}
\subsection{Le nouveau devient accessible localement sans avoir été concevable à l'avance}
{\small\color{metadata} \texttt{idea\_0009}\quad\textbullet\quad argument\quad\textbullet\quad conceptuelle\par}

L'Adjacent Possible n'est pas un catalogue de possibilités qui attendraient d'être sélectionnées. Il désigne les transformations accessibles depuis l'état historique d'un système. Ouvrir une possibilité fait apparaître de nouvelles fonctions et de nouvelles relations qui n'étaient pas seulement inconnues, mais impossibles à décrire auparavant.

La nouveauté est donc imprévisible sans être miraculeuse : elle émerge par bifurcations locales dans un système contraint. Si l'on pouvait en dresser la carte complète à l'avance, l'espace serait déjà donné et il n'y aurait plus d'Adjacent Possible au sens fort.

\begin{quote}
\textbf{Intérêt pour la thèse.} Cette carte articule création, contrainte et rareté. L'objet intéressant peut être pensé comme un point trouvé dans un voisinage de possibles, ni évident ni arbitraire.
\end{quote}
{\footnotesize\color{metadata}\textbf{Provenance :} De l'impossibilité de créer.pdf; Notes thèse.pdf. De l'impossibilité de créer, PDF p. 20-25\par}

\Needspace{8\baselineskip}
\subsection{Produire sans but explicite peut faire apparaître le but après coup}
{\small\color{metadata} \texttt{idea\_0012}\quad\textbullet\quad hypothèse\quad\textbullet\quad conceptuelle\par}

Dans le gribouillage, le producteur n'exécute pas nécessairement une forme déjà conçue. Il échantillonne des gestes avec peu de direction explicite, puis reconnait dans certains résultats une forme qui mérite d'être poursuivie. La sélection peut ainsi précéder l'intention : un but local apparaît après la production au lieu de la commander.

Le noodling musical montre aussi que l'exploration la moins dirigée n'est pas neutre : elle retombe sur des gestes et des enchaînements familiers qui forment une signature personnelle reconnaissable. L'échantillonnage révèle donc à la fois des possibles et les attracteurs déjà incorporés par le producteur.

\begin{quote}
\textbf{Intérêt pour la thèse.} Cette carte donne un exemple simple pour relier IA générative, création humaine et exploration d'un espace de possibles.
\end{quote}
{\footnotesize\color{metadata}\textbf{Provenance :} Notes thèse.pdf; Synopsis MIT Press.doc; PACHET HISTOIRE OREILLE BAT.pdf. Notes thèse, doodling comme échantillonnage peu contraint en dessin et en musique, PDF p. 1 Synopsis MIT Press, rendu PDF p. 37 et 40 Histoire d'une oreille, noodling sans but et signature générative personnelle, PDF p. 279-281.\par}
{\footnotesize\color{metadata}\textbf{Références :} \parencite{pachet2018oreille}\par}

\Needspace{8\baselineskip}
\subsection{Créer n'ajoute rien au monde physique}
{\small\color{metadata} \texttt{idea\_0042}\quad\textbullet\quad argument\quad\textbullet\quad conceptuelle\par}

Du point de vue physique, créer ne signifie jamais ajouter de la matière, de l'énergie ou une forme durable ex nihilo. Ce que nous appelons création est une transformation locale et provisoire dans un monde de conservation et de dissipation.

\begin{quote}
\textbf{Intérêt pour la thèse.} Cette carte donne un socle anti-romantique : la création est une fiction psychologique greffée sur des transformations matérielles.
\end{quote}
{\footnotesize\color{metadata}\textbf{Provenance :} De l'impossibilité de créer.pdf. PDF p. 9-14\par}

\Needspace{8\baselineskip}
\subsection{L'impossible est productif}
{\small\color{metadata} \texttt{idea\_0043}\quad\textbullet\quad argument\quad\textbullet\quad conceptuelle\par}

Les impossibilités les plus fécondes ne sont pas des murs extérieurs : elles naissent des ambitions internes d'un système. Plus un langage est expressif, plus il produit d'indécidable ; plus une méthode est efficace sur une classe de problèmes, plus elle renonce à l'universalité. La puissance se paie par une zone d'interdiction.

La création utilise intelligemment cette incompatibilité. Elle ne supprime pas la limite mais construit une réponse singulière qui la rend visible et sensible. Toute création durable peut ainsi être lue comme une réponse à une impossibilité bien posée.

\begin{quote}
\textbf{Intérêt pour la thèse.} Cette carte renverse la valeur de la limite : l'impossibilité n'est pas l'ennemie de la création, elle en est parfois le moteur.
\end{quote}
{\footnotesize\color{metadata}\textbf{Provenance :} De l'impossibilité de créer.pdf. PDF p. 14-20\par}

\Needspace{8\baselineskip}
\subsection{On ne se chatouille pas soi-même}
{\small\color{metadata} \texttt{idea\_0045}\quad\textbullet\quad argument\quad\textbullet\quad conceptuelle\par}

Le cerveau soustrait les effets sensoriels qu'il prédit de ses propres actions : c'est pourquoi on ne peut pas se chatouiller soi-même. Mais un léger délai entre l'action et son effet suffit à faire réapparaître l'étrangeté et le chatouillement.

La surprise exige donc moins un autre sujet qu'une altération de la boucle de prédiction. Le créateur, condamne à anticiper ses gestes, doit introduire du délai, de l'opacité ou une altérité artificielle pour que sa propre production puisse lui échapper.

\begin{quote}
\textbf{Intérêt pour la thèse.} Cette carte donne une analogie puissante : l'intéressant suppose peut-être un écart que le sujet ne peut pas se donner intégralement lui-même.
\end{quote}
{\footnotesize\color{metadata}\textbf{Provenance :} De l'impossibilité de créer.pdf. PDF p. 36-39\par}

\Needspace{8\baselineskip}
\subsection{Le créateur doit se dédoubler en générateur et juge}
{\small\color{metadata} \texttt{idea\_0046}\quad\textbullet\quad argument\quad\textbullet\quad conceptuelle\par}

La création exige deux fonctions difficiles à tenir ensemble : relâcher les filtres pour produire, puis juger sévèrement. Le couple, l'alter ego ou le système interactif permettent de distribuer ces fonctions. Une machine infatigable peut engendrer de nombreuses variations tandis que l'utilisateur reconnait rapidement celles qui lui ressemblent tout en allant là où il n'aurait pas su aller seul.

\begin{quote}
\textbf{Intérêt pour la thèse.} Cette carte relie création, altérité, collaboration et systèmes réflexifs.
\end{quote}
{\footnotesize\color{metadata}\textbf{Provenance :} De l'impossibilité de créer.pdf; ERCInteractiveReflexions2.docx; TBKLullyNOTES.doc. PDF p. 39-43 Projet REFLEX, hypothèse sur l'évaluation de ses propres artefacts, rendu PDF p. 5 TBKLullyNOTES, distribution des rôles utilisateur-Continuator, rendu PDF p. 6-7\par}

\Needspace{8\baselineskip}
\subsection{Les inventions impensables perdent leur force en advenant}
{\small\color{metadata} \texttt{idea\_0047}\quad\textbullet\quad argument\quad\textbullet\quad conceptuelle\par}

Certains objets sont désirables tant qu'ils restent impossibles. Leur réalisation ferme l'espace imaginaire qui les rendait fascinants et les rend presque triviaux.

\begin{quote}
\textbf{Intérêt pour la thèse.} Cette carte lie l'intéressant à la promesse, à l'attente et à l'indéterminé, non seulement à l'objet réalise.
\end{quote}
{\footnotesize\color{metadata}\textbf{Provenance :} De l'impossibilité de créer.pdf. PDF p. 43-45\par}

\Needspace{8\baselineskip}
\subsection{La contrainte littéraire fait surgir du nouveau sans miracle}
{\small\color{metadata} \texttt{idea\_0051}\quad\textbullet\quad exemple\quad\textbullet\quad articulation\par}

Chez Perec, écrire consiste à choisir une règle, observer, classer, recombiner. La nouveauté surgit de la réorganisation de matériaux connus sous contraintes nouvelles.

\begin{quote}
\textbf{Intérêt pour la thèse.} Cette carte fournit un exemple littéraire fort de création non magique.
\end{quote}
{\footnotesize\color{metadata}\textbf{Provenance :} De l'impossibilité de créer.pdf; Synopsis MIT Press.doc. PDF p. 62-68 Synopsis MIT Press, rendu PDF p. 22 et 36\par}

\Needspace{8\baselineskip}
\subsection{La surabondance peut fermer des possibles en éliminant la diversité des manières de faire}
{\small\color{metadata} \texttt{idea\_0074}\quad\textbullet\quad hypothèse\quad\textbullet\quad conceptuelle\par}

Produire davantage ne signifie pas ouvrir davantage de possibilités. Sous la pression des résultats immédiats, des formes différentes peuvent converger vers une seule manière efficace de fonctionner. L'abondance d'objets masque alors une réduction des cadres, des attentes et des fondations à partir desquels une porte nouvelle pourrait s'ouvrir.

La crise de la qualité ne viendrait donc pas d'une disparition du talent mais d'une homogénéisation silencieuse. Faute de pouvoir créer autrement, le système produit plus de variantes semblables ; la prolifération devient le symptôme de portes qui se ferment.

\begin{quote}
\textbf{Intérêt pour la thèse.} Cette hypothèse distingue quantité de nouveauté et diversité epistemique. Elle permet d'analyser pourquoi une offre immense peut appauvrir l'expérience de l'intéressant.
\end{quote}
{\footnotesize\color{metadata}\textbf{Provenance :} De l'impossibilité de créer.pdf. PDF p. 6-8 et 25-27\par}

\Needspace{8\baselineskip}
\subsection{L'adjacent statistique prolonge les couloirs existants sans ouvrir l'Adjacent Possible}
{\small\color{metadata} \texttt{idea\_0075}\quad\textbullet\quad argument\quad\textbullet\quad articulation\par}

Un modèle génératif explore efficacement ce qui est déjà fréquent, combinable et bien représenté : un adjacent statistique. L'Adjacent Possible de Kauffman désigne au contraire une transformation qui fait apparaître des fonctions et des relations impossibles à énumérer avant son ouverture.

La vitesse d'exploration d'un espace donné ne doit donc pas être confondue avec l'ouverture d'un nouvel espace. L'IA parcourt et densifie des couloirs existants ; elle ne franchit une porte au sens fort que si son usage dans un système plus large fait émerger une fonction qui n'était pas déjà définie par son espace de calcul.

Introduire un élément externe ne suffit toutefois pas : utilise seulement pour satisfaire une contrainte, il produit une anomalie spectaculaire sans ouvrir un nouveau régime cohérent. L'innovation demande que l'élément extérieur soit intégré au modèle cible et transforme ses relations, pas qu'il soit simplement ajoute.

\begin{quote}
\textbf{Intérêt pour la thèse.} Cette distinction fournit un critère précis pour discuter la nouveauté générative sans opposer sommairement humain et machine.
\end{quote}
{\footnotesize\color{metadata}\textbf{Provenance :} De l'impossibilité de créer.pdf; Paper citations updated.docx. PDF p. 22-25 Projet SenseOfDirection, mécanismes explicites d'innovation, rendu PDF p. 5\par}

\Needspace{8\baselineskip}
\subsection{Réifier une idée crée un objet de pensée qui efface rétrospectivement son invention}
{\small\color{metadata} \texttt{idea\_0076}\quad\textbullet\quad argument\quad\textbullet\quad conceptuelle\par}

Nommer et stabiliser une abstraction permet de la décrire, la comparer et opérer sur elle comme sur une chose. Cette réification ouvre des opérations auparavant impensables : transformer les problèmes eux-mêmes en objets a par exemple rendu concevables des machines visant des classes de problèmes.

Mais le succès de l'invention conceptuelle en efface la rupture. Une fois le nouvel objet de pensée installe, il devient difficile de penser sans lui et son apparition semble avoir été inévitable. La création d'un concept produit ainsi sa propre illusion rétrospective de trivialité.

\begin{quote}
\textbf{Intérêt pour la thèse.} Cette proposition décrit une forme de création qui ne porte pas d'abord sur des artefacts, mais sur les unités avec lesquelles une recherche peut penser et agir.
\end{quote}
{\footnotesize\color{metadata}\textbf{Provenance :} De l'impossibilité de créer.pdf. PDF p. 31-33\par}

\Needspace{8\baselineskip}
\subsection{La créativité peut venir d'une levée provisoire des filtres plutôt que d'une capacité ajoutée}
{\small\color{metadata} \texttt{idea\_0077}\quad\textbullet\quad hypothèse\quad\textbullet\quad conceptuelle\par}

Les mécanismes cognitifs qui catégorisent, simplifient et rendent la perception fluide masquent aussi des détails et des relations inhabituelles. Affaiblir provisoirement certains filtres peut donner accès à une information plus littérale et modifier la stratégie avec laquelle un sujet dessine ou résout un problème.

La nouveauté n'exige donc pas toujours une puissance supplémentaire. Elle peut apparaître lorsqu'un automatisme utile cesse momentanément d'imposer ses solutions. Cette fragilisation ne garantit ni originalité ni qualité, et les résultats expérimentaux invoques restent modestes et discutes : elle ouvre des possibilités, elle ne produit pas une œuvre.

\begin{quote}
\textbf{Intérêt pour la thèse.} Cette hypothèse inverse la conception cumulative de la créativité et en conserve les limites empiriques, sans transformer la desinhibition en recette.
\end{quote}
{\footnotesize\color{metadata}\textbf{Provenance :} De l'impossibilité de créer.pdf. PDF p. 39-41\par}

\Needspace{8\baselineskip}
\subsection{La nouveauté générative peut résider dans une nouvelle relation entre forme et matière}
{\small\color{metadata} \texttt{idea\_0080}\quad\textbullet\quad distinction\quad\textbullet\quad conceptuelle\par}

Les catégories d'original, reprise et plagiat supposent des œuvres finies que l'on peut comparer. Les systèmes génératifs dissocient plus facilement la matière et la forme : une orchestration migre, un style s'applique à un matériau sans source unique, une structure subsiste sans ses éléments reconnaissables.

Ces productions ne sont ni des inventions absolues ni de simples copies. Leur nouveauté est relationnelle et hylemorphique : elle tient à une recomposition des rapports entre formes et matières héritées. Il faut donc évaluer un régime d'apparition, pas seulement calculer la ressemblance entre deux objets.

Le \texttt{templagiarism} fournit un cas limite : le générateur remplit avec une nouvelle texture une structure empruntée à une œuvre cible. La sortie peut être cohérente, mais la forme ne naît pas de son matériau. Une nouveauté plus forte exige que forme et contenu se contraignent et se transforment réciproquement pendant la génération.

\begin{quote}
\textbf{Intérêt pour la thèse.} Cette distinction fournit des catégories plus fines pour analyser ce que l'IA rend nouveau sans lui attribuer une création ex nihilo.
\end{quote}
{\footnotesize\color{metadata}\textbf{Provenance :} De l'impossibilité de créer.pdf; Paper citations updated.docx. PDF p. 87-88 Projet SenseOfDirection, templagiarism et structure adaptée au contenu, rendu PDF p. 2-6\par}

\Needspace{8\baselineskip}
\subsection{Fougue, foucade et fulgurance distinguent potentiel, rupture et nouveauté accomplie}
{\small\color{metadata} \texttt{idea\_0081}\quad\textbullet\quad distinction\quad\textbullet\quad conceptuelle\par}

La fougue est une énergie orientée vers le futur qui promet du nouveau sans lui donner de forme. La foucade est une rupture locale et surprenante qui advient mais reste instable ou sans développement. La fulgurance est brève elle aussi, mais sa forme s'impose comme juste et rétrospectivement nécessaire dès son apparition.

Ces régimes ne se distinguent pas par une quantité d'originalité, mais par leur rapport au temps et à l'accomplissement. Une puissance générative abondante peut manifester de la fougue et produire des foucades sans atteindre la fulgurance que l'on souhaite retenir, transmettre ou rejouer.

\begin{quote}
\textbf{Intérêt pour la thèse.} Cette distinction évite de confondre capacité de produire, surprise immédiate et proposition qui tient. Elle peut devenir un outil d'évaluation des cartes elles-mêmes.
\end{quote}
{\footnotesize\color{metadata}\textbf{Provenance :} De l'impossibilité de créer.pdf. PDF p. 88-90\par}

\Needspace{8\baselineskip}
\subsection{Créer est peut-être d'abord imposer une fin à un processus qui pourrait continuer}
{\small\color{metadata} \texttt{idea\_0082}\quad\textbullet\quad hypothèse\quad\textbullet\quad conceptuelle\par}

Un processus créatif avance le plus souvent sans connaître sa conclusion. Rien dans ses décisions locales n'indique nécessairement quand s'arrêter, alors que la fin retroagit sur la signification de tout ce qui précède. Trouver un second thème, développer ou varier sont déjà des manières de préparer une clôture.

L'ennui peut ainsi venir moins de la pauvreté du contenu que de l'absence de perspective de fin. Achever une œuvre exige une décision extérieure au processus, un renoncement à ce qui pourrait encore advenir. La compétence créative ne serait donc pas seulement de produire des possibilités, mais de leur imposer une forme en sachant s'arrêter.

\begin{quote}
\textbf{Intérêt pour la thèse.} Cette hypothèse relie structure, attention et création. Elle corrige le privilège romantique accordé au commencement en faisant de la clôture une opération créative centrale.
\end{quote}
{\footnotesize\color{metadata}\textbf{Provenance :} De l'impossibilité de créer.pdf; ERCInteractiveReflexions2.docx; Paper citations updated.docx. PDF p. 96-101 Projet REFLEX, note de travail sur le stopping problem, rendu PDF p. 5 Projet SenseOfDirection, critique des séquences sans début ni fin, rendu PDF p. 1-3\par}

\Needspace{8\baselineskip}
\subsection{La génération a besoin d'un sens de la direction}
{\small\color{metadata} \texttt{idea\_0016}\quad\textbullet\quad argument\quad\textbullet\quad articulation\par}

Les systèmes génératifs imitent souvent des styles locaux, mais peinent à produire des objets structures avec début, fin, climax, cohérence et direction. Une suite de transitions plausibles ne suffit pas : la forme doit présenter des corrélations ordonnées à plusieurs échelles, de la reprise locale jusqu'aux relations entre grandes parties.

Le sens de la direction ne désigne donc pas nécessairement un plan explicite fixe à l'avance. Il est l'impression qu'une séquence conserve, transforme et prépare ses matériaux de telle sorte que les événements tardifs répondent aux précédents. Imposer le patron d'une œuvre existante donne une structure superficielle; il faut que la structure produite corresponde au contenu qu'elle organise.

Le problème possède aussi une forme computationnelle précise : une exigence portant sur la fin, le mètre ou une position future couple les choix présents à l'ensemble de leurs continuations possibles. Une bonne probabilité locale ne dit donc pas si le chemin pourra satisfaire la forme globale, ni s'il restera probable une fois cette forme imposée.

\begin{quote}
\textbf{Intérêt pour la thèse.} Cette carte relie l'intéressant à la structure temporelle globale : un objet intéressant doit savoir où il va, ou au moins donner cette impression.
\end{quote}
{\footnotesize\color{metadata}\textbf{Provenance :} Paper citations updated.docx; Hidden Biases in Conditioning Autoregressive Models.pdf. Projet SenseOfDirection, résumé et challenge, rendu PDF p. 1-3 Projet SenseOfDirection, corrélations longues et synthèse, rendu PDF p. 5-6 Hidden Biases, contraintes de forme globales et dépendance aux continuations futures, PDF p. 1-3 et 12-14\par}
{\footnotesize\color{metadata}\textbf{Références :} \parencite{pachet2026biases}\par}

\Needspace{8\baselineskip}
\subsection{Une génération contrainte peut être validée tout en biaisant la distribution cible}
{\small\color{metadata} \texttt{idea\_0017}\quad\textbullet\quad argument\quad\textbullet\quad scientifique\par}

Satisfaire une contrainte ne signifie pas échantillonner exactement le modèle initial conditionné par cette contrainte. Recherche heuristique, repondération locale, reranking, rejet ou architecture spécialisée peuvent produire des objets valides tout en modifiant la loi cible. Le biais est ici inférentiel et non hérite des données d'apprentissage : le modèle est fixe, mais la procédure de contrôle change ce qui est effectivement produit.

Cette distorsion prend deux formes indépendantes. La procédure peut conserver toutes les solutions admissibles mais leur attribuer de mauvais poids; elle peut aussi perdre du support, lorsqu'un préfixe élimine rend certaines solutions définitivement inaccessibles. Greedy, top-k, top-p et beam search peuvent provoquer ce second effet en donnant une probabilité nulle à une bifurcation pourtant nécessaire à une complétion valide.

Ce n'est pas seulement un défaut d'implémentation. Pour la classe générale des modèles autoregressifs succincts, le MAP exact est NP-difficile, la normalisation conditionnelle exacte est \#P-difficile, et même une approximation à moins d'un bit de surprisal moyen par token reste NP-difficile dans le pire cas. Ces résultats ne mesurent cependant pas le biais d'un système concret et ne disent pas que toute instance pratique est difficile.

\begin{quote}
\textbf{Intérêt pour la thèse.} Cette carte donne un ancrage technique à l'intuition philosophique : une contrainte n'est pas un simple filtre après coup. La manière de l'imposer transforme la probabilité des formes et parfois l'espace même de ce qui peut apparaître. Une production convaincante ne prouve donc ni couverture des possibles ni fidélité au modèle conditionné.
\end{quote}
{\footnotesize\color{metadata}\textbf{Provenance :} Hidden Biases in Conditioning Autoregressive Models.pdf. Hidden Biases, distinction entre génération pratique et conditionnement exact, PDF p. 1-3 Hidden Biases, dureté du conditionnement et de l'approximation, PDF p. 4-9 Hidden Biases, perte de support, inpainting et limites des résultats, PDF p. 9-14\par}
{\footnotesize\color{metadata}\textbf{Références :} \parencite{pachet2026biases}\par}

\Needspace{8\baselineskip}
\subsection{Créer est difficile parce que produire et évaluer l'intéressant ne coïncident pas}
{\small\color{metadata} \texttt{idea\_0067}\quad\textbullet\quad hypothèse\quad\textbullet\quad articulation\par}

Les propriétés que nous reconnaissons dans les objets intéressants ne sont pas forcément celles que nous produisons naturellement lorsque nous créons. à cet écart entre perception et action s'ajoute une difficulté réflexive : le créateur doit évaluer son artefact tout en le produisant et anticiper un jugement qui peut changer au contact du résultat.

La page blanche ne vient donc pas seulement d'un manque de technique. Même un sujet capable d'exécuter de nombreuses actions peut ignorer lesquelles conduiront à une forme qu'il jugera intéressante, puis ne pas savoir si son jugement actuel survivra à l'étape suivante. Une assistance pertinente doit rendre cet écart explorable sans prétendre connaître à l'avance la bonne solution.

\begin{quote}
\textbf{Intérêt pour la thèse.} Cette carte formule un problème central entre goût, production et assistance logicielle, directement lié à la question de l'intéressant pour soi.
\end{quote}
{\footnotesize\color{metadata}\textbf{Provenance :} ERCGrantPachetInterestingness.pdf; ERCInteractiveReflexions2.docx. Grant ERC PDF p. 5 Projet REFLEX, hypothèse sur l'anticipation du goût, rendu PDF p. 5\par}

\Needspace{8\baselineskip}
\subsection{Satisfaire toutes les contraintes syntaxiques ne permet pas de distinguer une bonne solution}
{\small\color{metadata} \texttt{idea\_0098}\quad\textbullet\quad argument\quad\textbullet\quad articulation\par}

Un solveur peut énumérer rapidement toutes les harmonisations conformes aux règles d'un style sans disposer d'aucun critère pour séparer les solutions musicales des solutions seulement légales. La correction syntaxique définit l'espace admissible mais ne l'ordonne pas selon la singularité, l'intérêt ou la valeur.

Le problème de l'intéressant commence donc là où la satisfaction des contraintes s'arrête. Il exige des propriétés structurelles, notamment à long terme, qui ne sont pas nécessairement connues avant l'exploration des solutions.

\begin{quote}
\textbf{Intérêt pour la thèse.} Cette proposition fournit un cas technique net de la différence entre forme correcte et forme qui mérite l'attention.
\end{quote}
{\footnotesize\color{metadata}\textbf{Provenance :} ERCGrantPachetInterestingness.pdf; Synopsis MIT Press.doc. PDF p. 5 Synopsis MIT Press, rendu PDF p. 9 et 22\par}

\Needspace{8\baselineskip}
\subsection{La meilleure solution peut être simple tout en exigeant un chemin de création singulier}
{\small\color{metadata} \texttt{idea\_0099}\quad\textbullet\quad hypothèse\quad\textbullet\quad articulation\par}

Dans la composition, les solutions retenues peuvent sembler plus simples que les essais abandonnés tout en ayant été les plus difficiles à trouver. Leur évidence finale ne conserve aucune trace des bifurcations qui ont permis de les atteindre. La difficulté porte sur le chemin de recherche, non sur la complexité visible du résultat.

Trouver cette solution demande aussi de se distancier de son propre objet afin de le juger comme une possibilité parmi d'autres. La trajectoire créative est singulière parce que chaque essai transforme à la fois l'espace exploré et les critères du créateur.

\begin{quote}
\textbf{Intérêt pour la thèse.} Cette hypothèse relie simplicité, difficultuosité et nécessité rétrospective sans faire de la complexité une mesure directe de la valeur.
\end{quote}
{\footnotesize\color{metadata}\textbf{Provenance :} ERCGrantPachetInterestingness.pdf. PDF p. 5\par}

\Needspace{8\baselineskip}
\subsection{La difficulté d'une contrainte dépend de l'état conjoint du modèle et de la forme}
{\small\color{metadata} \texttt{idea\_0100}\quad\textbullet\quad argument\quad\textbullet\quad scientifique\par}

Une contrainte simple ou même régulière n'est pas automatiquement facile à imposer. Sa tractabilité dépend du couple forme-modèle. Si le modèle possède un état suffisant borne, comme un modèle de Markov d'ordre fini, et si la contrainte est reconnue par un automate fini, leur produit fournit un état conjoint borne : programmation dynamique, échantillonnage conditionnel exact et MAP exact deviennent alors calculables en temps polynomial.

Pour un modèle autoregressif général, l'automate de la contrainte ne résume pas la dépendance du modèle a tout le préfixe. Chaque décision requiert alors la masse des continuations futures encore compatibles, sans statistique suffisante bornée générique. La difficulté ne sépare donc pas simplement local et global, ni Markov et contrainte : elle sépare les couples qui admettent une recursion d'état bornée de ceux qui exigent de sommer ou d'optimiser globalement sur les suffixes.

\begin{quote}
\textbf{Intérêt pour la thèse.} Cette proposition relie une limite technique précise des générateurs à la question philosophique du sens de la direction. Elle évite aussi une conclusion trop forte : la combinatoire d'une forme globale n'interdit pas toujours un contrôle exact; tout dépend de la représentation conjointe du modèle et de la contrainte.
\end{quote}
{\footnotesize\color{metadata}\textbf{Provenance :} ERCGrantPachetInterestingness.pdf; Paper citations updated.docx; Hidden Biases in Conditioning Autoregressive Models.pdf. PDF p. 6 Projet SenseOfDirection, corrélations longues et contraintes globales, rendu PDF p. 3-5 Hidden Biases, opposition entre état borne et contexte autoregressif général, PDF p. 3 et 10-14\par}
{\footnotesize\color{metadata}\textbf{Références :} \parencite{pachet2026biases}\par}

\Needspace{8\baselineskip}
\subsection{La contrainte ouvre un espace d'exploration entre ordre fixe et hasard}
{\small\color{metadata} \texttt{idea\_0103}\quad\textbullet\quad argument\quad\textbullet\quad conceptuelle\par}

L'alternative entre une séquence entièrement fixée et une sélection aléatoire est fausse. Des contraintes globales peuvent garantir une continuité, une répartition ou une forme tout en laissant plusieurs solutions et donc une possibilité réelle d'exploration. La variété n'est pas le hasard : elle devient perceptible sur le fond d'une régularité qui relie les choix.

La qualité de l'espace dépend du compromis entre deux échecs. Trop de contraintes ramènent à un ordre unique; trop peu rendent les variantes arbitraires. Une bonne contrainte ne remplace donc pas le choix : elle lui donne une portée en éliminant les trajectoires qui détruiraient la cohérence globale.

\begin{quote}
\textbf{Intérêt pour la thèse.} Cette carte relie une thèse esthétique sur la variété à un principe technique de génération et d'interaction.
\end{quote}
{\footnotesize\color{metadata}\textbf{Provenance :} Synopsis MIT Press.doc. Synopsis MIT Press, rendu PDF p. 18-23\par}

\Needspace{8\baselineskip}
\subsection{Une surprise musicale forte devient rétrospectivement nécessaire}
{\small\color{metadata} \texttt{idea\_0019}\quad\textbullet\quad hypothèse\quad\textbullet\quad conceptuelle\par}

Dans les chansons dites Jotney, un événement mélodique ou harmonique localement improbable ne brise pas la logique de la chanson : ce qui suit le justifie au point qu'il paraît rétrospectivement nécessaire. Cette surprise diffère donc d'une rupture spectaculaire. Elle déplace l'auditeur, puis lui fait entendre la contrainte invisible qui rendait ce déplacement juste.

L'hypothèse est qu'une œuvre peut concilier accessibilité et invention lorsque sa complexité reste assez invisible pour que la mélodie demeure autonome, mais assez structurante pour que l'analyse augmente, plutôt qu'elle ne dissipe, la singularité.

Le mécanisme dépasse le cas des chansons Jotney. Dans un effet de "aha", un événement ultérieur attribue soudain une fonction à une anomalie antérieure et force l'auditeur ou le spectateur à relire mentalement la séquence. Le futur ne supprime pas l'étrangeté : il la transforme rétrospectivement en indice ou en anticipation.

\begin{quote}
\textbf{Intérêt pour la thèse.} Cette carte fournit un exemple musical concret de l'objet intéressant : une solution qui paraît simple alors qu'elle résout plusieurs contraintes invisibles.
\end{quote}
{\footnotesize\color{metadata}\textbf{Provenance :} The Mystery of Jotney Songs.pdf; Synopsis MIT Press.doc. The Mystery of Jotney Songs, necessary surprise et justification rétrospective, PDF p. 1-3 Synopsis MIT Press, chapitre Aha, rendu PDF p. 31-34\par}

\Needspace{8\baselineskip}
\subsection{Une forme intéressante peut paraître simple parce qu'elle a rendu invisibles les contraintes improbables de sa genèse}
{\small\color{metadata} \texttt{idea\_0087}\quad\textbullet\quad hypothèse\quad\textbullet\quad conceptuelle\par}

Certaines formes semblent \texttt{autostables} : déplacer un élément altérerait l'ensemble, comme si chaque détail était devenu nécessaire aux autres. Cette évidence de surface ne signifie pas que la forme était facile à trouver. Elle peut au contraire être le résultat visible d'une genèse rare, difficile et hautement contrainte dont les essais ont disparu.

L'intérêt naît alors d'une double perception : la forme va de soi une fois entendue, mais l'auditeur pressent qu'elle était improbable avant d'exister. La simplicité du résultat et la difficulté de production ne s'opposent pas ; leur écart constitue une source de singularité.

Cette proposition ouvre une question empirique : quelles traces de la facture permettent effectivement d'inférer une genèse contrainte lorsque les essais et les impasses ont disparu ? Il ne faut pas présupposer que la difficulté réelle laisse toujours une marque perceptible, ni qu'une marque de difficulté prouve l'intérêt.

\begin{quote}
\textbf{Intérêt pour la thèse.} Cette proposition relie rareté, nécessité locale et difficultuosité sans définir l'intéressant par la seule complexité objective.
\end{quote}
{\footnotesize\color{metadata}\textbf{Provenance :} projet thèse philo.pdf; The Mystery of Jotney Songs.pdf; Notes thèse.pdf. PDF p. 3 Jotney, complexité invisible et nécessités cachées sous l'aisance, PDF p. 1-3 Notes thèse, question de la trace de la facture et de la difficultuosité, PDF p. 1\par}

\Needspace{8\baselineskip}
\subsection{Les théories de la créativité sont souvent sans théorèmes}
{\small\color{metadata} \texttt{idea\_0044}\quad\textbullet\quad objection\quad\textbullet\quad conceptuelle\par}

Les discours savants sur la créativité abondent, mais ils expliquent souvent après coup sans permettre de déduire, prédire ou produire un acte créateur. Ils ressemblent davantage à des artefacts explicatifs qu'à des théories fortes.

\begin{quote}
\textbf{Intérêt pour la thèse.} Cette carte permet de critiquer les cadres trop confortables : parler de créativité n'est pas encore la comprendre.
\end{quote}
{\footnotesize\color{metadata}\textbf{Provenance :} De l'impossibilité de créer.pdf. PDF p. 28-35\par}

\section{Musique et micro-transformations de l'attention}

\Needspace{8\baselineskip}
\subsection{La musique est un laboratoire de l'intéressant parce qu'elle rend l'attente et la mémoire constitutives de la forme}
{\small\color{metadata} \texttt{idea\_0086}\quad\textbullet\quad méthode\quad\textbullet\quad articulation\par}

Une forme musicale n'est jamais donnée tout entière : chaque événement est entendu sur le fond de ce qui vient d'être retenu et de ce qui est attendu. Une modulation ou une résolution peut ainsi déplacer les attentes qu'elle accomplit et modifier rétrospectivement le sens du début.

La musique n'est donc pas seulement un domaine d'application commode. Elle expose directement le mécanisme temporel que la thèse attribue à l'intéressant : mémoire, anticipation, surprise, reconnaissance et relance y sont inséparables de l'objet.

\begin{quote}
\textbf{Intérêt pour la thèse.} Cette carte justifie méthodologiquement le privilège accordé aux chansons sans réduire la philosophie de l'intéressant à une esthétique musicale.
\end{quote}
{\footnotesize\color{metadata}\textbf{Provenance :} projet thèse philo.pdf. PDF p. 2-3\par}

\Needspace{8\baselineskip}
\subsection{Une micro-émotion musicale qualifie une transformation attentionnelle, pas un objet harmonique}
{\small\color{metadata} \texttt{idea\_0112}\quad\textbullet\quad hypothèse\quad\textbullet\quad conceptuelle\par}

Les mots inventés dans \emph{Histoire d'une oreille} ne renomment pas simplement des accords déjà répertoriés. Ils isolent des transformations temporelles de l'attention : chute contrôlée vers une stabilité dans l'escatte, résolution vers une nouvelle instabilité dans l'escassepte, frottement d'une solution contre la syntaxe attendue dans la neuveclamine, ou changement de fonction d'une septième de dominante devenue stable dans le stabileme. Une même structure technique ne produit donc pas nécessairement la même micro-émotion, et une micro-émotion peut être réalisée par plusieurs structures.

Nommer l'effet rend une expérience fugitive localisable, comparable et retrouvable. Le concept agit comme un pointeur pour l'écoute : il permet de reconnaître une famille de trajectoires, d'en rechercher des variantes et de décrire ce que le passage fait à l'attention. Mais cette acquisition transforme aussi l'oreille. L'effet d'abord fascinant peut devenir un automatisme, être surexploité, puis retrouver sa force lorsqu'il est préparé avec parcimonie.

Les exemples forment ainsi moins un lexique de recettes qu'une phénoménologie opérationnelle. L'escatte, la misette, la cebimine, le stabileme, la descente triangulaire et le glisseme sont des sondes différentes pour décrire comment une forme locale tend, déplace, suspend ou requalifie l'écoute.

\paragraph{Exemples}
L'escatte est retrouvée dans \emph{For No One} ; l'escassepte dans « Viens je t'emmène » de France Gall ; la misette dans \emph{Only You} et \emph{New York State of Mind} ; la cebimine dans « Tous les bateaux » et \emph{Yesterday} ; le stabileme dans \emph{Oh Happy Day} ; la descente triangulaire dans \emph{Good Day Sunshine}, puis renouvelée par le détour de \emph{Honey Pie} ; le glisseme dans \emph{Oh Lori} et \emph{Just the Two of Us}. La diversité de ces cas importe davantage que l'étiquette harmonique commune.

\begin{quote}
\textbf{Intérêt pour la thèse.} Cette proposition fournit un niveau d'analyse entre le trait musical et le jugement global. Répondre a « pourquoi ce titre est intéressant ? » consiste alors à reconstruire une trajectoire d'attention composée de micro-transformations, sans réduire le titre à la présence d'un accord ou d'un hook.
\end{quote}
{\footnotesize\color{metadata}\textbf{Provenance :} PACHET HISTOIRE OREILLE BAT.pdf. Escatte et escassepte, PDF p. 12-17 ; misette et neuveclamine, p. 21-24 ; soulèvement relatif, p. 47-48 ; cebimine, p. 62-64 ; stabileme, p. 94-95 ; descente triangulaire et glisseme, p. 118-121 ; lexique, p. 311.\par}
{\footnotesize\color{metadata}\textbf{Références :} \parencite{pachet2018oreille}\par}

\Needspace{8\baselineskip}
\subsection{La surprise musicale dépend du degré de réalisation de l'attente}
{\small\color{metadata} \texttt{idea\_0038}\quad\textbullet\quad argument\quad\textbullet\quad scientifique\par}

L'écoute musicale construit en temps réel des attentes plus ou moins fortes et spécifiques. La surprise vient du degré de correspondance entre ce qui était attendu et ce qui advient.

Un même écart objectif ne produit donc pas toujours la même surprise. Une attente forte et précise rend sa non-réalisation saillante ; une attente vague ou ambiguë distribue l'erreur entre plusieurs interprétations possibles. Harmonie, mètre, mémoire et contexte construisent ensemble l'attente, et une surprise peut conduire à réinterpréter rétrospectivement ce qui la précédait.

La régularité n'est pas l'ennemie de la surprise mais sa condition locale : pour déjouer une attente, le compositeur doit d'abord permettre qu'elle se forme. Il doit être simple avant d'être complexe et régulier avant d'être irrégulier.

L'escatte et l'escassepte montrent que la résolution elle-même doit être qualifiée. Une tension peut aboutir à une stabilité qui absorbe rétrospectivement le mouvement, ou à une nouvelle instabilité qui renvoie l'attention vers l'attente non close. L'expertise déplace en outre ces effets : elle rend les surprises plus rares, mais peut ouvrir l'écoute à des structures auparavant inaccessibles.

\paragraph{Contexte bibliographique}
Le modèle s'inscrit dans une filiation plus ancienne. Meyer relie l'affect et le sens musical aux attentes apprises, ainsi qu'à leur inhibition, leur retard ou leur résolution. Narmour prolonge cette intuition avec le modèle implication-réalisation : certains intervalles mélodiques créent des implications dont la continuation peut constituer une réalisation ou un écart, en combinant des principes perceptifs ascendants et des schémas stylistiques appris.

Ces trois ouvrages sont ici des références structurantes, pas encore des sources locales ingérées. Leurs mécanismes précis devront être vérifiés par lecture directe avant d'être mobilisés comme preuves détaillées dans le manuscrit.

\begin{quote}
\textbf{Intérêt pour la thèse.} Cette carte donne un modèle précis de l'intéressant temporel : pas seulement surprise, mais écart mesuré entre prédiction et réalisation.
\end{quote}
{\footnotesize\color{metadata}\textbf{Provenance :} Expectationfinal-paper.pdf; aaai99B.pdf; PACHET HISTOIRE OREILLE BAT.pdf. Berger et Gang, manuscrit intégral, modèle DRE et etudes Haydn, PDF p. 1-34 et figures p. 39-48 Version condensée neurosymbolique, PDF p. 1-6 Histoire d'une oreille, escatte et escassepte comme résolutions distinctes, PDF p. 12-17 ; attentes moins accessibles avec l'expertise, p. 241-242.\par}
{\footnotesize\color{metadata}\textbf{Références :} \parencite{berger1999expectations,gang1999unified,meyer1956emotion,narmour1990basic,narmour1992complexity,pachet2018oreille}\par}

\Needspace{8\baselineskip}
\subsection{L'intéressant d'une séquence naît de dimensions qui se contraignent mutuellement}
{\small\color{metadata} \texttt{idea\_0039}\quad\textbullet\quad méthode\quad\textbullet\quad scientifique\par}

Une longue séquence intéressante n'est jamais seulement mélodique, harmonique, rythmique ou sonore. Son intérêt peut apparaître dans l'accord ou le conflit entre ces dimensions : un timbre rend une harmonie singulière, une orchestration donne une fonction à une mélodie, ou un réseau de personnages révèle une structure que l'intrigue seule ne montre pas. Modéliser chaque dimension indépendamment fait disparaître ces relations transversales.

Cette pluralité concerne aussi les représentations. Les attentes ne sont ni entièrement déductibles de règles, ni réductibles à des fréquences apprises. Un modèle peut faire négocier règles symboliques, concepts graduels, signal audio et apprentissage statistique; l'échec d'une représentation devient alors une information exploitable par les autres plutôt qu'une erreur terminale.

Dans NetNeg, cette coopération est opérationnelle : le réseau propose des notes selon les régularités apprises, puis des agents négocient entre ces préférences et les règles du contrepoint. Les essais séparés échouent de façons complémentaires : le réseau produit des continuations plausibles mais illégales, les règles seules des solutions légales mais pauvres. La qualité vient de leur compromis répété.

Le cas Jotney rend cette contrainte mutuelle particulièrement nette : la mobilité harmonique n'est intéressante que si la mélodie conserve sa cohérence, tandis que l'autonomie mélodique devient perceptible parce qu'elle traverse les changements d'harmonie. Mesurer une dimension sans l'autre manquerait donc précisément l'équilibre que l'analyse cherche à expliquer.

\begin{quote}
\textbf{Intérêt pour la thèse.} Cette carte formule une hypothèse sur l'intéressant et fournit une architecture possible pour la tester sans réduire l'objet à une seule dimension.
\end{quote}
{\footnotesize\color{metadata}\textbf{Provenance :} aaai99A.pdf; aaai99B.pdf; netneg.pdf; Paper citations updated.docx; The Mystery of Jotney Songs -full.pdf. NetNeg condense, apprentissage neuronal et négociation de contraintes, PDF p. 1-6 Modèle unifie attente/contexte/mètre, PDF p. 1-6 NetNeg développe, architecture, expériences et limites, PDF p. 1-18 Projet SenseOfDirection, idée 3 sur la multidimensionnalite, rendu PDF p. 6 Dossier Jotney, profil analytique et équilibre mélodie-harmonie, PDF p. 19-23 et 29-36\par}
{\footnotesize\color{metadata}\textbf{Références :} \parencite{gang1999unified,goldman1999netneg}\par}

\Needspace{8\baselineskip}
\subsection{Autonomie mélodique et mobilité harmonique peuvent se renforcer au lieu de s'opposer}
{\small\color{metadata} \texttt{idea\_0111}\quad\textbullet\quad hypothèse\quad\textbullet\quad conceptuelle\par}

Une chanson peut conserver une mélodie autonome - chantable sans arrangement, mémorable et cohérente - tout en faisant voyager constamment son harmonie. La mobilité harmonique n'affaiblit pas nécessairement l'identité mélodique : elle peut suivre ses inflexions, déplacer silencieusement son centre tonal et donner une nécessité à des notes qui sembleraient autrement ordinaires ou improbables.

L'équilibre est étroit. Trop de mouvements harmoniques rendent la construction démonstrative ou dense ; trop peu la rabattent sur la monotonie d'une boucle. Une modulation trop spectaculaire attire l'attention sur le procédé ; une harmonie trop statique ne soutient plus la trajectoire. Dans le régime nommé Jotney, le travail harmonique reste assez mobile pour porter l'invention et assez discret pour que la chanson paraisse simple.

Jotney doit donc être compris comme un régime compositionnel graduel, non comme un genre ou une étiquette géographique : une chanson peut être plus ou moins Jotney selon la manière dont ces deux dimensions se soutiennent. Les traits proposés (basse chromatique, dominantes secondaires, pivots discrets, mélodie stable sous reharmonisation) sont des indices pour enquêter sur cet équilibre, pas sa définition par checklist.

Cette proposition ne définit pas une supériorité universelle. D'autres traditions résolvent la tension compositionnelle par le rythme, le timbre, la modalité ou la dramaturgie mélodique. Jotney nomme une solution particulière au problème de réconcilier accessibilité populaire et invention structurelle.

\begin{quote}
\textbf{Intérêt pour la thèse.} La carte fournit une structure musicale plus précise que la formule générale du juste milieu entre ordre et surprise. L'intéressant peut naître de deux autonomies qui semblent incompatibles mais se contraignent mutuellement sans exhiber leur travail.
\end{quote}
{\footnotesize\color{metadata}\textbf{Provenance :} The Mystery of Jotney Songs.pdf; The Mystery of Jotney Songs -full.pdf. Jotney, définition de l'équilibre et contrastes stylistiques, PDF p. 1-3 Dossier Jotney, essais, profil analytique et programme de recherche, PDF p. 2-4, 19-23 et 29-36\par}

\Needspace{8\baselineskip}
\subsection{L'optale éprouve une nécessité par variation, la pseudoptale la simule par familiarité}
{\small\color{metadata} \texttt{idea\_0113}\quad\textbullet\quad distinction\quad\textbullet\quad conceptuelle\par}

Une phrase produit un sentiment d'optale lorsque l'auditeur tente mentalement des variations proches et constate qu'elles l'affaiblissent : l'original apparaît comme un optimum local, au point que l'améliorer exigerait de changer le problème entier. Cette nécessité n'est pas une propriété mystique de l'œuvre. Elle est le résultat sensible d'une exploration contrefactuelle de son voisinage.

La pseudoptale constitue un faux positif essentiel. La répétition peut creuser un chemin si familier que toute autre continuation paraît fautive, alors que des alternatives viables existent. L'impression que rien ne peut être changé confond alors deux causes opposées : la solidarité effective des parties ou la simple inertie de la mémoire.

Cette distinction rend l'analyse falsifiable au moins en principe. Il ne suffit pas de déclarer une forme inévitable après l'avoir beaucoup entendue ; il faut produire des variantes suffisamment proches, montrer ce qu'elles détruisent et rester attentif au fait que l'exercice lui-même modifie l'oreille qui juge.

\paragraph{Exemples}
L'ouverture harmonique trouvée par Claude Barthélemy pour une chanson à boire est décrite comme une optale : commencer hors du premier degré résout un problème que les solutions voisines laissaient plat. à l'inverse, le livre avertit que les mélodies apprises par répétition peuvent paraître intouchables seulement parce que leur trace mnémonique est devenue la route normale.

\begin{quote}
\textbf{Intérêt pour la thèse.} L'optale donne une opération concrète pour relier surprise et nécessité rétrospective. La pseudoptale empêche cette opération de transformer toute préférence acquise en preuve de perfection.
\end{quote}
{\footnotesize\color{metadata}\textbf{Provenance :} PACHET HISTOIRE OREILLE BAT.pdf. Écoute superposée, optale et pseudoptale, PDF p. 17-21 ; sentiment d'optale, p. 64-66 ; accord d'ouverture trouvé comme optale, p. 282-283.\par}
{\footnotesize\color{metadata}\textbf{Références :} \parencite{pachet2018oreille}\par}

\Needspace{8\baselineskip}
\subsection{Une phrase banale peut devenir intéressante après coup ou rester un simple substrat}
{\small\color{metadata} \texttt{idea\_0114}\quad\textbullet\quad distinction\quad\textbullet\quad conceptuelle\par}

Deux phrases initialement banales peuvent avoir des fonctions temporelles opposées. Une fausse évidence devient intéressante lorsqu'une reprise ou une transformation ultérieure révèle que sa simplicité préparait une contrainte : le présent requalifie alors rétrospectivement la première écoute. Un substrat nécessaire reste au contraire pauvre en lui-même ; il fournit seulement le temps, le contraste ou l'appui dont un événement futur a besoin.

La répétition ne vaut donc ni comme bien ni comme défaut général. Dans le « matraquage transformatif », elle rend perceptible une organisation qui ne pouvait pas l'être à la première exposition et donne à l'auditeur le sentiment d'avoir résolu un problème. Dans le substrat, elle peut seulement prolonger une préparation. La différence ne se voit qu'en suivant ce que la suite fait au matériau initial.

La « fausse naïveté » pousse cette stratégie près de l'échec : une œuvre laisse l'attention s'approcher de la rupture avant qu'une trouvaille tardive justifie la patience demandée. Ce risque explique pourquoi l'ennui local peut participer à l'intéressant sans que tout ennui devienne rétrospectivement vertueux.

\paragraph{Exemple}
Les mélodies de Nino Rota pour \emph{Amarcord} paraissent d'abord suivre des rails connus. Leurs répétitions légèrement transformées font ensuite entendre la fonction de phrases auparavant creuses, comme un coup de jeu dont la pertinence n'apparaît qu'après la chute d'un groupe de pièces. \emph{Lovely Rita}, pourtant connue par coeur, sert de contre-exemple : la répétition n'y révèle pas nécessairement une organisation cachée.

\begin{quote}
\textbf{Intérêt pour la thèse.} Cette distinction donne une grammaire plus précise de l'attente. Elle permet d'analyser pourquoi un titre apparemment simple gagne à la réécoute, tout en conservant la possibilité qu'un passage soit effectivement vide.
\end{quote}
{\footnotesize\color{metadata}\textbf{Provenance :} PACHET HISTOIRE OREILLE BAT.pdf. Fausse naïveté et interdifference, PDF p. 133-135 ; réalisation progressive, fausses évidences, matraquage transformatif et substrat nécessaire, p. 218-220.\par}
{\footnotesize\color{metadata}\textbf{Références :} \parencite{pachet2018oreille}\par}

\Needspace{8\baselineskip}
\subsection{Les hooks harmoniques comme micro-solutions intéressantes}
{\small\color{metadata} \texttt{idea\_0034}\quad\textbullet\quad exemple\quad\textbullet\quad articulation\par}

Un hook harmonique peut être compris comme une petite solution locale : une progression attire l'oreille par un déplacement, une couleur ou une tension, tout en restant suffisamment chantable et memorisable. Le document rapproche trois enchaînements d'accords de plusieurs chansons ; leur récurrence suggère qu'une surprise harmonique peut être réutilisable sans perdre son effet d'accroche.

La discussion associée suggère un mécanisme de production plus précis : les harmonies intéressantes ne seraient pas choisies comme des formules autonomes, mais émergeraient de la mélodie, de l'adéquation au contenu et d'un refus des solutions trop évidentes. Cette formulation reste une hypothèse. La liste et les témoignages ne montrent ni que ces progressions causent l'intérêt, ni qu'elles sont plus fréquentes dans les succès que dans un corpus témoin. Il faudrait une analyse harmonique comparative et des tests d'écoute pour transformer l'intuition en résultat.

Un éclat local ne doit pas non plus être confondu avec la cause autonome de l'intérêt du titre. Dans \emph{Ce n'est rien}, l'accord saillant n'est remarquable que parce que la chanson entière lui fournit rythme, attente et contraste. Dans Take 6, l'accumulation d'éclats peut au contraire saturer l'attention. L'analyse doit donc tester les deux directions causales : le détail organise-t-il l'ensemble, ou l'ensemble rend-il seulement ce détail perceptible et intéressant ?

\begin{quote}
\textbf{Intérêt pour la thèse.} Cette carte donne un niveau d'analyse plus fin que la chanson entière : l'intéressant peut se loger dans une micro-décision harmonique.
\end{quote}
{\footnotesize\color{metadata}\textbf{Provenance :} Hooks in hits.doc; mvdig003.htm; PACHET HISTOIRE OREILLE BAT.pdf. Hooks in hits.doc, p. 1 : trois progressions rapprochées de chansons populaires ; liste exploratoire sans analyse ni protocole comparatif. mvdig003.htm, messages MV88, MV92 et réponse du compositeur : analyses d'accords et témoignage selon lequel les harmonies dérivent de la mélodie et du refus de solutions trop évidentes ; matériau d'exemple, non preuve générale. Histoire d'une oreille, saturation d'effets chez Take 6, PDF p. 195-197 ; effets et structure dans Sign o' the Times, p. 246-247 ; éclat rendu possible par la chanson entière dans Ce n'est rien, p. 270-271.\par}
{\footnotesize\color{metadata}\textbf{Références :} \parencite{pachet2018oreille}\par}

\Needspace{8\baselineskip}
\subsection{Quand la musique devient liquide, l'attention devient rare}
{\small\color{metadata} \texttt{idea\_0040}\quad\textbullet\quad argument\quad\textbullet\quad conceptuelle\par}

Quand les copies musicales deviennent parfaites, gratuites et manipulables, la valeur se déplace vers ce qui ne se copie pas facilement : confiance, immédiateté, personnalisation, sélection, contexte et attention. La technique ne modifie donc pas seulement la distribution de la musique : comme l'enregistrement avait transformé la manière de jouer, la fluidité numérique transforme l'objet musical en pratique de sélection, de recombinaison et de participation.

Dans ce régime d'abondance, découvrir la bonne musique et construire un parcours d'écoute peuvent avoir davantage de valeur que posséder une copie. L'intéressant devient alors une relation organisée entre un auditeur, un contexte et une profusion d'objets disponibles.

Une discussion d'auditeurs donne un indice complémentaire : le bouton permettant de passer instantanément un morceau favorise le jugement initial, tandis qu'un support imposant une exposition plus longue peut laisser un morceau d'abord rejeté devenir favori. Ce témoignage ne démontre pas l'effet, mais formule une hypothèse testable sur la manière dont l'interface redistribue la patience nécessaire à l'intéressant.

\begin{quote}
\textbf{Intérêt pour la thèse.} Cette carte relie l'intéressant à une économie de l'abondance : quand les objets prolifèrent, le problème devient de produire des prises attentionnelles et des parcours singuliers.
\end{quote}
{\footnotesize\color{metadata}\textbf{Provenance :} WhereMusicWillBeComingFromNYTArticle.txt; mvdig003.htm. Kevin Kelly, Where Music Will Be Coming From, New York Times Magazine, 2002 : article intégral ; copie numérique, musique manipulable et déplacement de la valeur vers les qualités non copiables. mvdig003.htm, message sur l'écoute CD/cassette : témoignage selon lequel la possibilité de sauter immédiatement un morceau peut empêcher l'attachement tardif ; indice anecdotique d'un effet de l'interface sur l'attention.\par}
{\footnotesize\color{metadata}\textbf{Références :} \parencite{kelly2002music}\par}

\Needspace{8\baselineskip}
\subsection{L'écoute surattentive peut fabriquer l'intérêt qu'elle prétend découvrir}
{\small\color{metadata} \texttt{idea\_0115}\quad\textbullet\quad objection\quad\textbullet\quad conceptuelle\par}

L'attention n'est pas seulement un détecteur passif de qualités. Lorsqu'un auditeur admire déjà un artiste, il peut traiter chaque faiblesse comme un code, inventer la résolution d'une banalité et produire lui-même le plaisir qu'il attribue ensuite au morceau. La confiance, la contagion d'un album ou le désir de retrouver un effet ancien deviennent une assistance interpretative.

Ce pouvoir actif rend l'enquête réflexive mais aussi circulaire : plus on cherche la singularité, plus on donne de l'importance au passage examiné. Inversement, l'objet peut s'auto-évanouir lorsque sa reconstruction réussit, ou perdre son attrait lorsque l'énigme est effectivement résolue. L'analyse transforme donc la valeur qu'elle cherche à expliquer dans les deux directions.

La surattention n'invalide pas toute analyse. Elle impose de comparer les interprétations, d'utiliser des variantes et des contre-exemples, de distinguer une surprise qui s'impose sans assistance d'une résolution laborieusement fabriquée, et de documenter l'histoire de l'écoute au lieu de présenter le jugement final comme une propriété immédiate de l'objet.

\paragraph{Exemples}
Les albums tardifs de Take 6 et \emph{Free as a Bird} montrent une écoute qui pardonne et complète parce qu'elle veut retrouver une magie ancienne. \emph{Surf's Up} sert de cas contrasté : la surprise y est décrite comme s'imposant avant le travail de justification. La version complète de \emph{Her Majesty} fournit l'autre limite de l'enquête : résoudre l'énigme peut supprimer l'intérêt qui motivait sa recherche.

\begin{quote}
\textbf{Intérêt pour la thèse.} Cette objection protège la théorie des micro-émotions et de l'optale contre l'auto-confirmation. Elle fait aussi de la transformation de l'observateur par sa propre méthode un problème central de la thèse.
\end{quote}
{\footnotesize\color{metadata}\textbf{Provenance :} PACHET HISTOIRE OREILLE BAT.pdf. Auto-évanescence de l'objet analyse, PDF p. 191 ; contagion, p. 220 ; écoute surattentive et résolution forcée des fausses évidences, p. 221-222 ; énigme résolue qui détruit l'intérêt, p. 284-285.\par}
{\footnotesize\color{metadata}\textbf{Références :} \parencite{pachet2018oreille}\par}

\section{IA, interaction et méthode réflexive}

\Needspace{8\baselineskip}
\subsection{L'IA comme laboratoire philosophique de la naissance des idées}
{\small\color{metadata} \texttt{idea\_0018}\quad\textbullet\quad méthode\quad\textbullet\quad articulation\par}

Construire un système d'IA oblige à rendre explicites des opérations que le discours philosophique peut laisser indéterminées : comment produire une variante, la sélectionner, reconnaître une forme ou constater un échec. La construction ne démontre pas à elle seule une théorie de l'esprit, mais elle transforme une notion en dispositif dont les présupposés et les limites deviennent observables.

L'IA est donc ici moins un objet auquel appliquer la philosophie qu'une pratique expérimentale capable de faire surgir de nouvelles questions philosophiques sur la naissance et la sélection des idées.

Cette fonction de laboratoire ne transforme pas automatiquement un résultat scientifique en proposition philosophique. Elle fournit des opérations, des échecs et des contrastes qu'une enquête réflexive doit encore reconstruire et argumenter.

\begin{quote}
\textbf{Intérêt pour la thèse.} Cette carte justifie la légitimité d'une thèse philosophique appuyée sur une pratique de recherche en IA.
\end{quote}
{\footnotesize\color{metadata}\textbf{Provenance :} projet thèse philo.pdf. PDF p. 1 et 3-4\par}

\Needspace{8\baselineskip}
\subsection{Comprendre une pratique scientifique exige de distinguer ses résultats des idées qu'elle rend pensables}
{\small\color{metadata} \texttt{idea\_0110}\quad\textbullet\quad méthode\quad\textbullet\quad articulation\par}

Une activité scientifique produit des systèmes, des algorithmes, des protocoles et des résultats dont les critères de validité appartiennent à leur discipline. Ces productions ne constituent pas directement une philosophie de l'intéressant. Elles peuvent toutefois rendre visibles des opérations récurrentes - générer, contraindre, sélectionner, reconnaître, revenir en arrière, terminer - ainsi que les présupposés qui orientaient leur construction sans avoir été formulés comme tels.

"Comprendre ce que je fais" consiste alors en une enquête de second niveau. Il ne s'agit ni de résumer les publications, ni de convertir rétrospectivement chaque système en preuve philosophique. Il s'agit de reconstruire les questions que la pratique a rencontrées, de distinguer les solutions techniques des intuitions conceptuelles, puis d'examiner quels passages entre elles sont légitimes.

Faire ces passages de manière explicite peut constituer une contribution propre de la thèse. Le résultat n'est pas seulement une collection d'idées sur l'intéressant, mais une méthode pour transformer une trajectoire de construction scientifique en matériau philosophique sans confondre les deux régimes de preuve.

\begin{quote}
\textbf{Intérêt pour la thèse.} Cette proposition donne une fonction architecturale aux cartes \texttt{scientific} et \texttt{articulation}. Les premières conservent ce qui a effectivement été construit ou établi ; les secondes formulent et soumettent à la critique l'interprétation philosophique de cette activité.
\end{quote}
{\footnotesize\color{metadata}\textbf{Provenance :} projet thèse philo.pdf. Projet de thèse, comprendre quarante ans de construction de systèmes, PDF p. 1 Projet de thèse, reprise philosophique des anciens travaux computationnels, PDF p. 3-4\par}

\Needspace{8\baselineskip}
\subsection{Le système intéressant agit comme miroir interactif}
{\small\color{metadata} \texttt{idea\_0006}\quad\textbullet\quad argument\quad\textbullet\quad articulation\par}

Un système réflexif apprend en cours d'interaction des régularités du geste de l'utilisateur et les lui renvoie sous une forme transformée. Il ne produit pas une œuvre intéressante sur commande et n'expose pas seulement davantage d'options : chaque action humaine déclenche une imitation à laquelle l'utilisateur peut à son tour réagir. L'intérêt vient de la reconnaissance de sa propre activité, déplacée assez fortement pour appeler une nouvelle action.

Cette boucle diminue aussi le besoin d'un modèle semantique complet. L'ambiguïté et la dépendance au contexte peuvent être corrigées localement par les réactions de l'utilisateur : la signification utile au système se construit dans les reprises et les rectifications plutôt que d'être entièrement spécifiée avant l'interaction.

Le critère de succès n'est pas l'effet bref d'une démonstration. à chaque tour du Continuator, l'utilisateur fournit davantage de matériau, le modèle se précise et la réponse suivante modifie ce que l'utilisateur joue. L'interaction devient intéressante si cette évolution soutient l'attention sur plusieurs cycles, au lieu de répéter indéfiniment le même tour de passe-passe.

\begin{quote}
\textbf{Intérêt pour la thèse.} Le Continuator devient un cas exemplaire : l'intéressant n'est ni dans l'utilisateur seul, ni dans le système seul, mais dans la boucle.
\end{quote}
{\footnotesize\color{metadata}\textbf{Provenance :} Interesting Interactions (sent to Luc).doc; ERCInteractiveReflexions2.docx; Synopsis MIT Press.doc; ERCGrantPachetInterestingness.pdf. Grant ERC PDF p. 4 Synopsis MIT Press, rendu PDF p. 4 et 36 Projet REFLEX, rendu PDF p. 1 et 4-6 ; argument sur la boucle p. 17 Interesting Interactions, rendu PDF p. 7\par}

\Needspace{8\baselineskip}
\subsection{La création comme effet secondaire de l'interaction}
{\small\color{metadata} \texttt{idea\_0007}\quad\textbullet\quad argument\quad\textbullet\quad articulation\par}

Dans les systèmes réflexifs, le contenu créé n'est pas seulement le résultat d'une intention préalable. Il émerge comme effet secondaire d'une interaction soutenue avec un système qui imite, prolonge, contrarie ou relance. L'objet se constitue donc à travers une série de réponses mutuelles : il n'est ni généré au toucher d'un bouton, ni entièrement contenu dans ce que l'utilisateur avait en tête.

\begin{quote}
\textbf{Intérêt pour la thèse.} Cette carte déplace la création de l'acte souverain vers une écologie d'interactions.
\end{quote}
{\footnotesize\color{metadata}\textbf{Provenance :} Interesting Interactions (sent to Luc).doc; ERCInteractiveReflexions2.docx; TBKLullyNOTES.doc. Projet REFLEX, résumé et description du projet, rendu PDF p. 1 et 4-6 Interesting Interactions, rendu PDF p. 7 TBKLullyNOTES, couple utilisateur-Continuator et plaisir de faire, rendu PDF p. 6-8\par}

\Needspace{8\baselineskip}
\subsection{L'influence sociale n'abolit pas les propriétés propres qui orientent la popularité}
{\small\color{metadata} \texttt{idea\_0013}\quad\textbullet\quad hypothèse\quad\textbullet\quad scientifique\par}

La pression sociale rend la popularité contingente, mais elle ne produit pas une distribution uniforme des succès. Certains objets attirent plus souvent certains publics, ce qui suggère que des propriétés de l'objet interagissent avec les sujets et les contextes. Le \texttt{popularity feature problem} consiste à identifier ces propriétés sans les confondre avec une essence universelle de la valeur.

\begin{quote}
\textbf{Intérêt pour la thèse.} Cette question évite deux réductions : tout expliquer par les propriétés internes, ou tout expliquer par le social.
\end{quote}
{\footnotesize\color{metadata}\textbf{Provenance :} ERCGrantPachetInterestingness.pdf; ERCInteractiveReflexions2.docx. Grant ERC PDF p. 1 Projet REFLEX, discussion de Salganik, rendu PDF p. 4\par}

\Needspace{8\baselineskip}
\subsection{PathBuilder : l'intéressant entre deux choses aimées}
{\small\color{metadata} \texttt{idea\_0014}\quad\textbullet\quad hypothèse\quad\textbullet\quad scientifique\par}

Si quelqu'un aime A et B, il devrait exister des chemins intéressants entre A et B, à condition de disposer d'une relation de similarité assez riche. L'intérêt ne vient pas seulement des points, mais du parcours entre eux. Un tel chemin doit combiner des contraintes locales de continuité et des propriétés globales de la séquence; la similarité seule ne suffit pas à garantir une progression.

\begin{quote}
\textbf{Intérêt pour la thèse.} Cette carte transforme la recommandation en problème de navigation : faire émerger des transitions intéressantes plutôt que prédire un objet préféré.
\end{quote}
{\footnotesize\color{metadata}\textbf{Provenance :} TBKLullyNOTES.doc; Synopsis MIT Press.doc. TBKLullyNOTES, hypothèse PathBuilder, rendu PDF p. 1 Synopsis MIT Press, rendu PDF p. 18-21\par}

\Needspace{8\baselineskip}
\subsection{Democratic composition : enregistrer tout le processus de création}
{\small\color{metadata} \texttt{idea\_0068}\quad\textbullet\quad méthode\quad\textbullet\quad scientifique\par}

Pour comprendre la création, il faut enregistrer les essais, erreurs, abandons et choix mélodiques, harmoniques et timbraux d'une composition complète, puis suivre l'œuvre jusqu'à sa production et sa réception. L'objet final efface la topologie du parcours : les bifurcations, retours et solutions refusées qui rendent la décision intelligible.

La trace doit inclure le contexte et le moment des actions, pas seulement leurs résultats. Elle permet alors d'étudier la création comme chemin singulier plutôt que de reconstruire après coup une intention linéaire.

Cette distinction interdit aussi de confondre deux explications légitimes : une analyse peut montrer comment l'œuvre finie agit sur un auditeur sans décrire le processus qui l'a produite. Inversement, le souvenir de l'auteur ne constitue pas à lui seul une trace fiable de ce processus ; il doit être confronté aux actions et bifurcations enregistrées au moment de la création.

\begin{quote}
\textbf{Intérêt pour la thèse.} Cette carte donne une méthode empirique pour rendre visible ce que la preuve ou l'œuvre achevée effacent.
\end{quote}
{\footnotesize\color{metadata}\textbf{Provenance :} ERCGrantPachetInterestingness.pdf; Synopsis MIT Press.doc; ERCInteractiveReflexions2.docx; mvdig003.htm. Grant ERC PDF p. 4-5 Synopsis MIT Press, prologue, rendu PDF p. 3 Projet REFLEX, protocole des etudes pilotes, rendu PDF p. 6-8 et 17 mvdig003.htm, message MV94 : le compositeur distingue l'analyse du produit fini de celle du processus d'écriture et juge le récit rétrospectif de l'artiste peu fiable.\par}

\Needspace{8\baselineskip}
\subsection{La popularité musicale exige des traits générés, pas seulement mesurés}
{\small\color{metadata} \texttt{idea\_0069}\quad\textbullet\quad méthode\quad\textbullet\quad scientifique\par}

Les traits audio génériques ne suffisent pas à expliquer la popularité ou la qualité d'une harmonisation. Un apprentissage de descripteurs doit pouvoir inventer des propriétés adaptées à une tâche supervisée, notamment des relations mélodico-harmoniques et des structures longues.

La réduction usuelle d'un morceau à un sac de trames détruit précisément ces relations temporelles. Générer de bons traits exige donc de représenter la séquence entière au lieu d'agréger des mesures locales avant l'apprentissage.

\begin{quote}
\textbf{Intérêt pour la thèse.} Cette carte relie l'intéressant à la construction active des descripteurs.
\end{quote}
{\footnotesize\color{metadata}\textbf{Provenance :} ERCGrantPachetInterestingness.pdf. Grant ERC PDF p. 5-6\par}

\Needspace{8\baselineskip}
\subsection{On ne peut comprendre l'intéressant sans étudier ensemble la perception et la création}
{\small\color{metadata} \texttt{idea\_0094}\quad\textbullet\quad méthode\quad\textbullet\quad articulation\par}

Étudier seulement les objets préférés décrit ce qui a retenu l'attention sans montrer comment une forme a été trouvée. Étudier seulement la production ignore les jugements qui orientent, interrompent ou relancent chaque tentative. Les deux processus forment une boucle : le créateur produit selon certaines possibilités, puis perçoit son propre résultat avec des critères qui peuvent différer de ceux de son action.

Une théorie active de l'intéressant doit donc observer conjointement les artefacts, les choix successifs et les réactions des sujets, dans des situations naturelles comme dans des interactions artificielles contrôlées.

\begin{quote}
\textbf{Intérêt pour la thèse.} Cette proposition donne un principe d'organisation aux expériences et interdit de réduire l'intéressant à un classement statique d'objets finis.
\end{quote}
{\footnotesize\color{metadata}\textbf{Provenance :} ERCGrantPachetInterestingness.pdf; ERCInteractiveReflexions2.docx. PDF p. 2 et 4 Projet REFLEX, résumé et hypothèses, rendu PDF p. 1 et 5 ; protocole p. 6-8\par}

\Needspace{8\baselineskip}
\subsection{Un contrôle devient significatif lorsqu'il préserve des invariants sémantiques}
{\small\color{metadata} \texttt{idea\_0102}\quad\textbullet\quad argument\quad\textbullet\quad articulation\par}

Les contrôles d'un dispositif musical héritent souvent des possibilités de la technique : volume, position, vitesse ou paramètres atomiques. Ils deviennent musicalement significatifs seulement si le système connaît des relations entre les parties et propage l'action de l'utilisateur de manière à préserver certaines intentions globales.

Dans MusicSpace, déplacer un instrument peut ainsi entraîner le déplacement d'un autre pour conserver un groupe ou un équilibre. L'invariant ne dicte pas une configuration unique; il définit une famille de variations dans laquelle l'action reste cohérente. Le contrôle exprime alors une intention plutôt qu'une manipulation de bouton.

\begin{quote}
\textbf{Intérêt pour la thèse.} Cette proposition donne un critère de conception concret pour une participation qui transforme l'œuvre sans rendre toute action équivalente.
\end{quote}
{\footnotesize\color{metadata}\textbf{Provenance :} Synopsis MIT Press.doc. Synopsis MIT Press, chapitre Active Listening, rendu PDF p. 11-16 et 20\par}

\Needspace{8\baselineskip}
\subsection{Une interaction intéressante maintient à la fois l'autonomie et l'ajustement mutuel}
{\small\color{metadata} \texttt{idea\_0104}\quad\textbullet\quad hypothèse\quad\textbullet\quad articulation\par}

Un partenaire artificiel purement autonome cesse de répondre, tandis qu'un partenaire entièrement soumis ne fait que répéter l'utilisateur. L'interaction devient intéressante lorsqu'une production possède sa dynamique propre tout en satisfaisant des contraintes issues de l'autre participant.

Le texte rapproche ce double ajustement du Continuator et de l'écriture contrainte : un enchaînement peut être probable selon sa dynamique locale et, simultanément, convenir à une contrainte externe. L'intérêt réside dans cette compatibilité rare, qui permet de reconnaître une réponse sans pouvoir la réduire à un écho.

Le paramètre d'attachement proposé pour le Continuator rend ce compromis opérationnel : il règle à quel point la réponse reste liée au matériau de l'utilisateur. L'autonomie et l'interactivité ne sont donc pas deux propriétés à maximiser séparément, mais les extrémités d'un réglage dont la zone productive peut évoluer au cours du dialogue.

\begin{quote}
\textbf{Intérêt pour la thèse.} Cette hypothèse rend testable le vague idéal d'une machine "partenaire" en faisant varier indépendamment autonomie du modèle et force de l'ajustement à l'utilisateur.
\end{quote}
{\footnotesize\color{metadata}\textbf{Provenance :} Synopsis MIT Press.doc; Interesting Interactions (sent to Luc).doc; TBKLullyNOTES.doc. Synopsis MIT Press, chapitre Interesting Interactions, rendu PDF p. 36 Interesting Interactions, paramètre d'attachement, rendu PDF p. 7 TBKLullyNOTES, autonomie et tâche réalisée par le couple, rendu PDF p. 2 et 6\par}

\Needspace{8\baselineskip}
\subsection{Naturaliser l'intéressant, c'est expliquer une relation psychique sans l'assigner à l'objet}
{\small\color{metadata} \texttt{idea\_0125}\quad\textbullet\quad distinction\quad\textbullet\quad conceptuelle\par}

Étudier l'intéressant revient à constituer comme objet d'enquête un phénomène central de la vie psychique qui demeure souvent abandonné à l'évidence privée : quelque chose nous arrête, mobilise notre attention, nous fait revenir, construire et parfois nous attache. Cette objectivation ne consiste pas à déclarer que l'intérêt est une propriété mesurable déposée dans l'objet. Elle consiste à rechercher les causes, conditions, transformations et régularités d'une relation entre une forme et un sujet.

La préface de la partie III de l'\emph{Éthique} fournit un geste méthodologique décisif. Spinoza refuse de traiter les affects humains comme un « empire dans un empire » qui échapperait aux lois communes de la nature. Actions, appétits et affects doivent être compris par leurs causes avec la même exigence que les autres phénomènes naturels. La thèse peut reprendre ce geste : le fait qu'une expérience soit subjective et située ne la rend ni surnaturelle, ni arbitraire, ni réfractaire à l'explication.

À propos de la haine, de la colère et de l'envie, Spinoza écrit qu'« elles ont des causes déterminées qui servent à les expliquer » (\emph{Éthique}, III, préface, traduction Saisset). Le parallèle proposé est fort mais limité : comme la colère, l'intérêt n'arrive pas sans raison ; il doit être rapporté aux causes qui le font naître, durer, se transformer ou disparaître. Cette causalité n'est toutefois pas déposée dans l'objet seul. Spinoza indique ailleurs que les idées des affections enveloppent à la fois « la nature de notre corps et la nature présente du corps extérieur » (\emph{Éthique}, III, proposition XXVII, démonstration). L'affect donne ainsi un précédent pour penser une causalité de la rencontre.

Transposée à l'intéressant, cette causalité est distribuée entre la structure et les prises offertes par l'objet ; la mémoire, les compétences, les attentes et l'état du sujet ; les associations produites par son histoire ; enfin les circonstances sociales et temporelles de la rencontre. Expliquer pourquoi un objet devient intéressant ne consiste donc ni à découvrir en lui une essence de l'intérêt, ni à invoquer une préférence subjective sans cause. Il faut reconstruire la configuration relationnelle qui produit cet effet pour ce sujet à ce moment.

Il faut donc distinguer deux sens de l'objectivation. \textbf{Objectiver le phénomène} est légitime : décrire des trajectoires, isoler des variables, comparer des sujets et des formes, construire des modèles et produire des tests intersubjectifs. \textbf{Objectiver l'intéressant dans l'objet} serait une erreur : aucune propriété intrinsèque de la forme ne suffit à expliquer l'intérêt indépendamment des compétences, de la mémoire, du corps, de l'histoire et de l'horizon social du sujet. Le modèle \texttt{I(F, S | H, t)} permet précisément de naturaliser la relation sans la réifier dans l'un de ses termes.

Le rapprochement avec Spinoza n'impose pas de classer immédiatement l'intéressant parmi les émotions élémentaires. Dans la définition III de la partie III, l'affect engage une modification de la puissance d'agir et l'idée de cette modification. L'intéressant pourrait plutôt désigner un \textbf{régime affectivo-cognitif de transition} : la rencontre avec une forme mobilise l'attention et modifie la puissance de percevoir, de comprendre, d'anticiper ou d'agir. L'ennui, l'anxiété, le plaisir, la frustration ou l'attachement seraient alors des devenirs possibles de cette relation, et non de simples synonymes de l'intérêt.

\paragraph{Conséquence méthodologique}
Naturaliser l'intéressant exige de coordonner plusieurs descriptions sans décider à l'avance que l'une absorbera les autres : expérience à la première personne, comportements observables, transformations des compétences, structure de l'objet, histoire du sujet et modèles computationnels. Une explication réussie doit montrer comment ces niveaux se contraignent, plutôt que remplacer l'expérience par un seul indice physiologique ou statistique.

\begin{quote}
\textbf{Intérêt pour la thèse.} Cette proposition donne une ambition philosophique précise au projet : faire pour un phénomène négligé de la vie psychique un geste analogue à celui que Spinoza accomplit pour les affects, tout en conservant le caractère relationnel, historique et social de l'intéressant. Elle transforme l'opposition vague entre subjectivité et science en un programme causal, comparatif et critique.

Elle justifie également la neutralité axiologique provisoire : expliquer pourquoi une chanson, un jeu ou un modèle de langage nous attache ne revient ni à l'approuver ni à le condamner. L'évaluation éthique devient plus exigeante lorsqu'elle porte sur des mécanismes compris plutôt que sur une exception psychique supposée inexplicable.
\end{quote}
{\footnotesize\color{metadata}\textbf{Provenance :} BUT DE LA THESE.md. Spinoza, Éthique, partie III, préface : les affects ont des causes déterminées qui permettent de les expliquer ; traduction Saisset, 1861. Spinoza, Éthique, partie III, proposition XXVII, démonstration : l'idée de l'affection enveloppe à la fois la nature de notre corps et celle du corps extérieur. Macherey, Introduction à l'Éthique de Spinoza. La troisième partie : la vie affective, 1995.\par}
{\footnotesize\color{metadata}\textbf{Références :} \parencite{spinoza1966ethique,spinoza1861oeuvres,macherey1995vieaffective}\par}

\Needspace{8\baselineskip}
\subsection{Un modèle computationnel peut mettre une hypothèse philosophique à l'épreuve sans épuiser son objet}
{\small\color{metadata} \texttt{idea\_0088}\quad\textbullet\quad méthode\quad\textbullet\quad articulation\par}

Formaliser l'attention, la surprise ou la sélection oblige à choisir des variables, des mécanismes et des critères observables. Le modèle peut alors révéler une contradiction, produire une prédiction ou montrer qu'une intuition reste trop générale. Mais cette explicitation locale ne prouve pas que l'expérience de l'intéressant se réduise au mécanisme implémenté.

Les anciens projets computationnels doivent donc servir d'arrière-plan expérimental et de réserve d'hypothèses, puis être reformulés dans un cadre philosophique qui interroge leurs présupposés au lieu de les prendre pour une théorie complète.

Le cadre peut chercher un compromis entre conditions subjectives du jugement, naturalisation de certaines structures transcendantales et émergence de propriétés globales dans des systèmes dynamiques. Ce compromis est un programme à argumenter, non une réduction déjà acquise : il doit encore distinguer ce qui pourrait être universel de ce qui dépend d'un contexte perceptif comme la musique tonale occidentale.

Schmidhuber reconnaît lui-même que les millions de surprises et de modifications internes de son agent sont difficiles à interpréter, et que les capacités de cet agent différent fortement de celles d'un observateur humain. Une simulation qui produit la dynamique attendue ne suffit donc pas à identifier ce qu'un humain trouvera intéressant.

\begin{quote}
\textbf{Intérêt pour la thèse.} Cette position permet d'utiliser quarante ans de construction de systèmes sans confondre preuve de concept, modèle cognitif et ontologie de l'intéressant.
\end{quote}
{\footnotesize\color{metadata}\textbf{Provenance :} projet thèse philo.pdf; Project philosophy thesis.pdf; interestingness.pdf. PDF p. 3 Version anglaise du projet, PDF p. 1 : articulation proposée entre philosophie kantienne du sujet, naturalisation du transcendantal chez Petitot et propriétés émergentes des systèmes dynamiques. Schmidhuber, What's Interesting?, limites d'interprétation du système, PDF p. 12 et 21\par}
{\footnotesize\color{metadata}\textbf{Références :} \parencite{schmidhuber1997interesting,petitot1993phenomenologie,petitot2002naturaliser}\par}

\section{Virtuosité, difficulté et incorporation}

\Needspace{8\baselineskip}
\subsection{Virtuosité et difficultuosité}
{\small\color{metadata} \texttt{idea\_0010}\quad\textbullet\quad distinction\quad\textbullet\quad conceptuelle\par}

La virtuosité n'est pas seulement une difficulté vaincue. Elle met en scène la difficulté, la rend perceptible, parfois spectaculaire, parfois trompeuse. La difficultuosité serait le régime où la difficulté devient elle-même matériau d'intérêt.

\begin{quote}
\textbf{Intérêt pour la thèse.} Cette carte aide à distinguer effort réel, effort visible, performance, séduction et valeur créative.
\end{quote}
{\footnotesize\color{metadata}\textbf{Provenance :} De l'impossibilité de créer.pdf; ESSAI-La Virtuosité à portée des caniches-F. PACHET.pdf; ERCGrantPachetInterestingness.pdf. De l'impossibilité de créer, PDF p. 46-58 La Virtuosité à portée des caniches, PDF p. 13-16 Grant ERC PDF p. 4 et 6\par}

\Needspace{8\baselineskip}
\subsection{L'appropriation par participation s'oppose à l'appropriation par confrontation}
{\small\color{metadata} \texttt{idea\_0020}\quad\textbullet\quad distinction\quad\textbullet\quad conceptuelle\par}

L'effet IKEA maximise l'appropriation tout en minimisant le risque : une micro-difficulté calibrée suffit pour que le sujet survalorise ce qu'il a activé ou assemble. L'effet Ménard suit le mouvement inverse : refaire à l'identique devient signifiant lorsque le chemin est volontairement coûteux, non optimisé et exposé à l'échec.

L'originalité de l'artefact ne suffit donc pas à distinguer production et création. Il faut considérer la facture et la relation au processus : participation guidée d'un côté, confrontation consciente à une impossibilité de l'autre.

\begin{quote}
\textbf{Intérêt pour la thèse.} Cette carte permet de séparer production matérielle, effort subjectif, nouveauté objective et effet de cadre.
\end{quote}
{\footnotesize\color{metadata}\textbf{Provenance :} De l'impossibilité de créer.pdf. PDF p. 73-79\par}

\Needspace{8\baselineskip}
\subsection{La virtuosité séduit par l'apparence d'une liberté sans friction}
{\small\color{metadata} \texttt{idea\_0048}\quad\textbullet\quad argument\quad\textbullet\quad conceptuelle\par}

La virtuosité attire parce qu'elle donne à voir une liberté qui semble avoir absorbé la difficulté. Ce qui séduit n'est pas seulement le résultat, mais la sensation que tout aurait pu varier sans rupture.

\begin{quote}
\textbf{Intérêt pour la thèse.} Cette carte sert de pont entre l'essai sur la création et \texttt{Caniches}.
\end{quote}
{\footnotesize\color{metadata}\textbf{Provenance :} De l'impossibilité de créer.pdf; ESSAI-La Virtuosité à portée des caniches-F. PACHET.pdf. PDF p. 46-53 La Virtuosité à portée des caniches, PDF p. 9-13\par}

\Needspace{8\baselineskip}
\subsection{La difficultuosité vit dans une zone de crédibilité}
{\small\color{metadata} \texttt{idea\_0049}\quad\textbullet\quad distinction\quad\textbullet\quad conceptuelle\par}

Une performance impressionne quand la difficulté paraît extrême mais encore humaine. Trop facile, elle ennuie ; trop impossible, elle bascule vers l'artefact et cesse de susciter l'admiration.

\begin{quote}
\textbf{Intérêt pour la thèse.} Cette carte aide à penser pourquoi l'IA ou le trucage peuvent produire moins d'intérêt malgré une complexité plus grande.
\end{quote}
{\footnotesize\color{metadata}\textbf{Provenance :} De l'impossibilité de créer.pdf; ESSAI-La Virtuosité à portée des caniches-F. PACHET.pdf. De l'impossibilité de créer, PDF p. 46 et 54-58 La Virtuosité à portée des caniches, PDF p. 13-15\par}

\Needspace{8\baselineskip}
\subsection{Le flow est le signal tardif d'une difficulté incorporée, pas un moteur créatif}
{\small\color{metadata} \texttt{idea\_0050}\quad\textbullet\quad argument\quad\textbullet\quad conceptuelle\par}

Le flow apparaît lorsque les compétences acquises sont ajustées à une difficulté réelle : trop facile, la tâche ennuie ; trop difficile, elle angoisse. Il ne génère donc pas la création et n'explique pas la formation des compétences. Il signale subjectivement que le travail antérieur a rendu la difficulté traversable.

À compétences fixes, le diagramme compétences-exigences ordonne donc trois régions lorsque la difficulté augmente : \textbf{ennui → intéressant/flow → anxiété}. La région intermédiaire correspond à une tâche assez difficile pour solliciter les compétences, mais assez accessible pour que l'activité conserve une prise. L'intéressant n'est pas ici le contraire de l'ennuyeux : il est la zone qui sépare la sous-sollicitation de la sur-sollicitation.

Si la tâche reste fixe tandis que les compétences augmentent, le parcours devient \textbf{anxiété → intéressant/flow → ennui}. Une tâche se déplace ainsi dans le diagramme sans changer matériellement. Cette zone de flow fournit ainsi une définition possible de l'intéressant. Elle doit être distinguée de l'état subjectif d'absorption : l'intérêt peut se concentrer sur le seuil où une difficulté devient presque maîtrisable sans que l'activité soit encore parfaitement fluide.

Une interface sans friction peut imiter cette sensation tout en désengageant les compétences. Supprimer l'obstacle ne produit pas le flow : cela supprime la résistance que le sujet aurait pu incorporer et transforme la fluidité en illusion de maîtrise.

\begin{quote}
\textbf{Intérêt pour la thèse.} Cette carte combat la promesse technologique de création sans effort et fournit une variable relationnelle pour comparer la variation de la tâche et l'apprentissage du sujet.
\end{quote}
{\footnotesize\color{metadata}\textbf{Provenance :} De l'impossibilité de créer.pdf. PDF p. 59-62\par}

\Needspace{8\baselineskip}
\subsection{Le vrai virtuose paie la fluidité par une vie de travail}
{\small\color{metadata} \texttt{idea\_0058}\quad\textbullet\quad argument\quad\textbullet\quad conceptuelle\par}

La vraie virtuosité ne se réduit ni à la vitesse ni à l'absence de faute. Des années de travail incorporent des contraintes techniques et des heuristiques jusqu'à ce que le geste devienne disponible sans calcul conscient. Mais cette fluidité reste singulière : McLaughlin, Di Meola ou Paco de Lucia peuvent être également virtuoses sans produire les mêmes trajectoires ni résoudre les mêmes résistances.

L'apprentissage a aussi un effet moral et perceptif. Plus l'expert maîtrise, plus il voit les contraintes et l'étendue de ce qui lui échappe. La fluidité authentique produit de l'humilité parce qu'elle conserve la mémoire du prix paye.

\begin{quote}
\textbf{Intérêt pour la thèse.} Cette carte oppose la fluidité acquise à la pseudo-fluidité spectaculaire.
\end{quote}
{\footnotesize\color{metadata}\textbf{Provenance :} ESSAI-La Virtuosité à portée des caniches-F. PACHET.pdf. PDF p. 10-13 et 17-19\par}

\Needspace{8\baselineskip}
\subsection{La chalepodoxie : réputation sociale de difficulté}
{\small\color{metadata} \texttt{idea\_0059}\quad\textbullet\quad définition\quad\textbullet\quad conceptuelle\par}

La chalepodoxie désigne l'appropriation sociale d'une difficulté non maîtrisée : donner l'apparence de naviguer dans le difficile pour en tirer prestige et plaisir. Elle met en relation un domaine réputé difficile, un sujet qui n'en a pas surmonte la résistance et un tiers qu'il souhaite impressionner.

Elle ne suppose pas nécessairement un mensonge conscient. Le sujet peut éprouver sincèrement le plaisir de la fluidité ; ce plaisir porte toutefois sur le fait d'être vu comme capable, sans l'obligation de le devenir. La difficulté fonctionne alors comme capital symbolique plutôt que comme épreuve.

\begin{quote}
\textbf{Intérêt pour la thèse.} Cette carte nomme un mécanisme central de \texttt{Caniches}.
\end{quote}
{\footnotesize\color{metadata}\textbf{Provenance :} ESSAI-La Virtuosité à portée des caniches-F. PACHET.pdf. PDF p. 13-16\par}

\Needspace{8\baselineskip}
\subsection{On désire l'état du virtuose plus que son objet}
{\small\color{metadata} \texttt{idea\_0060}\quad\textbullet\quad argument\quad\textbullet\quad conceptuelle\par}

Ce qui attire dans la virtuosité n'est pas seulement l'objet produit, mais l'état suppose du virtuose : aisance, maîtrise, continuité parfaite entre intention et exécution. Le spectacle du geste fluide fait croire cet état proche et disponible, alors même que la raison sait qu'il résulte d'années de travail.

L'image abolit fictivement la distance temporelle de l'apprentissage. En masquant la répétition, la fatigue et l'échec, elle transforme l'admiration en injonction mimétique : voir la maîtrise devient croire à son accessibilité immédiate.

\begin{quote}
\textbf{Intérêt pour la thèse.} Cette carte explique pourquoi la pseudo-virtuosité peut séduire même lorsqu'elle ne produit pas de vraie pensée.
\end{quote}
{\footnotesize\color{metadata}\textbf{Provenance :} ESSAI-La Virtuosité à portée des caniches-F. PACHET.pdf. PDF p. 16-19\par}

\Needspace{8\baselineskip}
\subsection{La virtuosité est perçue par comparaison avec les capacités corporelles latentes du public}
{\small\color{metadata} \texttt{idea\_0078}\quad\textbullet\quad hypothèse\quad\textbullet\quad conceptuelle\par}

Le public n'évalue pas nécessairement la difficulté par une mesure externe. Il peut simuler implicitement le geste et comparer la performance à ce que son propre corps pourrait produire. Les femelles canaris, qui possèdent l'appareil latent du chant sans chanter dans les conditions ordinaires, pourraient ainsi percevoir de l'intérieur le caractère impossible des trilles qu'elles préfèrent.

La difficultuosité est alors une relation entre un geste perçu et un corps possible. Une performance séduit parce que le public sait, au moins en puissance, ce qu'elle coûterait et jusqu'où elle dépasse une norme corporelle partagée.

\begin{quote}
\textbf{Intérêt pour la thèse.} Cette hypothèse donne un mécanisme perceptif et incarne la distinction entre difficulté objective et difficulté ressentie.
\end{quote}
{\footnotesize\color{metadata}\textbf{Provenance :} De l'impossibilité de créer.pdf. PDF p. 51-53\par}

\Needspace{8\baselineskip}
\subsection{La virtuosité se reconnaît à la capacité d'absorber l'imprévu sans perdre la cohérence du geste}
{\small\color{metadata} \texttt{idea\_0089}\quad\textbullet\quad argument\quad\textbullet\quad conceptuelle\par}

Exécuter sans faute une prouesse fixée à l'avance ne suffit pas à définir la virtuosité. Le virtuose peut recevoir une contrainte impromptue, rencontrer un obstacle ou choisir une autre trajectoire sans que son langage se désagrège. La maîtrise se manifeste par la disponibilité de nombreuses options cohérentes, pas seulement par la perfection d'une séquence répétée.

Cette capacité vient de contraintes et d'heuristiques incorporées par une longue pratique. Elle permet de répondre à l'imprévisible dans le style propre du sujet, au lieu d'appliquer consciemment une recette à chaque événement.

\begin{quote}
\textbf{Intérêt pour la thèse.} Cette proposition rapproche virtuosité et création : le nouveau peut surgir d'une réponse située lorsque la maîtrise rend plusieurs chemins réellement praticables.
\end{quote}
{\footnotesize\color{metadata}\textbf{Provenance :} ESSAI-La Virtuosité à portée des caniches-F. PACHET.pdf. PDF p. 10-13\par}

\Needspace{8\baselineskip}
\subsection{L'expertise dissipe le mystère d'une performance mais peut approfondir l'admiration}
{\small\color{metadata} \texttt{idea\_0090}\quad\textbullet\quad hypothèse\quad\textbullet\quad conceptuelle\par}

L'apprentissage rend visibles les procédés qui semblaient magiques et peut faire tomber une difficulté apparente. Mais il révèle aussi les contraintes que la performance a effectivement surmontées, les choix possibles et la fragilité du geste. Le connaisseur admire alors moins l'effet spectaculaire et davantage le travail précis qui le rend possible.

Cette double transformation explique pourquoi l'expertise n'entraîne pas nécessairement le désenchantement. Elle dissipe la pseudo-profondeur tout en augmentant la sensibilité aux résistances réelles. La mémoire du prix devient un instrument de discernement face aux fluidités trop faciles.

\begin{quote}
\textbf{Intérêt pour la thèse.} Cette hypothèse montre comment la compréhension peut détruire un intérêt fondé sur l'opacité tout en produisant un intérêt plus informé et plus durable.
\end{quote}
{\footnotesize\color{metadata}\textbf{Provenance :} ESSAI-La Virtuosité à portée des caniches-F. PACHET.pdf; The Mystery of Jotney Songs.pdf. PDF p. 17-19 et 33 Jotney, analyse qui remplace le tour par la forme et augmente la singularité, PDF p. 3\par}

\section{Culture, langage et politique de l'attention}

\Needspace{8\baselineskip}
\subsection{L'homogénéisation culturelle appelle une innovation structurée}
{\small\color{metadata} \texttt{idea\_0026}\quad\textbullet\quad argument\quad\textbullet\quad conceptuelle\par}

La pression industrielle produit beaucoup de contenus imitatifs et homogènes. Une production orientée par la demande anticipée tend à réutiliser les traits des succès antérieurs et peut accélérer cette convergence. Le problème de l'IA créative n'est donc pas seulement d'automatiser la production, mais d'introduire des mécanismes d'innovation structurante sans confondre nouveauté et anomalie locale.

\begin{quote}
\textbf{Intérêt pour la thèse.} Cette carte donne un enjeu historique et économique au problème de l'intéressant.
\end{quote}
{\footnotesize\color{metadata}\textbf{Provenance :} Paper citations updated.docx. Projet SenseOfDirection, overview, rendu PDF p. 2\par}

\Needspace{8\baselineskip}
\subsection{L'injonction démocratique à créer}
{\small\color{metadata} \texttt{idea\_0041}\quad\textbullet\quad argument\quad\textbullet\quad conceptuelle\par}

Les industries technologiques promettent de démocratiser la création en supprimant les difficultés. Cette promesse transforme la création en slogan de masse et confond expression facile, circulation et valeur créative.

\begin{quote}
\textbf{Intérêt pour la thèse.} Cette carte pose l'adversaire principal : le discours contemporain selon lequel tout le monde peut créer pourvu qu'un outil rende l'acte fluide.
\end{quote}
{\footnotesize\color{metadata}\textbf{Provenance :} De l'impossibilité de créer.pdf. PDF p. 3-9\par}

\Needspace{8\baselineskip}
\subsection{L'IA apprend nos répétitions plutôt que la singularité des œuvres}
{\small\color{metadata} \texttt{idea\_0054}\quad\textbullet\quad argument\quad\textbullet\quad conceptuelle\par}

Un modèle génératif apprend ce qui revient assez souvent pour former une régularité. Un corpus d'œuvres radicalement singulières et sans points communs lui apparaîtrait comme du bruit. Sa performance ne prouve donc pas qu'il a absorbé l'essence des œuvres ; elle révèle l'ampleur de l'auto-similarité de nos productions.

Utilisée comme anti-modèle, l'IA peut cartographier ce qui vient facilement, les pentes statistiques d'un style et les solutions déjà disponibles. Créer avec elle consisterait alors moins à accepter ses sorties qu'à chercher précisément ce qu'elle ne sait pas généraliser.

Le transfert de style illustre cette limite : il peut reproduire une texture locale tout en empruntant à une cible la structure globale qui lui manque. Cette apparence de direction ne prouve pas que le modèle a appris à engendrer une forme adaptée à son propre contenu.

\begin{quote}
\textbf{Intérêt pour la thèse.} Cette carte prolonge la question de la création impossible vers celle de la dépossession contemporaine.
\end{quote}
{\footnotesize\color{metadata}\textbf{Provenance :} De l'impossibilité de créer.pdf; Paper citations updated.docx. PDF p. 83-87 Projet SenseOfDirection, critique du style transfer, rendu PDF p. 2-5\par}

\Needspace{8\baselineskip}
\subsection{Le créateur veut circuler sans être dissous}
{\small\color{metadata} \texttt{idea\_0055}\quad\textbullet\quad argument\quad\textbullet\quad conceptuelle\par}

Le créateur veut être vu, lu, entendu, mais craint que son œuvre, en circulant, cesse de lui appartenir. La diffusion est à la fois accomplissement et dépossession.

\begin{quote}
\textbf{Intérêt pour la thèse.} Cette carte situe la création dans une économie d'usages, de reprises et de transformations.
\end{quote}
{\footnotesize\color{metadata}\textbf{Provenance :} De l'impossibilité de créer.pdf. PDF p. 91-93\par}

\Needspace{8\baselineskip}
\subsection{La création contemporaine est distribuée}
{\small\color{metadata} \texttt{idea\_0056}\quad\textbullet\quad argument\quad\textbullet\quad conceptuelle\par}

La figure romantique du créateur solitaire masque une réalité distribuée : artistes, outils, publics, infrastructures, corpus et machines apprenantes participent tous à la circulation et à la transformation des œuvres.

\begin{quote}
\textbf{Intérêt pour la thèse.} Cette carte peut servir à déplacer la thèse de l'individu créateur vers l'écologie de création.
\end{quote}
{\footnotesize\color{metadata}\textbf{Provenance :} De l'impossibilité de créer.pdf. PDF p. 91-96\par}

\Needspace{8\baselineskip}
\subsection{Appeler une œuvre contenu la transforme en substance interchangeable}
{\small\color{metadata} \texttt{idea\_0073}\quad\textbullet\quad argument\quad\textbullet\quad conceptuelle\par}

Le mot \texttt{contenu} n'est pas une désignation neutre. Il regroupe le roman, le quatuor, la vidéo et le tutoriel comme autant de substances destinées à remplir les mêmes infrastructures. cette catégorie privilégie ce qui se mesure en volume, durée ou engagement et rend secondaires la singularité, la facture et la valeur symbolique de chaque œuvre.

La liquéfaction est donc conceptuelle avant d'être technique : en traitant les œuvres comme des matières premières interchangeables, le vocabulaire des plateformes prépare leur circulation, leur découpage et leur monétisation.

\begin{quote}
\textbf{Intérêt pour la thèse.} Cette proposition montre qu'une catégorie descriptive peut transformer son objet. Elle relie langage, économie de l'attention et effacement de la singularité.
\end{quote}
{\footnotesize\color{metadata}\textbf{Provenance :} De l'impossibilité de créer.pdf. PDF p. 4-6\par}

\Needspace{8\baselineskip}
\subsection{Céline : créer comme ruse contre la langue}
{\small\color{metadata} \texttt{idea\_0052}\quad\textbullet\quad exemple\quad\textbullet\quad conceptuelle\par}

Créer peut consister à lutter contre une langue déjà formée, à la déformer par des ruses rythmiques, syntaxiques ou stylistiques pour rouvrir un espace d'expression.

\begin{quote}
\textbf{Intérêt pour la thèse.} Cette carte nuance la contrainte : elle n'est pas seulement règle externe, elle peut devenir adversaire intime de l'écriture.
\end{quote}
{\footnotesize\color{metadata}\textbf{Provenance :} De l'impossibilité de créer.pdf. PDF p. 68-71\par}

\Needspace{8\baselineskip}
\subsection{Les fantasmagories séduisent par plaisir de la glisse}
{\small\color{metadata} \texttt{idea\_0057}\quad\textbullet\quad argument\quad\textbullet\quad conceptuelle\par}

Les généalogies intellectuelles, les biais cognitifs ou les crises sociales peuvent expliquer d'où viennent les fantasmagories, mais pas entièrement pourquoi elles captent l'attention. Leur efficacité peut tenir au plaisir qu'elles procurent : une cohérence rapide, un sentiment de surplomb et la possibilité d'enchaîner des causes sans rencontrer de résistance.

L'hypothèse est donc aussi esthétique et hédoniste. Le locuteur et son public ne sont pas seulement trompés par un contenu ; ils recherchent une expérience de glisse qui ressemble, à faible coût, à l'état du virtuose.

\begin{quote}
\textbf{Intérêt pour la thèse.} Cette carte précise pourquoi il ne suffit pas de corriger des erreurs factuelles : il faut comprendre le plaisir formel du discours.
\end{quote}
{\footnotesize\color{metadata}\textbf{Provenance :} ESSAI-La Virtuosité à portée des caniches-F. PACHET.pdf. PDF p. 5-9\par}

\Needspace{8\baselineskip}
\subsection{Le langage est le terrain privilégié des impostures virtuoses}
{\small\color{metadata} \texttt{idea\_0061}\quad\textbullet\quad argument\quad\textbullet\quad conceptuelle\par}

Le langage permet très tôt à chacun de produire des séquences fluides. Pour simuler la virtuosité, il suffit alors de relâcher des contraintes dont la violation est moins immédiatement sanctionnée que dans le sport ou la musique : vérité, cohérence à long terme, réfutabilité, précision et responsabilité.

La parole conserve l'allure d'une performance tout en supprimant la possibilité de la fausse note. C'est pourquoi le langage devient un terrain privilégié pour confondre continuité verbale et maîtrise cognitive.

\begin{quote}
\textbf{Intérêt pour la thèse.} Cette carte relie virtuosité, sophistique contemporaine et critique de la pensée sans résistance.
\end{quote}
{\footnotesize\color{metadata}\textbf{Provenance :} ESSAI-La Virtuosité à portée des caniches-F. PACHET.pdf. PDF p. 20-23\par}

\Needspace{8\baselineskip}
\subsection{Penser, c'est accepter la contrainte du réel}
{\small\color{metadata} \texttt{idea\_0062}\quad\textbullet\quad distinction\quad\textbullet\quad conceptuelle\par}

La pensée véritable accepte l'épreuve des faits, de la logique, de l'histoire et de la forme. La fantasmagorie conserve l'allure de l'improvisation mais se donne ensuite l'apparence d'une structure.

\begin{quote}
\textbf{Intérêt pour la thèse.} Cette carte donne un critère pour distinguer fluidité argumentative et travail conceptuel.
\end{quote}
{\footnotesize\color{metadata}\textbf{Provenance :} ESSAI-La Virtuosité à portée des caniches-F. PACHET.pdf. PDF p. 20-23\par}

\Needspace{8\baselineskip}
\subsection{La substitution harmonique comme modèle de discours pseudo-profond}
{\small\color{metadata} \texttt{idea\_0063}\quad\textbullet\quad exemple\quad\textbullet\quad conceptuelle\par}

Un discours peut remplacer des formulations simples par des syntagmes plus lourds et plus dissonants, sans ajouter d'information. Comme une substitution harmonique, il conserve la fonction tout en donnant une impression de richesse.

\begin{quote}
\textbf{Intérêt pour la thèse.} Cette carte donne un outil d'analyse très concret pour les pseudo-virtuosités discursives.
\end{quote}
{\footnotesize\color{metadata}\textbf{Provenance :} ESSAI-La Virtuosité à portée des caniches-F. PACHET.pdf. PDF p. 24-25\par}

\Needspace{8\baselineskip}
\subsection{L'encerclement mélodique du discours}
{\small\color{metadata} \texttt{idea\_0064}\quad\textbullet\quad exemple\quad\textbullet\quad conceptuelle\par}

Un discours peut tourner autour d'une idée par reformulations successives et montées en généralité. En musique, l'encerclement n'a de sens que parce qu'il doit atteindre une note cible imposée par l'harmonie. Sans cible conceptuelle définie, condition de falsification ou objection capable de l'interrompre, le discours peut circuler indéfiniment sans jamais sonner faux.

L'impression de travail conceptuel vient alors du mouvement autour de la notion, non d'une progression démonstrative. La butée est précisément ce qui distingue la variation risquée de la paraphrase protégée.

\begin{quote}
\textbf{Intérêt pour la thèse.} Cette carte précise comment une forme musicale peut devenir instrument critique de la prose idéologique.
\end{quote}
{\footnotesize\color{metadata}\textbf{Provenance :} ESSAI-La Virtuosité à portée des caniches-F. PACHET.pdf. PDF p. 26-27\par}

\Needspace{8\baselineskip}
\subsection{La pédale de distorsion verbale}
{\small\color{metadata} \texttt{idea\_0065}\quad\textbullet\quad exemple\quad\textbullet\quad conceptuelle\par}

Des mots explosifs peuvent fonctionner comme une pédale de distorsion : ils saturent le signal, gomment les hésitations, prolongent l'effet et donnent à une attaque faible l'intensité d'un geste maîtrise.

\begin{quote}
\textbf{Intérêt pour la thèse.} Cette carte enrichit la grammaire musicale des impostures discursives.
\end{quote}
{\footnotesize\color{metadata}\textbf{Provenance :} ESSAI-La Virtuosité à portée des caniches-F. PACHET.pdf. PDF p. 28-32\par}

\Needspace{8\baselineskip}
\subsection{Une démocratie tient par résistance à la glisse}
{\small\color{metadata} \texttt{idea\_0066}\quad\textbullet\quad argument\quad\textbullet\quad conceptuelle\par}

On ne combat pas une fantasmagorie seulement par des faits, car elle procure un plaisir de fluidité et de maîtrise. Une réponse exclusivement corrective laisse donc intacte la gratification qui porte le discours. Il faut rendre de nouveau désirable une autre satisfaction : celle d'une pensée qui accepte la lenteur, la résistance, l'interruption et la possibilité d'être corrigée.

La résistance est une condition politique, pas seulement cognitive. Une démocratie devient vulnérable lorsque la continuité du débit vaut preuve et que l'hésitation est interprétée comme faiblesse.

\begin{quote}
\textbf{Intérêt pour la thèse.} Cette carte donne à l'analyse de la virtuosité une portée politique.
\end{quote}
{\footnotesize\color{metadata}\textbf{Provenance :} ESSAI-La Virtuosité à portée des caniches-F. PACHET.pdf. PDF p. 33-35\par}

\Needspace{8\baselineskip}
\subsection{Déformer la langue peut forcer à penser ou seulement forcer à admirer}
{\small\color{metadata} \texttt{idea\_0091}\quad\textbullet\quad distinction\quad\textbullet\quad conceptuelle\par}

Le grand styliste rend la langue étrangère à elle-même parce qu'une idée ou une expérience résiste aux formes disponibles. La déformation ouvre un passage à la limite du dicible et oblige le lecteur à réorganiser sa pensée.

La chalepodoxie peut employer les mêmes signes visibles - jargon, abstraction, syntaxe brisée, obscurité - sans rencontrer cette résistance. La complexité ne porte alors aucune proposition plus exigeante ; elle sert à produire une réputation de profondeur. Dans un cas, la difficulté est le coût d'une idée ; dans l'autre, elle est l'effet destiné au public.

\begin{quote}
\textbf{Intérêt pour la thèse.} Cette distinction évite d'assimiler obscurité et invention, tout en reconnaissant qu'une forme véritablement nouvelle peut d'abord être difficile à comprendre.
\end{quote}
{\footnotesize\color{metadata}\textbf{Provenance :} ESSAI-La Virtuosité à portée des caniches-F. PACHET.pdf. PDF p. 30-33\par}

\Needspace{8\baselineskip}
\subsection{Les réseaux sociaux enseignent à paraître maître en punissant l'hésitation}
{\small\color{metadata} \texttt{idea\_0092}\quad\textbullet\quad hypothèse\quad\textbullet\quad conceptuelle\par}

Les formats numériques sélectionnent et mettent en scène des séquences continues, rapides et sans accroc. Ils montrent le résultat plutôt que les essais et valorisent la réponse immédiate plutôt que l'exposition à la correction. L'utilisateur apprend ainsi qu'il faut parler sans hésiter, produire un récit complet et donner un visage de maîtrise.

Les réseaux ne créent pas seuls la confusion entre fluidité et vérité, mais ils en font une norme pratique répétée quotidiennement. L'hésitation, qui pourrait signaler la rencontre d'une difficulté réelle, devient un défaut de performance à supprimer.

\begin{quote}
\textbf{Intérêt pour la thèse.} Cette hypothèse décrit un mécanisme institutionnel par lequel une forme expressive modifie les critères publics de crédibilité et d'intérêt.
\end{quote}
{\footnotesize\color{metadata}\textbf{Provenance :} ESSAI-La Virtuosité à portée des caniches-F. PACHET.pdf. PDF p. 9-13 et 33-35\par}

\Needspace{8\baselineskip}
\subsection{La promesse de compréhension peut fasciner sans produire de compréhension}
{\small\color{metadata} \texttt{idea\_0124}\quad\textbullet\quad hypothèse\quad\textbullet\quad conceptuelle\par}

Certaines choses fascinent sans que le sujet puisse dire ce qu'il comprend ni pourquoi il continue à écouter. Elles entretiennent la promesse qu'une clef décisive va bientôt réorganiser le monde, puis diffèrent continuellement sa livraison. Chaque proposition familière semble préparer une révélation ; chaque difficulté ou obscurité peut être interprétée comme l'indice d'une profondeur encore inaccessible.

Ce régime simule le travail constructif de l'intéressant. Le sujet ne construit pas nécessairement un modèle plus précis ; il maintient plutôt une croyance de second ordre : « je suis sur le point de comprendre ». La promesse suffit à relancer l'attention, même si les prédictions, distinctions et capacités d'action du sujet ne se transforment pas. La fascination peut ainsi survivre à l'absence répétée de gain cognitif.

Dans la géométrie du flow, cette rhétorique maintient artificiellement le sujet près de la frontière de l'anxiété. Des banalités, anecdotes ou formules familières lui donnent assez de prises pour éviter l'opacité complète ; la cadence, le vocabulaire technique ou l'annonce d'une révélation empêchent simultanément la chute dans l'ennui. Le discours promet la zone intéressante sans nécessairement permettre la construction qui devrait s'y accomplir.

\begin{quote}
\textbf{Intérêt pour la thèse.} Cette hypothèse empêche d'identifier trop vite fascination, flow et construction cognitive. Elle fournit un cas limite pour tester la définition de l'intéressant : la promesse d'un gain peut-elle soutenir durablement l'attention lorsque le gain lui-même reste nul ou invérifiable ?
\end{quote}
{\footnotesize\color{metadata}\textbf{Provenance :} ESSAI-La Virtuosité à portée des caniches-F. PACHET.pdf. La Virtuosité à la portée des caniches, plaisir de la glisse, langage et pseudo-profondeur, PDF p. 5-9 et 20-25 Tony Robbins et les frères Bogdanov sont ajoutés comme cas à documenter ; leur interprétation proposée ici reste une hypothèse de la thèse\par}
{\footnotesize\color{metadata}\textbf{Références :} \parencite{robbinsPersonalPower,baez2004bogdanoff,bogdanoff2002etat}\par}
\clearpage
\nocite{allouche1999thueMorse,baez2004bogdanoff,berger1999expectations,bogdanoff2002etat,chomsky1957syntactic,colton2000interestingness,feynman1999pleasure,florman1996existential,gang1999unified,goldman1999netneg,ishizu2011brainbeauty,kelly2002music,macherey1995vieaffective,meyer1956emotion,narmour1990basic,narmour1992complexity,oudeyer2007intrinsic,oudeyer2007typology,pachet2018oreille,pachet2026biases,petitot1993phenomenologie,petitot2002naturaliser,robbinsPersonalPower,schmidhuber1997interesting,schmidhuber2009compression,silberschatz1995subjective,silberschatz1996patterns,spiliopoulou1999rules,spinoza1861oeuvres,spinoza1966ethique,steels2005coordinating,steels2015talkingheads,zeki2018mathematicalbeauty}
\printbibliography[title={Bibliographie des cartes}]

\end{document}
